\documentclass[11pt,a4paper]{article}

\usepackage{amsmath,amssymb,amsthm}
\usepackage{geometry}
\usepackage{hyperref}
\usepackage{booktabs}
\usepackage{array}
\usepackage{enumitem}
\usepackage{titlesec}
\usepackage{xcolor}
\usepackage{fancyhdr}
\usepackage[expansion=false]{microtype}
\usepackage{parskip}
\usepackage{listings}
\usepackage{mdframed}

\geometry{margin=2.5cm}

\definecolor{deepblue}{RGB}{31,56,100}
\definecolor{midblue}{RGB}{46,84,150}
\definecolor{codebg}{RGB}{245,247,250}
\definecolor{coderule}{RGB}{200,210,225}

\titleformat{\section}{\large\bfseries\color{deepblue}}{{\thesection}}{1em}{}
\titleformat{\subsection}{\normalsize\bfseries\color{midblue}}{{\thesubsection}}{1em}{}
\titleformat{\subsubsection}{\normalsize\bfseries}{{\thesubsubsection}}{1em}{}

\lstset{
  backgroundcolor=\color{codebg},
  basicstyle=\ttfamily\small,
  breaklines=true,
  frame=single,
  rulecolor=\color{coderule},
  framesep=6pt,
  xleftmargin=8pt,
  xrightmargin=8pt,
  aboveskip=8pt,
  belowskip=8pt,
}

\newmdenv[
  backgroundcolor=codebg,
  linecolor=coderule,
  linewidth=0.5pt,
  innerleftmargin=10pt,
  innerrightmargin=10pt,
  innertopmargin=8pt,
  innerbottommargin=8pt,
  skipabove=8pt,
  skipbelow=8pt,
]{eqbox}

\pagestyle{fancy}
\fancyhf{}
\fancyhead[L]{\small\color{gray}Linearized GR Sandbox}
\fancyhead[R]{\small\color{gray}Design Document\,---\,v35}
\fancyfoot[C]{\small\thepage}
\renewcommand{\headrulewidth}{0.4pt}

\newcommand{\vv}[1]{\mathbf{#1}}
\newcommand{\Bg}{B_g}
\newcommand{\Phig}{\Phi_g}
\newcommand{\Ag}{\vv{A}_g}
\newcommand{\uEM}{u_{\mathrm{EM}}}
\newcommand{\order}[1]{\mathcal{O}(#1)}
\newcommand{\dd}{\mathrm{d}}
\newcommand{\half}{\tfrac{1}{2}}

\newcolumntype{L}[1]{>{\raggedright\arraybackslash}p{#1}}
\newcolumntype{C}[1]{>{\centering\arraybackslash}p{#1}}

% Numbered algorithm steps with consistent formatting
\newlist{algsteps}{enumerate}{1}
\setlist[algsteps]{label=\textbf{Step \arabic*.},
                   leftmargin=*,
                   labelwidth=5em,
                   itemsep=6pt}

%--------------------------------------------------------------------
\begin{document}
%--------------------------------------------------------------------

\begin{titlepage}
  \centering
  \vspace*{3cm}
  {\Huge\bfseries\color{deepblue} General Relativity Sandbox\par}
  \vspace{0.8cm}
  {\large\itshape Design and Implementation of a Linearized GR\,+\,Electromagnetism Simulation\par}
  \vspace{0.4cm}
  {\normalsize\color{gray} Comparing Full Linearized Einstein--Maxwell vs.\ GEM/EM Approaches\par}
  \vspace{1.2cm}
  {\large\bfseries\color{midblue} Revision v35\par}
  \vspace{0.2cm}
  {\small\color{gray}
  \textbf{Implementation alignment with §\ref{sec:numerical}'s Yee
  specification, plus opt-in Esirkepov current deposition}.
  Section~\ref{sec:numerical} already specifies the Yee staggered
  grid and the discrete continuity equation as the gauge-preservation
  defense; the v32--v34 implementation departed from that spec by
  using cell-centered storage, which broke discrete continuity and
  produced ``grid-heating'' self-forces strong enough that Stage~11
  (mutual gravity via FDTD) is unreachable.  v35 brings the storage
  layout, deposition kernels, and force-evaluation operators into
  line with what Section~\ref{sec:numerical} prescribes, and adds
  Esirkepov current deposition~\cite{Esirkepov2001} as an opt-in
  path beyond the truncation-order accuracy of
  Section~\ref{sec:numerical}'s CIC-on-Yee recipe, giving exact
  discrete continuity and structural cancellation of the moving-
  particle self-force.\par
  Section~\ref{sec:yee_pivot} contains the motivation and migration
  plan (save points S1--S7); intermediate cell-centered version
  preserved on the \texttt{v34-cell-centered} git branch.\par
  All v34 physics corrections (EIH 1PN sign$+$Shapiro term,
  proper-time Eq.~\ref{eq:proper_time_full}, analytic-background
  mode) carry forward unchanged: they are independent of the grid
  layout.\par}
  \vspace{1.5cm}
  \hrule
\end{titlepage}

\tableofcontents
\newpage

%====================================================================
\section*{Changes from v32}
%====================================================================

This revision makes two changes to the v32 document:

\begin{enumerate}
  \item \textbf{Implementation plan, §\ref{sec:impl}}: a new stage is inserted between the original Stages 5 and 6, introducing \emph{sampled background field arrays} as a prerequisite for the particle-dynamics tests.  All subsequent stages are renumbered (old Stages 6--16 become new Stages 7--17).  The motivation is testability: with sampled backgrounds available before the perturbation-source tests, a single test particle can be placed in any prescribed background (Keplerian orbit around a softened point mass, gyroscope in a frame-dragging field, charge in an external Coulomb potential, etc.) without first having to set up a second body and the FDTD source-deposition path.  The full background-modes coverage formerly in old Stage 14 (Schwarzschild, Kerr, Reissner--Nordstr\"om, Kerr--Newman, binary) remains in the new Stage 15, where it expands the library of generators introduced in the new Stage 6.
  \item \textbf{§\ref{sec:abc} clarification}: the round-trip reflection $R\approx10^{-3}$ derived there is the plane-wave normal-incidence limit; the realized in-interior residual for a 2D circular wavefront is observed to be several percent due to two effects --- the $(1-\sigma\,\dd t)$ linearization differing from $\exp(-\sigma\,\dd t)$ at the $\sigma_{\max}\,\dd t\sim0.46$ used in practice, and oblique-incidence waves into the grid corners being weakly damped by the axis-aligned $d_{i,j}=\max(d_x,d_y)$ profile.  Both effects would be mitigated by a PML; for the qualitative pedagogical purpose of the sandbox, the simple damping layer is adequate.  See the new paragraph at the end of §\ref{sec:abc} for the empirical figures.
\end{enumerate}

All physics chapters (§\ref{sec:gem_metric} through §\ref{sec:abc} and beyond) are unchanged from v32.  Cross-references to \texttt{eq:} labels remain valid.

\newpage

%====================================================================
\section{Implementation Alignment with §\ref{sec:numerical} (v35): Yee + Esirkepov}
\label{sec:yee_pivot}
%====================================================================

\subsection{Why this work is being done}

The Yee staggered grid is already specified in
Section~\ref{sec:numerical} (``Field grid'' and ``Why this staggering:
exact gradient and curl operations'').  Source deposition placed at
the correct sublattices is also specified there, as is the discrete
continuity equation as the primary defense against gauge drift
(``The discrete continuity equation: the primary defense'' in
Section~\ref{sec:numerical}'s gauge-condition subsection).
\textbf{This section is not introducing a new architecture; it is
documenting the migration of the v32--v34 implementation to match
what the doc already specifies, plus one addition (Esirkepov current
deposition) that goes beyond the truncation-order accuracy of
Section~\ref{sec:numerical}'s CIC-on-Yee recipe.}

\paragraph{What the v32--v34 implementation actually did.}  Rather
than the Yee layout specified in Section~\ref{sec:numerical}, the
implementation stored all six potentials, all six source arrays, and
all derived field quantities on a single cell-centered grid at
$(i+\tfrac{1}{2}, j+\tfrac{1}{2})\Delta x$, with CIC deposition into a
single $\rho_{\mathrm{matter}}$ array at cell centers and a
CIC-of-FD operator for the force gradient.  This was the simpler
path to a working scalar wave equation (Stages 1--4), and Stages 5--9
each rely on at most a \emph{static} particle in a \emph{fixed}
background where the Hockney--Eastwood theorem still gives zero
self-force regardless of the underlying grid placement.  Stage 10's
self-force-cancellation test (drift $< 10^{-5}$ cells over $10^4$
timesteps) passes in the cell-centered code at $m = 10^{-3}$.

\paragraph{Where the cell-centered implementation breaks down.}  Once
particles \emph{move} and \emph{deposit dynamic sources}, the
cell-centered code shows pathological behavior.  Both $\rho$ and
$\vv{J}$ live at the same cell-centered nodes there, so the discrete
continuity equation that Section~\ref{sec:numerical} writes for the
Yee layout (with $\rho$ at corners and $J_i$ at edges and a
staggered divergence stencil) is not even satisfied to truncation
order in the cell-centered code.  The 2-point centered FD divergence
at the cell-center $\rho$ position uses $J$ values two cells away,
breaking the consistency that would damp gauge drift.  The numerical
residual takes the form of a ``grid-heating'' self-force on the
moving particle, with the wrong sign — energy is pumped \emph{into}
the orbit rather than radiated out.  Scaling, in our natural units:
\begin{equation}
  \frac{a_{\mathrm{heat}}}{a_{\mathrm{phys}}} \;\sim\; \alpha\,\frac{r}{m_{\mathrm{src}}},
  \qquad \alpha = \mathcal{O}(1)\text{ dimensionless heating coefficient},
\end{equation}
where $r$ is the orbital scale and $m_{\mathrm{src}}$ is the
\emph{source} mass on the other side of the mutual interaction.

For a single test particle around a fixed analytic background of
mass $M$, this gives $a_{\mathrm{heat}}/a_{\mathrm{phys}} \sim \alpha r / M$, which
\emph{is} suppressed by going to small $m_{\mathrm{test}}$ (Stage 10 worked
at $m_{\mathrm{test}} = 10^{-6}$ in the weak-coupling regime).  But for
\emph{mutual} gravity between two particles of mass $m$ at separation
$r$, both bodies' self-heating scales like the physical mutual
acceleration, and the ratio is $\sim \alpha r/m$ \emph{independent of
$m$}.  Scaling masses down does not help.  Linearization breaks
down before heating becomes negligible.  Stage 11 (two-body
gravitational interaction via FDTD) is therefore unreachable in the
cell-centered framework.

\paragraph{What this implies for the project.}  The cell-centered
implementation effectively limits FDTD's role to visualization of
fields and to quantitative tests with analytic backgrounds; the
mutual-gravity Stage~11 case is unreachable.  That undermines the
project's framing as a dynamic-field PIC simulation, since the FDTD
ends up carrying weight only as scaffolding around an analytic
backbone.  v35 fixes this by aligning the implementation with the
Section~\ref{sec:numerical} Yee specification.

\subsection{Esirkepov current deposition: the one addition beyond §\ref{sec:numerical}}

Section~\ref{sec:numerical}'s gauge subsection notes that ``using a
consistent CIC kernel for both $\rho$ and $\vv{J}$ deposition
satisfies [the discrete continuity equation] to the accuracy of the
kernel,'' i.e.\ to $\mathcal{O}((k\Delta x)^2)$ truncation error.  That
truncation is the source of the residual grid-heating self-force on
moving particles; on the Yee grid the heating is much smaller than on
cell-centered storage (the discrete divergence stencil now matches
the source-deposition staggering), but not zero.

v35 adds \emph{Esirkepov current deposition}~\cite{Esirkepov2001} as
an opt-in code path beyond what Section~\ref{sec:numerical}
specifies.  Given a particle's positions $\vv{x}^n$ and
$\vv{x}^{n+1}$, Esirkepov's local recurrence computes the deposited
$J_x$ values at horizontal-edge cells and $J_y$ values at
vertical-edge cells such that the discrete continuity equation of
Section~\ref{sec:numerical}
\begin{equation}
\label{eq:discrete_continuity_yee}
  \frac{\rho^{n+1}_{i,j} - \rho^n_{i,j}}{\Delta t}
  + \frac{J_x^{n+1/2}_{i+1/2,j} - J_x^{n+1/2}_{i-1/2,j}}{\Delta x}
  + \frac{J_y^{n+1/2}_{i,j+1/2} - J_y^{n+1/2}_{i,j-1/2}}{\Delta y}
  \;=\; 0
\end{equation}
holds \emph{exactly} (not just to truncation order) at every corner
$(i,j)$.  The recurrence is local (involves only the cells touched by
the particle's trajectory in one timestep) and is roughly 50 lines in
2D.  Once Eq.~(\ref{eq:discrete_continuity_yee}) is satisfied exactly,
the Hockney--Eastwood-style residual self-force on a moving particle
is structurally zero, and the parabolic / hyperbolic cleaning steps
described in Section~\ref{sec:numerical} are not needed to keep gauge
drift bounded.

The opt-in flag lets us keep simple CIC-on-Yee deposition (as
specified in Section~\ref{sec:numerical}) as the default for tests
that don't involve dynamic mutual sources, and switch on Esirkepov
for stages that do (Stages 10--11+).  This matches
Section~\ref{sec:numerical}'s framing of parabolic/hyperbolic cleaning
as optional refinements: the simple recipe works for the common case;
the exact-continuity recipe is brought in when needed.

\subsection{What this enables}

\begin{itemize}
  \item \textbf{Stage 11 becomes feasible.}  Two equal-mass bodies in
        mutual orbit, each sourcing the other via the FDTD perturbation
        field (and reading the resulting force via the Yee gradient
        operator at the particle's CIC-interpolated position), should
        give a stable bound orbit of the correct period.  This is the
        original sandbox vision.
  \item \textbf{Quantitative radiation reaction.}  With
        charge-conserving deposition, the dynamic self-force has the
        correct sign (energy LOSS, not energy gain) and an amplitude
        of the right order.  Orbital decay for an inspiraling binary
        becomes measurable, not buried by spurious heating.
  \item \textbf{Mass-scale independence.}  Tests no longer require
        artificially small $m_{\mathrm{test}}$ in the weak-coupling
        regime; the heating residual is structurally zero.  The
        relevant constraint becomes the linearization condition
        $r_s/r \ll 1$, which is much more permissive.
  \item \textbf{Genuinely dynamic-field PIC.}  The FDTD is again
        load-bearing for quantitative dynamics, not just visualization.
\end{itemize}

\subsection{Migration plan}
\label{sec:yee_migration_plan}

The refactor proceeds in save points, each of which leaves a building
tree with all previously-passing tests still passing.  Tag at every
save point so a regression can be bisected cleanly.

\begin{center}
\begin{tabular}{@{}lL{8.5cm}L{2.5cm}@{}}
\toprule
\textbf{S} & \textbf{What lands} & \textbf{Should pass} \\
\midrule
S0 & v35 doc revision (this section), architectural pivot agreed and
documented; cell-centered v34 preserved on
\texttt{v34-cell-centered} branch. & doc only \\
S1 & Add Yee sublattice constants and documentation to internal
headers (\texttt{sim\_internal.h}, comments).  Storage unchanged.
Code unchanged. & 1--10 \\
S2 & \texttt{src/field.c}: per-field Laplacian uses correct
sublattice neighbors and boundary indexing.  No functional change for
test 1--4 (scalar-only) at the discretization-order accuracy. & 1--4 \\
S3 & \texttt{src/deposit.c}: split CIC kernel into corner-deposit
(rho), x-edge-deposit (J\_x), y-edge-deposit (J\_y) variants.  No
Esirkepov yet --- just W$_2$ on each sublattice. & 1--5 \\
S4 & \texttt{src/background.c}: per-field sampled-background placement.
\texttt{gr\_bg\_eval\_analytic} and \texttt{gr\_bg\_eval\_A\_g} return
values at the correct sublattice. & 1--6 \\
S5 & \texttt{src/particle.c}: replace CIC-of-FD gradient with Yee
gradient (each $\partial_i \Phi$ on its native sublattice, then
CIC-interpolated to the particle).  Replace ad-hoc curl with Yee
curl (Eq.~\ref{eq:yee_curl}). & 1--9 \\
S6 & Esirkepov current deposition as a separate code path (opt-in
flag, default off so existing tests continue using simple W$_2$).
Verify that with Esirkepov on, the Stage 10 self-force test (now at
$m_{\mathrm{test}} = 1$, not $10^{-6}$) gives drift $< 10^{-5}$ cells
over $10^4$ steps. & 1--10 \\
S7 & Stage 11: \texttt{pic\_binary} scenario (two equal masses,
mutual gravity through the FDTD path, Esirkepov on).  Verify orbital
period matches the analytic two-body period to $\sim 1\%$. & 1--11 \\
S8 (defer)
   & Stage 12+: EM Effective GEM sources, spin sources (W$_3$
   deposition), background-spacetime modes. & 12+ \\
\bottomrule
\end{tabular}
\end{center}

\paragraph{Compatibility.}  The public API surface
(\texttt{grlite.h}) changes minimally: field-pointer queries
continue to return \texttt{float*} to the underlying storage, with the
caller responsible for knowing the field's sublattice.  Scenarios
that use only the analytic-background path (Stages 7--9 in v34) are
re-validated unchanged.  The WASM EXPORTED\_FUNCTIONS list grows by
one entry (\texttt{gr\_sim\_set\_esirkepov}).

\paragraph{Risks.}  Anticipated debugging hotspots:

\begin{itemize}
  \item Index off-by-one between corner-indexed and edge-indexed
        arrays (each edge array has $N_x \times (N_y$ or $N_y-1)$
        valid cells depending on direction).
  \item Sign conventions in the Esirkepov recurrence.  Cross-check
        against the published 2D pseudocode in~\cite{Esirkepov2001}
        and at least one reference implementation.
  \item Damping layer's sublattice consistency: each field needs its
        $d_{i,j}$ multiplier evaluated at its own sublattice
        position, otherwise the damping looks different to different
        fields at the boundary.
  \item Background-array sampling formerly placed all generators at
        cell centers.  With Yee, the same analytic formula must be
        sampled at corner, x-edge, and y-edge positions, with
        different sub-cell offsets each time.
\end{itemize}

\paragraph{What does NOT change.}  All v34 physics corrections from
Section~\ref{sec:relcorr} (EIH 1PN with sign$+$Shapiro term),
Section~\ref{sec:proptime} (gravitomagnetic proper-time formula),
and Section~\ref{sec:bg_mode} (analytic-background evaluation mode)
remain valid and carry forward unchanged: the physics is independent
of the grid layout used to discretize the field equations.

\newpage

%====================================================================
\section{Project Goal}
%====================================================================

The aim is to build an interactive general relativity sandbox that runs in real time on a desktop computer.  A full numerical relativity simulation using the ADM (Arnowitt--Deser--Misner) framework with a Cauchy boundary is computationally infeasible for this purpose, as it typically requires supercomputer resources.  Instead, the simulation is based on \emph{linearized} general relativity, which captures the most physically interesting effects while remaining tractable.

In the linearized theory, the metric perturbation satisfies a wave equation---the d'Alembert equation---which can be time-stepped efficiently on a grid.  Massive particles and charged test particles then follow trajectories determined by the geodesic equation in this perturbed spacetime, with electromagnetic forces added via the standard Lorentz force.

The sandbox also incorporates electromagnetism, using the full coupled Einstein--Maxwell system in the linearized limit.  This means the electromagnetic field contributes to the gravitational field via its stress-energy tensor, and the gravitational field in turn affects the propagation of electromagnetic waves.

%====================================================================
\section{Physical Framework}
%====================================================================

\subsection{Linearized General Relativity}

In linearized GR, the metric is written as a small perturbation around flat (Minkowski) spacetime:
\begin{equation}
  g_{\mu\nu} = \eta_{\mu\nu} + h_{\mu\nu}, \qquad |h_{\mu\nu}| \ll 1.
\end{equation}
In the Lorenz gauge, the trace-reversed perturbation $\bar{h}^{\mu\nu} \equiv h^{\mu\nu} - \tfrac{1}{2}\eta^{\mu\nu}h$ satisfies a wave equation sourced by the full stress-energy tensor:
\begin{equation}
  \Box\,\bar{h}^{\mu\nu} = -\frac{16\pi G}{c^4}\,T^{\mu\nu}_{\mathrm{total}}.
\end{equation}
This is structurally identical to the wave equation for the electromagnetic four-potential, so FDTD methods from computational electromagnetics apply directly.

\subsection{The GEM Metric}
\label{sec:gem_metric}

In the Lorenz gauge, choosing the GEM potentials $\Phig$ and $\Ag$ to represent the relevant components of $\bar{h}^{\mu\nu}$, the full linearized metric perturbation is:
\begin{equation}
\label{eq:gem_metric}
  h_{00} = -\frac{2\Phig}{c^2}, \qquad
  h_{0i} = h_{i0} = \frac{2A_{g,i}}{c}, \qquad
  h_{ij} = -\frac{2\Phig}{c^2}\,\delta_{ij},
\end{equation}
giving the line element:
\begin{equation}
\label{eq:gem_line_element}
  \dd s^2 = -\!\left(1+\frac{2\Phig}{c^2}\right)c^2\dd t^2
            + \frac{4}{c}\,\Ag\cdot\dd\vv{x}\,\dd t
            + \left(1-\frac{2\Phig}{c^2}\right)\!\left(\dd x^2+\dd y^2+\dd z^2\right).
\end{equation}
The trace-reversed perturbation $\bar{h}_{\mu\nu} = h_{\mu\nu} - \tfrac{1}{2}\eta_{\mu\nu}h$ has components:
\begin{equation}
  \bar{h}_{00} = -\frac{4\Phig}{c^2}, \qquad
  \bar{h}_{0i} = \frac{4A_{g,i}}{c}, \qquad
  \bar{h}_{ij} = 0,
\end{equation}
which is why the GEM wave equations decouple cleanly: $\Box\Phig \propto \rho$ and $\Box\Ag \propto \vv{J}$, with no spatial stress sourcing the spatial metric at this order.

\subsubsection*{What each metric component produces}

Each of the three distinct components of the metric in Eq.~(\ref{eq:gem_metric}) is responsible for a different class of physical effects:

\paragraph{$h_{00} = -2\Phig/c^2$ (time-time component):}
This is the Newtonian gravitational potential in relativistic clothing.  It is the only component present for a static non-rotating mass.  Its physical consequences include:
\begin{itemize}
  \item \textbf{Gravitational attraction}: the gradient $-\nabla\Phig$ gives the Newtonian force.
  \item \textbf{Gravitational time dilation}: clocks in a gravitational well ($\Phig < 0$) run slow by factor $\sqrt{1+2\Phig/c^2}$, contributing $+\Phig/c^2$ to $\dd\tau/\dd t$ at first order.
  \item \textbf{Shapiro delay}: light (and any massless signal) is slowed in a gravitational potential, producing the delay $\Delta t = -2\int\Phig\,\dd\ell/c^3$.
  \item \textbf{Half of light bending}: the $h_{00}$ component alone would deflect light by $2GM/(c^2 b)$ --- the Newtonian result.
\end{itemize}

\paragraph{$h_{ij} = -2\Phig/c^2\,\delta_{ij}$ (spatial diagonal components):}
The spatial metric is isotropic and also proportional to $\Phig$ --- the same scalar as the time-time component.  This component has no Newtonian analogue and is invisible to slow particles ($v\ll c$), but has important consequences for relativistic motion:
\begin{itemize}
  \item \textbf{Second half of light bending}: the spatial metric contributes an equal and opposite deflection to $h_{00}$, doubling the total to $4GM/(c^2 b)$ --- the correct GR result.  This is why Newtonian gravity predicts half the correct deflection.
  \item \textbf{Post-Newtonian force corrections}: the $(v^2/c^2)\vv{g}$ and $4(\vv{v}\cdot\vv{g})\vv{v}/c^2$ terms in the geodesic expansion (Section~\ref{sec:relcorr}) arise from $h_{ij}$ contracting with the spatial velocity.
  \item \textbf{Perihelion precession}: the post-Newtonian force corrections from $h_{ij}$ are responsible for the extra orbital precession beyond the Newtonian result.
  \item \textbf{ISCO}: the additional attractive force from $h_{ij}$ at high velocities removes the minimum from the effective potential at small radii, producing the innermost stable circular orbit.
\end{itemize}

\paragraph{$h_{0i} = 2A_{g,i}/c$ (off-diagonal time-space components):}
These components are nonzero only for rotating or moving masses, arising from the mass current $\vv{J}_{\mathrm{matter}}$ sourcing $\Ag$.  They have no Newtonian analogue.  Their physical consequences are:
\begin{itemize}
  \item \textbf{Gravitomagnetic force}: the off-diagonal metric produces the $4\vv{v}\times\vv{B}_g$ term in the GEM Lorentz force, responsible for Lense--Thirring orbital plane precession and the factor-of-4 difference from the EM analogy.
  \item \textbf{Gyroscope precession}: a spinning test body in a gravitomagnetic field precesses at rate $\boldsymbol{\Omega}_{\mathrm{LT}} = \nabla\times\Ag$, the Lense--Thirring effect.
  \item \textbf{Gravitomagnetic time dilation}: the $h_{0i}$ component contributes $4\Ag\cdot\vv{v}/c^2$ to the proper time rate (Section~\ref{sec:proptime}), with opposite sign for co- and counter-rotating orbits --- the gravitomagnetic clock effect.
  \item \textbf{Frame dragging}: the dragging of inertial frames by a rotating mass is entirely encoded in $h_{0i}$; the ergosphere of the Kerr metric is the region where $h_{0i}$ terms dominate and nothing can remain stationary.
\end{itemize}

The three contributions are independent in the sense that $h_{00}$ and $h_{ij}$ are present for any mass regardless of rotation, while $h_{0i}$ requires a mass current.  Setting $\Ag = 0$ (non-rotating background) eliminates all gravitomagnetic effects and leaves only the Newtonian potential and its relativistic corrections.  The simulation can therefore interpolate between the Newtonian regime ($v\ll c$, $\Ag = 0$), the post-Newtonian regime ($v^2/c^2$ corrections from $h_{ij}$), and the full gravitomagnetic regime (all three metric components active).

\subsection{Gravitoelectromagnetism (GEM)}

In the slow-motion, weak-field limit the linearized GR equations reduce to a form mathematically identical to Maxwell's equations.  The gravitoelectric field $\vv{g}$ (analogous to $\vv{E}$) and gravitomagnetic field $\vv{B}_g$ (analogous to $\vv{B}$) satisfy:
\begin{align}
  \nabla\cdot\vv{g}   &= -4\pi G\rho, \\
  \nabla\cdot\vv{B}_g &= 0, \\
  \nabla\times\vv{g}  &= -\frac{\partial\vv{B}_g}{\partial t}, \\
  \nabla\times\vv{B}_g &= -\frac{4\pi G}{c^2}\vv{J} + \frac{1}{c^2}\frac{\partial\vv{g}}{\partial t}.
\end{align}
The GEM Lorentz force on a test particle with rest mass $m$ and electric charge $q$ is, in terms of the derived fields:
\begin{equation}
  \vv{F} = m\!\left(\vv{g} + 4\vv{v}\times\vv{B}_g\right) + q\!\left(\vv{E} + \vv{v}\times\vv{B}\right).
\end{equation}
The factor of 4 in the gravitomagnetic term arises from the spin-2 nature of the graviton and is responsible for the Lense--Thirring precession rate differing from a naive EM analogy.

Since the simulation stores potentials rather than fields, it is useful to have the equivalent expression directly in terms of potentials.  Substituting $\vv{g} = -\nabla\Phig - \partial_t\Ag$, $\vv{B}_g = \nabla\times\Ag$, $\vv{E} = -\nabla\phi - \partial_t\vv{A}$, $\vv{B} = \nabla\times\vv{A}$ and using the vector identity $\vv{v}\times(\nabla\times\vv{A}) = \nabla(\vv{v}\cdot\vv{A}) - (\vv{v}\cdot\nabla)\vv{A}$:
\begin{align}
  \vv{F}_{\mathrm{GEM}} &= m\!\left(-\nabla\Phig - \partial_t\Ag
    + 4\nabla(\vv{v}\cdot\Ag) - 4(\vv{v}\cdot\nabla)\Ag\right), \label{eq:force_gem_pot} \\
  \vv{F}_{\mathrm{EM}}  &= q\!\left(-\nabla\phi - \partial_t\vv{A}
    + \nabla(\vv{v}\cdot\vv{A}) - (\vv{v}\cdot\nabla)\vv{A}\right). \label{eq:force_em_pot}
\end{align}
These expressions involve only the potentials and their spatial and temporal derivatives evaluated at the particle position.  No curl operation over the full grid is required.  The physical force is gauge-invariant: the apparent gauge dependence of the individual potential terms cancels exactly when all terms are combined, as can be verified by applying a gauge transformation $\phi\to\phi-\partial_t\Lambda$, $\vv{A}\to\vv{A}+\nabla\Lambda$ and confirming $\delta\vv{F}=0$.

\subsection{Particle Properties: Mass and Electric Charge}

Each particle carries two \emph{independent} intrinsic properties: a rest mass $m$ and an electric charge $q$.  They couple to entirely separate fields:
\begin{itemize}
  \item \textbf{Rest mass $m$} sources the GEM field via $\rho_{\mathrm{matter}} = \sum_i \gamma m_i\,\delta(\vv{x}-\vv{x}_i)$ and $\vv{J}_{\mathrm{matter}} = \sum_i \gamma m_i\vv{v}_i\,\delta(\vv{x}-\vv{x}_i)$.  It also determines the inertial response through $\vv{p} = \gamma m\vv{v}$.
  \item \textbf{Electric charge $q$} sources the EM field via $\rho_q = \sum_i q_i\,\delta(\vv{x}-\vv{x}_i)$ and $\vv{J}_q = \sum_i q_i\vv{v}_i\,\delta(\vv{x}-\vv{x}_i)$.  It determines the response to the electromagnetic Lorentz force $q(\vv{E}+\vv{v}\times\vv{B})$.
\end{itemize}
A particle with both $m\neq 0$ and $q\neq 0$ contributes simultaneously to all four potentials $(\Phig,\Ag,\phi,\vv{A})$ and experiences forces from all four fields.

\subsection{Relativistic Particle Integration}

Particles are integrated relativistically via the 3-momentum $\vv{p} = \gamma m\vv{v}$.  The basic update scheme is:
\begin{equation}
  \vv{p} \;\leftarrow\; \vv{p} + \vv{F}_{\mathrm{total}}\,\dd t, \qquad
  \vv{v} = \frac{\vv{p}}{\sqrt{m^2 + |\vv{p}|^2/c^2}}, \qquad
  \vv{x} \;\leftarrow\; \vv{x} + \vv{v}\,\dd t.
\end{equation}
This correctly handles the anisotropy of relativistic inertia (longitudinal mass $\gamma^3 m$ vs.\ transverse mass $\gamma m$).  The force $\vv{F}_{\mathrm{total}}$ requires relativistic corrections at high velocities---see Section~\ref{sec:relcorr}.

\subsection{Electromagnetic Stress-Energy Coupling}

The total stress-energy tensor sourcing the GEM field is:
\begin{equation}
  T^{\mu\nu}_{\mathrm{total}} = T^{\mu\nu}_{\mathrm{matter}} + T^{\mu\nu}_{\mathrm{EM}},
\end{equation}
Since the simulation stores potentials rather than derived fields, it is useful to express $T^{\mu\nu}_{\mathrm{EM}}$ directly in terms of $\phi$ and $\vv{A}$ and their derivatives.  Define the shorthand combinations:
\begin{equation}
  \vv{P} \equiv \nabla\phi + \partial_t\vv{A}, \qquad
  \vv{C} \equiv \nabla\times\vv{A},
\end{equation}
so that $\vv{E} = -\vv{P}$ and $\vv{B} = \vv{C}$.  These are not stored as grid arrays; they are computed on demand from finite differences of the potential arrays.  The EM stress-energy components in terms of $\vv{P}$ and $\vv{C}$ are:
\begin{align}
  T^{00}_{\mathrm{EM}} &= \frac{\varepsilon_0}{2}\!\left(|\vv{P}|^2 + c^2|\vv{C}|^2\right), \label{eq:T00_pot}\\
  T^{0i}_{\mathrm{EM}} &= -\frac{1}{\mu_0 c}\!\left(\vv{P}\times\vv{C}\right)^i, \label{eq:T0i_pot}\\
  T^{ij}_{\mathrm{EM}} &= \varepsilon_0\!\left(-P^i P^j + c^2\sum_k C^{ik}C^{jk}\right) - \tfrac{1}{2}\delta^{ij}\varepsilon_0\!\left(|\vv{P}|^2+c^2|\vv{C}|^2\right), \label{eq:Tij_pot}
\end{align}
where $C^{ij} \equiv \partial_i A_j - \partial_j A_i$ is the antisymmetric derivative tensor (the spatial components of $F_{\mu\nu}$), so that $C^i = (\nabla\times\vv{A})^i = \tfrac{1}{2}\epsilon^{ijk}C_{jk}$.  The expression $\sum_k C^{ik}C^{jk}$ in Eq.~(\ref{eq:Tij_pot}) involves only antisymmetric combinations of first derivatives of $\vv{A}$ --- no pseudovectors appear.

Note that $\vv{C} = \nabla\times\vv{A}$ is the antisymmetric combination $(\partial_i A_j - \partial_j A_i)$ contracted with the Levi-Civita symbol, which is a component of the covariant field strength tensor $F_{\mu\nu}$.  It is not a grid array; it is evaluated on the fly from finite differences of edge potential values at the locations where it is needed.

The GEM effective sources use only $T^{00}$ and $T^{0i}$.  Using perturbation potentials only (to avoid double-counting the background metric; see Section~\ref{sec:energy}):
\begin{align}
  \rho_{\mathrm{EM}} &= \frac{T^{00}_{\mathrm{EM}}}{c^2}
    = \frac{\varepsilon_0}{2c^2}\!\left(|\nabla\phi^{\mathrm{pert}} + \partial_t\vv{A}^{\mathrm{pert}}|^2
      + c^2|\nabla\times\vv{A}^{\mathrm{pert}}|^2\right), \label{eq:rho_em_pot}\\
  \vv{J}_{\mathrm{EM}} &= \frac{\vv{S}_{\mathrm{EM}}}{c^2}
    = -\frac{1}{\mu_0 c^2}\!\left(\nabla\phi^{\mathrm{pert}} + \partial_t\vv{A}^{\mathrm{pert}}\right)
      \times\!\left(\nabla\times\vv{A}^{\mathrm{pert}}\right). \label{eq:J_em_pot}
\end{align}

In 2D, $\nabla\times\vv{A}^{\mathrm{pert}} = (\partial_x A_y^{\mathrm{pert}} - \partial_y A_x^{\mathrm{pert}})\hat{z}$ reduces to a single scalar $F_{12}^{\mathrm{pert}}$, and the cross product in Eq.~(\ref{eq:J_em_pot}) becomes a 90° rotation of the 2D vector $(\nabla\phi^{\mathrm{pert}}+\partial_t\vv{A}^{\mathrm{pert}})$ scaled by $F_{12}^{\mathrm{pert}}$.

The $T^{00}$ component means EM energy density gravitates.  The $T^{0i}$ component means EM energy flux (Poynting vector) acts as an effective mass current.  The spatial stress $T^{ij}$ is not needed for the GEM source but contributes to radiation pressure and EM momentum flux if those are tracked.

%====================================================================
\section{Two Implementation Approaches}
%====================================================================

\subsection{Method~1: Full Linearized Einstein--Maxwell}

This method integrates all 10 independent components of $\bar{h}^{\mu\nu}$ together with the EM wave equation in curved spacetime.

\subsubsection*{Per-timestep iteration}
\begin{enumerate}
  \item Assemble $T^{\mu\nu}_{\mathrm{matter}}$ from particle positions, velocities, and rest masses:
  \begin{equation}
    T^{00} = \gamma^2\rho c^2,\quad T^{0i} = \gamma^2\rho c\,v^i,\quad T^{ij} = \gamma^2\rho\,v^i v^j.
  \end{equation}

  \item Assemble $T^{\mu\nu}_{\mathrm{EM}}$ from current field values (Section~2.5).

  \item Deposit EM sources from charged particles:
  \begin{equation}
    \rho_q(\vv{x}) = \sum_i q_i\,\delta(\vv{x}-\vv{x}_i), \qquad
    \vv{J}_q(\vv{x}) = \sum_i q_i\vv{v}_i\,\delta(\vv{x}-\vv{x}_i).
  \end{equation}

  \item Sum: $T^{\mu\nu}_{\mathrm{total}} = T^{\mu\nu}_{\mathrm{matter}} + T^{\mu\nu}_{\mathrm{EM}}$.

  \item Step all 10 metric perturbation components:
  \begin{equation}
    \Box\,\bar{h}^{\mu\nu} = -\frac{16\pi G}{c^4}\,T^{\mu\nu}_{\mathrm{total}}.
  \end{equation}

  \item Step the EM wave equation in curved spacetime (with Ricci coupling):
  \begin{equation}
    \Box A^\mu + R^\mu{}_\nu A^\nu = \mu_0 J^\mu_q,
  \end{equation}
  where $R^\mu{}_\nu$ is the Ricci tensor from $\bar{h}^{\mu\nu}$.

  \item Extract fields from potentials:
  \begin{align}
    g^i  &= -\partial_i\bar{h}^{00}, &
    B_g^k &= \varepsilon^{ijk}\partial_i\bar{h}^{0j}, \\
    E^i  &= -\partial_t A^i - \partial^i A^0, &
    B^k  &= \varepsilon^{ijk}\partial_i A^j.
  \end{align}

  \item Compute geodesic acceleration using the electromagnetic field tensor $F^\mu{}_\nu$.  The Faraday tensor is defined as
  \begin{equation}
    F^{\mu\nu} = \partial^\mu A^\nu - \partial^\nu A^\mu,
    \qquad
    F^{\mu\nu} =
    \begin{pmatrix}
      0      & -E_x/c & -E_y/c & -E_z/c \\
      E_x/c  &  0     & -B_z   &  B_y   \\
      E_y/c  &  B_z   &  0     & -B_x   \\
      E_z/c  & -B_y   &  B_x   &  0
    \end{pmatrix},
  \end{equation}
  with $F^\mu{}_\nu = \eta_{\nu\alpha}F^{\mu\alpha}$.  The force term $(q/m)F^\mu{}_\nu U^\nu$ then reproduces the relativistic Lorentz force in the spatial components.  The geodesic equation is:
  \begin{equation}
    \frac{\dd U^\mu}{\dd\tau} = -\Gamma^\mu{}_{\alpha\beta}\,U^\alpha U^\beta + \frac{q}{m}F^\mu{}_\nu U^\nu.
  \end{equation}

  \item Update particle 4-velocities and positions.
\end{enumerate}

\subsection{Method~2: GEM + EM with Effective Sources}

This method replaces the 10-component tensor equation with two scalar/vector wave equations, absorbing $T^{\mu\nu}_{\mathrm{EM}}$ into effective sources.

\subsubsection*{Per-timestep iteration}
\begin{enumerate}
  \item Deposit matter sources:
  \begin{equation}
    \rho_{\mathrm{matter}} = \sum_i \gamma m_i\,\delta(\vv{x}-\vv{x}_i), \qquad
    \vv{J}_{\mathrm{matter}} = \sum_i \gamma m_i\vv{v}_i\,\delta(\vv{x}-\vv{x}_i).
  \end{equation}

  \item Deposit EM sources from charged particles.

  \item Compute EM effective GEM sources from perturbation potentials (see Eqs.~\ref{eq:rho_em_pot}--\ref{eq:J_em_pot}):
  \begin{align}
    \rho_{\mathrm{EM}} &= \frac{\varepsilon_0}{2c^2}\!\left(|\nabla\phi^{\mathrm{pert}} + \partial_t\vv{A}^{\mathrm{pert}}|^2
      + c^2|\nabla\times\vv{A}^{\mathrm{pert}}|^2\right), \\
    \vv{J}_{\mathrm{EM}} &= -\frac{1}{\mu_0 c^2}\!\left(\nabla\phi^{\mathrm{pert}} + \partial_t\vv{A}^{\mathrm{pert}}\right)
      \times\!\left(\nabla\times\vv{A}^{\mathrm{pert}}\right).
  \end{align}

  \item Sum effective sources: $\rho_{\mathrm{eff}} = \rho_{\mathrm{matter}} + \rho_{\mathrm{EM}}$,\; $\vv{J}_{\mathrm{eff}} = \vv{J}_{\mathrm{matter}} + \vv{J}_{\mathrm{EM}}$.

  \item Step GEM wave equations:
  \begin{equation}
    \Box\,\Phig = -4\pi G\rho_{\mathrm{eff}}, \qquad
    \Box\,\Ag   = -\frac{4\pi G}{c^2}\vv{J}_{\mathrm{eff}}.
  \end{equation}

  \item Step EM wave equation with $c_{\mathrm{local}}(\vv{x}) = c\!\left(1 + 2\Phig(\vv{x})/c^2\right)$:
  \begin{equation}
    \Box_{c_{\mathrm{local}}} A^\mu = \mu_0 J^\mu_q.
  \end{equation}

  \item Compute total force on each particle directly from potentials via Eqs.~(\ref{eq:em_force_potential})--(\ref{eq:gem_force_potential}), plus on-the-fly curl for Boris rotation (Section~\ref{sec:leapfrog}).  No full-grid field extraction is performed.

  \item Update momentum and position:
  \begin{equation}
    \vv{p} \leftarrow \vv{p} + \vv{F}_{\mathrm{total}}\,\dd t, \qquad
    \vv{v} = \frac{\vv{p}}{\sqrt{m^2+|\vv{p}|^2/c^2}}, \qquad
    \vv{x} \leftarrow \vv{x} + \vv{v}\,\dd t.
  \end{equation}
\end{enumerate}

%====================================================================
\section{Comparison of the Two Methods}
%====================================================================

\subsection{Field equations solved}

\begin{center}
\begin{tabular}{@{}L{3.2cm}L{5.0cm}L{5.0cm}@{}}
\toprule
\textbf{Quantity} & \textbf{Method~1 (Full Tensor)} & \textbf{Method~2 (GEM/EM)} \\
\midrule
Metric perturbation & All 10 components of $\bar{h}^{\mu\nu}$ & 2 potentials: $\Phig$, $\Ag$ \\[4pt]
EM field & $\Box A^\mu + R^\mu{}_\nu A^\nu = \mu_0 J^\mu_q$ & $\Box A^\mu = \mu_0 J^\mu_q$ with $c_{\mathrm{local}}$ \\[4pt]
EM--gravity coupling & Full $T^{\mu\nu}_{\mathrm{EM}}$ & $\rho_{\mathrm{EM}}=u_{\mathrm{EM}}/c^2$, $\vv{J}_{\mathrm{EM}}=\vv{S}/c^2$ only \\[4pt]
Particle equation & Geodesic via $\Gamma^\mu{}_{\alpha\beta}$ and $F^\mu{}_\nu$ & 3-force with relativistic corrections \\
\bottomrule
\end{tabular}
\end{center}

\subsection{Physical effects captured}

\begin{center}
\begin{tabular}{@{}L{4.5cm}L{3.0cm}L{5.5cm}@{}}
\toprule
\textbf{Effect} & \textbf{Method~1} & \textbf{Method~2} \\
\midrule
Newtonian gravity                     & Yes & Yes \\
Frame dragging (Lense--Thirring)      & Yes & Yes \\
Gravitational waves (tensor modes)    & Yes---all polarizations & Partial---scalar/vector only \\
Gravitomagnetic Lorentz force         & Yes & Yes (factor of 4 included) \\
EM energy gravitates                  & Yes---via $T^{00}_{\mathrm{EM}}$ & Yes---via $\rho_{\mathrm{EM}}=u_{\mathrm{EM}}/c^2$ \\
Poynting flux as mass current         & Yes---via $T^{0i}_{\mathrm{EM}}$ & Yes---via $\vv{J}_{\mathrm{EM}}=\vv{S}/c^2$ \\
Maxwell stress tensor sourcing metric & Yes---via $T^{ij}_{\mathrm{EM}}$ & No---dropped \\
Gravitational lensing                 & Exact---via Ricci coupling & Approximate---via $c_{\mathrm{local}}$ \\
Shapiro delay                         & Yes & Yes---via $c_{\mathrm{local}}$ \\
Relativistic inertia anisotropy       & Yes---geodesic & Yes---$p$-update \\
Velocity-dependent grav.\ force       & Yes---full $\Gamma^\mu{}_{\alpha\beta}U^\alpha U^\beta$ & Yes---Section~\ref{sec:relcorr} corrections \\
Correct photon deflection (factor 2)  & Yes & Yes---with Section~\ref{sec:relcorr} \\
Particles with both $m$ and $q$       & Yes & Yes \\
Intrinsic spin                        & Yes---MPD equations & Yes---Section~\ref{sec:spin} \\
\bottomrule
\end{tabular}
\end{center}

\subsection{Approximations in Method~2}

Method~2 is valid when:
\begin{itemize}
  \item Velocities are small compared to $c$ without Section~\ref{sec:relcorr} corrections; with those corrections, relativistic speeds are well handled.
  \item EM fields are not so strong that $T^{ij}_{\mathrm{EM}}$ contributes significantly to the spatial metric $h^{ij}$.
  \item The Ricci correction $R^\mu{}_\nu A^\nu$ is small---weak-field lensing is approximated by $c_{\mathrm{local}}$.
  \item The EM field is not too directional: $|\vv{S}|/c \ll u_{\mathrm{EM}}$.
\end{itemize}

\subsection{Coupling architecture}

\begin{eqbox}
\begin{equation*}
\begin{aligned}
&\text{GEM solver:}\quad \Box\bar{h} = \mathrm{source}(\rho_{\mathrm{eff}},\,\vv{J}_{\mathrm{eff}}) \\
&\hspace{4cm}\uparrow \\
&\hspace{2cm} \rho_{\mathrm{matter}} + u_{\mathrm{EM}}/c^2 \\
&\hspace{2cm} \vv{J}_{\mathrm{matter}} + \vv{S}_{\mathrm{Poynting}}/c^2 \\[6pt]
&\text{EM solver:}\quad \Box A = \mathrm{source}(\vv{J}_q),\quad c\to c_{\mathrm{local}} \\[6pt]
&\text{Particles:}\quad \dd\vv{p}/\dd t = \vv{F}_{\mathrm{GEM}} + \vv{F}_{\mathrm{EM}} + \vv{F}_{\mathrm{rel}} + \vv{F}_{\mathrm{spin}}
\end{aligned}
\end{equation*}
\end{eqbox}

%====================================================================
\section{Particles with Intrinsic Spin}
\label{sec:spin}
%====================================================================

\subsection{Overview}

A particle with intrinsic angular momentum (spin) obeys the Mathisson--Papapetrou--Dixon (MPD) equations rather than the geodesic equation.  In the GEM/EM approximation these reduce to a tractable set of 3-vector equations integrated alongside the momentum update.

Each spinning particle carries two additional properties beyond $m$ and $q$:
\begin{itemize}
  \item Spin 3-vector $\vv{S}$, with $|\vv{S}|$ conserved in the absence of radiation damping.
  \item Electromagnetic $g$-factor $g_s$ ($g_s=2$ for a Dirac fermion, $g_s=1$ for a classical spinning charged sphere).
\end{itemize}

\subsection{The spin 4-vector constraint}

The spin is carried by a 4-vector $S^\mu$ subject to the Tulczyjew--Dixon supplementary condition:
\begin{equation}
\label{eq:td}
  S^{\mu\nu} p_\nu = 0,
\end{equation}
where $S^{\mu\nu}$ is the antisymmetric spin tensor and $p^\mu$ is the \emph{4-momentum}, not the 4-velocity $U^\mu$.  This is an important distinction.  For a free particle in flat spacetime $p^\mu = mU^\mu$ and the two are proportional, so it makes no difference.  But in the MPD equations the canonical momentum receives a spin-curvature correction:
\begin{equation}
  p^\mu = mU^\mu - \frac{1}{2}R^\mu{}_{\nu\alpha\beta}U^\nu S^{\alpha\beta} + \ldots
\end{equation}
so $p^\mu$ and $mU^\mu$ are not exactly parallel when the particle is spinning in curved spacetime.

The condition $S^{\mu\nu}p_\nu = 0$ implies that in the lab frame:
\begin{equation}
  S^0 = \frac{\vv{S}\cdot\vv{p}}{\gamma mc},
\end{equation}
where $\vv{p} = \gamma m\vv{v}$ is the relativistic 3-momentum.  The timelike component $S^0$ is not an independent degree of freedom --- it is recomputed from $\vv{S}$ and $\vv{p}$ whenever needed and never evolved separately.  Using $\vv{p}$ rather than $m\vv{v}$ in this expression is exact; no linearized approximation is required.

\subsection{Constraint correction}

After each timestep $\vv{p}$ and $\vv{S}$ are updated independently, and $\vv{S}$ may acquire a spurious component parallel to $\vv{p}$ that violates the rest-frame orthogonality condition.  The 3+1 split of Eq.~(\ref{eq:td}) with $p^\mu = (E/c,\,\vv{p})$ and $S^\mu = (S^0,\,\vv{S})$ gives:
\begin{equation}
  S^0 = \frac{\vv{S}\cdot\vv{p}}{E/c} = \frac{\vv{S}\cdot\vv{p}}{\gamma mc},
\end{equation}
where $E = \gamma mc^2$ is the particle energy.  Since $S^0$ is never evolved independently but always reconstructed from $\vv{S}$ and $\vv{p}$, the constraint is automatically satisfied for $S^0$.  The drift accumulates entirely as a spurious component of $\vv{S}$ parallel to $\vv{p}$, which is removed by the covariant projection $S^\mu \leftarrow S^\mu - (S^\nu p_\nu/p^\alpha p_\alpha)\,p^\mu$.  Using $p^\alpha p_\alpha = -m^2c^2$ and $|\vv{p}|^2 + m^2c^2 = \gamma^2m^2c^2$, the spatial components give:
\begin{equation}
\label{eq:spin_correction}
  \vv{S} \;\leftarrow\; \vv{S} - \frac{\vv{S}\cdot\vv{p}}{|\vv{p}|^2 + m^2c^2}\,\vv{p},
  \qquad \vv{p} = \gamma m\vv{v}.
\end{equation}
This is the exact, covariant, single-formula correction requiring no threshold switching.  In the nonrelativistic limit $\vv{p}\approx m\vv{v}$, $\gamma\approx 1$:
\begin{equation}
  \vv{S} \leftarrow \vv{S} - \frac{\vv{S}\cdot\vv{v}}{c^2}\,\vv{v},
\end{equation}
confirming consistency.  The momentum $\vv{p} = \gamma m\vv{v}$ is available directly from the leapfrog integrator at half-integer timesteps.

\subsection{Spin-curvature and spin-field forces}

In 3D the spin gradient forces involve the gradient of $\vv{B}_g = \nabla\times\Ag$ and $\vv{B} = \nabla\times\vv{A}$ dotted with the spin and magnetic dipole moment:
\begin{align}
  \vv{F}_{\mathrm{spin,grav}} &= \nabla(\vv{S}\cdot\nabla\times\Ag), \\
  \vv{F}_{\mathrm{spin,EM}}   &= \nabla(\boldsymbol{\mu}\cdot\nabla\times\vv{A}), \qquad \boldsymbol{\mu} = \frac{g_s q}{2m}\vv{S}.
\end{align}
These gradients are evaluated at the particle position using the adjoint two-tent interpolation (Section~\ref{sec:deposition}).  In 2D $\nabla\times\Ag = B_{g,z}\hat{z}$ and $\nabla\times\vv{A} = B_z\hat{z}$ are scalars, and the spin is also scalar $s$, so the forces reduce to $s\,\nabla B_{g,z}$ and $\mu\,\nabla B_z$.

\subsection{Spin precession}

\begin{align}
  \boldsymbol{\Omega}_{\mathrm{grav}}   &= \nabla\times\Ag, \\
  \boldsymbol{\Omega}_{\mathrm{EM}}     &= \frac{g_s q}{2m}\nabla\times\vv{A}, \\
  \boldsymbol{\Omega}_{\mathrm{Thomas}} &= -\frac{\gamma^2}{\gamma+1}\,\frac{\vv{v}\times\vv{a}}{c^2},
\end{align}
\begin{equation}
  \frac{\dd\vv{S}}{\dd t} = \vv{S}\times\!\left(\boldsymbol{\Omega}_{\mathrm{grav}} + \boldsymbol{\Omega}_{\mathrm{EM}} + \boldsymbol{\Omega}_{\mathrm{Thomas}}\right).
\end{equation}
The quantities $\nabla\times\Ag$ and $\nabla\times\vv{A}$ are evaluated on the fly at the particle position from surrounding edge potentials via the Yee curl (Eq.~\ref{eq:yee_curl}), not stored as grid arrays.

\subsection{Spin as a field source}

A spinning charged particle carries magnetic dipole moment $\boldsymbol{\mu} = g_s(q/2m)\vv{S}$ and a spinning massive particle carries gravitomagnetic dipole moment $\boldsymbol{\sigma} = 2\vv{S}$.  These contribute magnetization currents to the source deposition:
\begin{align}
  \vv{J}_q    &\mathrel{+}= \nabla\times\vv{M}(\vv{x}), \quad \vv{M}(\vv{x}) = \sum_i \boldsymbol{\mu}_i\,\delta(\vv{x}-\vv{x}_i), \\
  \vv{J}_{\mathrm{grav}} &\mathrel{+}= \nabla\times\boldsymbol{\Sigma}(\vv{x}), \quad \boldsymbol{\Sigma}(\vv{x}) = \sum_i \boldsymbol{\sigma}_i\,\delta(\vv{x}-\vv{x}_i).
\end{align}
The $1/r^3$ near-field and $1/r$ radiation fall-offs of dipole fields emerge automatically from the d'Alembert operator.

\subsection{Radiation from precessing dipoles}

\begin{itemize}
  \item \textbf{Conservative:} ignore radiation reaction; $|\vv{S}|$ is conserved exactly.
  \item \textbf{Phenomenological damping:} $\dd\vv{S}/\dd t = \vv{S}\times\boldsymbol{\Omega}_{\mathrm{total}} - \alpha\vv{S}$.
\end{itemize}

\subsection{Complete algorithm including spin (Method~2)}
\label{sec:alg_spin}

The following enumeration gives the complete per-timestep algorithm for Method~2 including all spin effects developed so far, before adding the relativistic force corrections of Section~\ref{sec:relcorr}.

\begin{algsteps}
  \item \textbf{Deposit matter and charge sources.}
  \begin{equation}
    \rho_{\mathrm{matter}} = \sum_i \gamma m_i\,W(\vv{x}-\vv{x}_i), \quad
    \vv{J}_{\mathrm{matter}} = \sum_i \gamma m_i\vv{v}_i\,W(\vv{x}-\vv{x}_i),
  \end{equation}
  \begin{equation}
    \rho_q = \sum_i q_i\,W(\vv{x}-\vv{x}_i), \quad
    \vv{J}_q = \sum_i q_i\vv{v}_i\,W(\vv{x}-\vv{x}_i),
  \end{equation}
  where $W$ is the cloud-in-cell (CIC) kernel approximating $\delta(\vv{x}-\vv{x}_i)$ (see Section~\ref{sec:numerical}).

  \item \textbf{Deposit spin dipole sources.}
  \begin{align}
    \boldsymbol{\mu}_i   &= \frac{g_{s,i} q_i}{2m_i}\vv{S}_i, \qquad
    \boldsymbol{\sigma}_i = 2\vv{S}_i, \\
    \vv{J}_q             &\mathrel{+}= \nabla\times\!\left[\sum_i \boldsymbol{\mu}_i\,W_D(\vv{x}-\vv{x}_i)\right], \\
    \vv{J}_{\mathrm{grav}} &\mathrel{+}= \nabla\times\!\left[\sum_i \boldsymbol{\sigma}_i\,W_D(\vv{x}-\vv{x}_i)\right],
  \end{align}
  where $W_D$ is the quadratic B-spline kernel (see Section~\ref{sec:numerical}) used for dipole deposition.

  \item \textbf{Compute EM effective GEM sources} from perturbation potentials (Eqs.~\ref{eq:rho_em_pot}--\ref{eq:J_em_pot}):
  \begin{align}
    \rho_{\mathrm{EM}} &= \frac{\varepsilon_0}{2c^2}\!\left(|\nabla\phi^{\mathrm{pert}} + \partial_t\vv{A}^{\mathrm{pert}}|^2
      + c^2|\nabla\times\vv{A}^{\mathrm{pert}}|^2\right), \\
    \vv{J}_{\mathrm{EM}} &= -\frac{1}{\mu_0 c^2}\!\left(\nabla\phi^{\mathrm{pert}} + \partial_t\vv{A}^{\mathrm{pert}}\right)
      \times\!\left(\nabla\times\vv{A}^{\mathrm{pert}}\right).
  \end{align}

  \item \textbf{Sum effective sources.}
  \begin{equation}
    \rho_{\mathrm{eff}} = \rho_{\mathrm{matter}} + \rho_{\mathrm{EM}}, \qquad
    \vv{J}_{\mathrm{eff}} = \vv{J}_{\mathrm{matter}} + \vv{J}_{\mathrm{EM}}.
  \end{equation}

  \item \textbf{Step GEM wave equations.}
  \begin{equation}
    \Phig^{n+1} = 2\Phig^n - \Phig^{n-1} + (c\,\dd t)^2\!\left[\nabla^2\Phig^n - 4\pi G\rho_{\mathrm{eff}}^n\right],
  \end{equation}
  and similarly for each component of $\Ag$.

  \item \textbf{Step EM wave equations} using $c_{\mathrm{local}}^n(\vv{x}) = c(1+2\Phig^n(\vv{x})/c^2)$:
  \begin{equation}
    \phi^{n+1} = 2\phi^n - \phi^{n-1} + (c_{\mathrm{local}}\,\dd t)^2\!\left[\nabla^2\phi^n + \rho_q^n/\varepsilon_0\right],
  \end{equation}
  and similarly for each component of $\vv{A}$.

  \item \textbf{Compute force on each particle} from potentials directly (Eqs.~\ref{eq:force_gem_pot}--\ref{eq:force_em_pot}).  All terms involve only spatial and temporal finite differences of the potential arrays at and near the particle position, interpolated to the particle using the CIC weights.  Additionally compute $\nabla\times\Ag$ and $\nabla\times\vv{A}$ on-the-fly at the particle position for the Boris rotation and spin forces, using the Yee curl stencil evaluated at surrounding cell centers and bilinearly interpolated.  The full combined force including spin gradient terms is:
  \begin{align}
    \vv{F} &= m\!\left(-\nabla\Phig - \partial_t\Ag + 4\nabla(\vv{v}\cdot\Ag) - 4(\vv{v}\cdot\nabla)\Ag\right) \nonumber \\
           &+ q\!\left(-\nabla\phi - \partial_t\vv{A} + \nabla(\vv{v}\cdot\vv{A}) - (\vv{v}\cdot\nabla)\vv{A}\right) \nonumber \\
           &+ \nabla(\vv{S}\cdot\nabla\times\Ag) + \nabla(\boldsymbol{\mu}\cdot\nabla\times\vv{A}),
  \end{align}
  where the gradients of $\nabla\times\Ag$ and $\nabla\times\vv{A}$ for the spin gradient forces use the adjoint two-tent interpolation (Section~\ref{sec:deposition}).

  \item \textbf{Update momentum and position} (Boris pusher; see Section~\ref{sec:numerical}):
  \begin{equation}
    \vv{p} \leftarrow \vv{p} + \vv{F}\,\dd t, \qquad
    \vv{v} = \frac{\vv{p}}{\sqrt{m^2+|\vv{p}|^2/c^2}}, \qquad
    \vv{x} \leftarrow \vv{x} + \vv{v}\,\dd t.
  \end{equation}
  Record $\vv{p}_{\mathrm{old}}$ and $\vv{v}_{\mathrm{old}}$ before this step for use in Step~10.

  \item \textbf{Update spin.}
  \begin{align}
    \vv{a}                      &= (\vv{v} - \vv{v}_{\mathrm{old}})/\dd t, \\
    \boldsymbol{\Omega}_{\mathrm{Thomas}} &= -\frac{\gamma^2}{\gamma+1}\frac{\vv{v}\times\vv{a}}{c^2}, \\
    \boldsymbol{\Omega}_{\mathrm{total}} &= (\nabla\times\Ag)^{\mathrm{particle}} + \frac{g_s q}{2m}(\nabla\times\vv{A})^{\mathrm{particle}} + \boldsymbol{\Omega}_{\mathrm{Thomas}}, \\
    \vv{S}                      &\leftarrow \vv{S} + (\vv{S}\times\boldsymbol{\Omega}_{\mathrm{total}})\,\dd t.
  \end{align}
  where $(\nabla\times\Ag)^{\mathrm{particle}}$ and $(\nabla\times\vv{A})^{\mathrm{particle}}$ are the on-the-fly Yee curl values computed at the particle position from surrounding edge potentials (already available from the force evaluation step).
  Apply constraint correction (Eq.~\ref{eq:spin_correction}) and renormalize spin magnitude:
  \begin{equation}
    \vv{S} \leftarrow \vv{S} - \frac{\vv{S}\cdot\vv{p}}{|\vv{p}|^2 + m^2c^2}\,\vv{p}, \qquad
    \vv{S} \leftarrow \vv{S}\,\frac{|\vv{S}_0|}{|\vv{S}|}.
  \end{equation}
\end{algsteps}

%====================================================================
\section{Relativistic Force Corrections}
\label{sec:relcorr}
%====================================================================

\subsection{The problem with GEM at high velocities}

The GEM Lorentz force is derived under the slow-motion assumption.  At relativistic speeds the spatial metric perturbation $h^{ij}$ also contributes to particle deflection.  The clearest illustration is photon deflection: the correct GR prediction is twice the Newtonian (GEM) value, with the extra factor coming from spatial Christoffel terms that GEM drops.

\subsection{The linearized metric and its two metric coefficients}

The full linearized metric is:
\begin{equation}
  \dd s^2 = -\!\left(1+\frac{2\Phig}{c^2}\right)c^2\,\dd t^2
            + \left(1-\frac{2\Phig}{c^2}\right)\!\left(\dd x^2+\dd y^2+\dd z^2\right).
\end{equation}
The two coefficients act in opposite directions:
\begin{align}
  \text{Time:}  &\quad \left(1+\frac{2\Phig}{c^2}\right) < 1 \text{ near mass} \;\Rightarrow\; \text{clocks run slow.} \\
  \text{Space:} &\quad \left(1-\frac{2\Phig}{c^2}\right) > 1 \text{ near mass} \;\Rightarrow\; \text{rulers stretch.}
\end{align}
For a slow particle only the time component matters; the spatial component is negligible.  For a photon both contribute equally, doubling the deflection relative to the Newtonian prediction.

\subsection{The Shapiro delay}
\label{sec:shapiro}

A photon follows a null geodesic ($\dd s^2=0$).  Solving for the coordinate speed of light:
\begin{equation}
  \label{eq:c_local}
  \frac{\dd x}{\dd t} \approx c\!\left(1 + \frac{2\Phig}{c^2}\right) \equiv c_{\mathrm{local}}.
\end{equation}
The factor of 2 in $c_{\mathrm{local}}$ comes from both metric coefficients contributing to the null condition.  The Shapiro delay is:
\begin{equation}
  \label{eq:shapiro_delay}
  \Delta t_{\mathrm{Shapiro}} = -\frac{2}{c^3}\int_{\mathrm{path}}\Phig\,\dd\ell.
\end{equation}
This is implemented in the simulator by stepping the EM wave equation with $c_{\mathrm{local}}$ rather than $c$.  The perturbation of photon paths (gravitational lensing) enters through the trajectory integration, where the same $c_{\mathrm{local}}$ modification of the wavefronts causes bending.

\subsection{Expansion of the geodesic equation}
\label{sec:geodesic_expansion}

The relativistic correction terms used here are those of the
\textbf{Einstein--Infeld--Hoffmann (EIH) equation}~\cite{Will1993,AliHaimoud2019}, the standard
post-Newtonian (1PN) equation of motion for self-gravitating particles
in harmonic gauge.  For a test particle in a (possibly moving) gravitational
source, EIH may be written
\begin{equation}
\label{eq:eih_full}
  \frac{\dd\vv{v}}{\dd t}
  = -(1 + v^2/c^2 + 4\Phig/c^2)\nabla\Phig
    - \partial_t\Ag
    + \vv{v}\times(\nabla\times\Ag)
    + (3\partial_t\Phig + 4\vv{v}\cdot\nabla\Phig)\vv{v}/c^2,
\end{equation}
which is correct through $\order{(v/c)^2}$ in the coordinate acceleration.
The gravitomagnetic vector potential $\Ag$ itself scales as
$|\Ag| \sim (v/c)^3$ at the source level (since
$\Delta\Ag \propto T^{0i} \sim \rho v$), and the terms containing
$\partial_t\Ag$ and $\vv{v}\times(\nabla\times\Ag)$ contribute to
$\dd\vv{v}/\dd t$ at the same 1PN order as the $(v^2/c^2)\nabla\Phig$ and
$\Phig\nabla\Phig$ corrections; the Lense--Thirring frame-dragging effect
arises from these gravitomagnetic terms.

Expressed in terms of the local gravitational vector
$\vv{g} = -\nabla\Phig - \partial_t\Ag$ and reorganised by order in $v/c$:
\begin{equation}
\label{eq:geodesic_expansion}
  a^i = \underbrace{g^i}_{\order{1}}
      + \underbrace{4\varepsilon^{ijk}v_j B_{g,k}}_{\order{v/c}\text{ (Lense--Thirring)}}
      + \underbrace{\frac{v^2}{c^2}g^i + \frac{4\Phig}{c^2}g^i - \frac{4(\vv{v}\cdot\vv{g})v^i}{c^2}}_{\order{v^2/c^2}\text{ (EIH 1PN)}}
      + \order{v^3/c^3},
\end{equation}
where $B_{g,k} = (\nabla\times\Ag)_k$ is evaluated on-the-fly from the
potentials.  In implementation the cross-product $\varepsilon^{ijk}v_j B_{g,k}$
is replaced by its potential form $\partial_i(\vv{v}\cdot\Ag) - (\vv{v}\cdot\nabla)A_{g,i}$
via the vector identity, avoiding explicit computation of $\vv{B}_g$ as a
grid array (Section~\ref{sec:force_potential}).

\textit{Note (v34 correction vs.\ v33).}  v33 of this document had the
velocity-coupling term with a $+$ sign, $+4(\vv{v}\cdot\vv{g})v^i/c^2$, and
omitted the $4\Phig\vv{g}/c^2$ ``Shapiro'' radial term.  Both have been
corrected here against the canonical EIH derivation
(e.g.\ \cite{AliHaimoud2019} eq.~37; \cite{Will1993} \S6.2).  The sign error was confirmed
empirically: with the v33 sign the Stage~8 perihelion test precesses
\emph{retrograde}; with the EIH sign and the Shapiro term included the test
recovers the Schwarzschild precession to better than $1\%$ in the
deep-1PN regime, after subtracting a small ``kinematic'' baseline that arises
from running relativistic 3-momentum dynamics with a non-relativistic force
law (Section~\ref{sec:stage_precession}).  The Shapiro radial term
originates not from the spatial metric coefficient alone but from the
$\order{v^4/L}$ correction to $\Gamma^i_{00}$, which pulls in the second-order
expansion of $g_{00}$.

The $\order{v^2/c^2}$ terms produce three physically distinct effects:
\begin{itemize}
  \item $(v^2/c^2)\vv{g}$: a fast particle experiences stronger
    gravitational deflection.  Combined with the spatial-metric
    contribution to $c_{\mathrm{local}}$, the factor-of-2 photon-deflection
    result is recovered in the $v\to c$ limit.
  \item $(4\Phig/c^2)\vv{g}$: an additional radial enhancement of the
    Newtonian force that vanishes far from the source.  Contributes to the
    static binding energy and to perihelion precession.
  \item $-4(\vv{v}\cdot\vv{g})\vv{v}/c^2$: a velocity-coupling force along
    $\vv{v}$ that does no work over a closed orbit but breaks the
    Kepler-orbit closure condition.
\end{itemize}
Together these are the test-particle limit of the full EIH equation; the
combination, integrated over a bound orbit, reproduces the canonical
$\Delta\phi_{\mathrm{prec}} = 6\pi G_{\mathrm{eff}} M / (c_{\mathrm{eff}}^2\,a(1-e^2))$
Schwarzschild precession.

\subsection{Practical implementation tiers}

\begin{center}
\begin{tabular}{@{}lL{8cm}L{2cm}@{}}
\toprule
\textbf{Tier} & \textbf{Gravitational force expression (EIH 1PN basis)} & \textbf{Valid range} \\
\midrule
0: Newtonian & $-m\nabla\Phig$ & $v/c < 0.01$ \\
1: GEM       & $m\bigl(-\nabla\Phig - \partial_t\Ag + 4\nabla(\vv{v}\cdot\Ag) - 4(\vv{v}\cdot\nabla)\Ag\bigr)$ & $v/c < 0.1$ \\
2: EIH 1PN (test particle) & Tier~1 plus $(v^2/c^2)\vv{g} + (4\Phig/c^2)\vv{g} - 4(\vv{v}\cdot\vv{g})\vv{v}/c^2$ & $v/c < 0.5$ \\
3: Full EIH  & All $\order{(v/c)^2}$ acceleration terms; Lense--Thirring enters via $\Ag$ (which scales as $(v/c)^3$ at the source) & $v/c \leq 1$ \\
\bottomrule
\end{tabular}
\end{center}

Tier~2 in v34 differs from v33 by (i) the sign on the velocity-coupling
term and (ii) the addition of the $4\Phig\vv{g}/c^2$ Shapiro radial term;
see Section~\ref{sec:geodesic_expansion} for the derivation against the
EIH lecture~25 form.

The EM Lorentz force $q(-\nabla\phi - \partial_t\vv{A} + \nabla(\vv{v}\cdot\vv{A}) - (\vv{v}\cdot\nabla)\vv{A})$ is already exactly relativistic and requires no velocity-dependent corrections.  In all tiers the momentum update uses the relativistic $p$-update.

\subsection{Where each effect enters the simulator}

\begin{center}
\begin{tabular}{@{}L{4.5cm}L{4cm}L{4cm}@{}}
\toprule
\textbf{Effect} & \textbf{Source} & \textbf{Simulator entry point} \\
\midrule
Gravitational time dilation & $g_{00}$ expansion & Proper time update \\
Shapiro delay / light bending & Both metric coefficients via $c_{\mathrm{local}}$ & EM wave equation \\
Doubled photon deflection & Spatial metric $(1-2\Phig/c^2)$ & $c_{\mathrm{local}}$; $(v^2/c^2)\vv{g}$ for massive particles \\
Perihelion precession & EIH 1PN: combination of $(v^2/c^2)\vv{g}$, $(4\Phig/c^2)\vv{g}$, $-4(\vv{v}\cdot\vv{g})\vv{v}/c^2$ & Tier-2 force terms \\
Lense--Thirring frame dragging & Gravitomagnetic $\Ag\sim(v/c)^3$ at source & GEM Tier-1 ($\Ag$) potential terms \\
\bottomrule
\end{tabular}
\end{center}

\subsection{Complete algorithm including spin and relativistic corrections (Method~2)}
\label{sec:alg_rel}

This enumeration extends Section~\ref{sec:alg_spin} by adding the Tier~3 relativistic force corrections.  Steps~1--7 are unchanged.  Steps~8 and~10 are modified as shown.

\begin{algsteps}
  \item \textbf{Deposit matter and charge sources.} (As in Section~\ref{sec:alg_spin}, Step~1.)

  \item \textbf{Deposit spin dipole sources.} (As in Section~\ref{sec:alg_spin}, Step~2.)

  \item \textbf{Compute EM effective GEM sources.} (As before.)

  \item \textbf{Sum effective sources.} (As before.)

  \item \textbf{Step GEM wave equations.} (As before.)

  \item \textbf{Step EM wave equations} with $c_{\mathrm{local}} = c(1+2\Phig/c^2)$. (As before.)

  \item \textbf{Compute total force on each particle} from potentials directly.  The gravitomagnetic cross-product terms $\vv{v}\times\vv{B}_g$ and $\vv{v}\times\vv{B}$ are folded into the potential-form force via the identity $\vv{v}\times(\nabla\times\vv{A}) = \nabla(\vv{v}\cdot\vv{A}) - (\vv{v}\cdot\nabla)\vv{A}$.  The $\vv{g}$ appearing in the EIH 1PN correction terms is computed on-the-fly at the particle position as $\vv{g} = -\nabla\Phig - \partial_t\Ag$.  The $4\Phig\vv{g}/c^2$ (``Shapiro'') term uses the total $\Phig$ interpolated at the particle:
  \begin{eqbox}
  \begin{align*}
    \vv{F}_{\mathrm{grav}}  &= m\!\left[-\nabla\Phig - \partial_t\Ag
                                + 4\nabla(\vv{v}\cdot\Ag) - 4(\vv{v}\cdot\nabla)\Ag
                                + \frac{v^2}{c^2}\vv{g}
                                + \frac{4\Phig}{c^2}\vv{g}
                                - \frac{4(\vv{v}\cdot\vv{g})\vv{v}}{c^2}\right], \\[4pt]
    \vv{F}_{\mathrm{EM}}    &= q\!\left[-\nabla\phi - \partial_t\vv{A}
                                + \nabla(\vv{v}\cdot\vv{A}) - (\vv{v}\cdot\nabla)\vv{A}\right], \\[4pt]
    \vv{F}_{\mathrm{spin}}  &= \nabla(\vv{S}\cdot\nabla\times\Ag)
                              + \nabla(\boldsymbol{\mu}\cdot\nabla\times\vv{A}), \\[4pt]
    \vv{F}_{\mathrm{total}} &= \vv{F}_{\mathrm{grav}} + \vv{F}_{\mathrm{EM}} + \vv{F}_{\mathrm{spin}}.
  \end{align*}
  \end{eqbox}
  The bracketed expression in $\vv{F}_{\mathrm{grav}}$ is the test-particle limit of the EIH equation (Eq.~\ref{eq:eih_full}) written in terms of $\vv{g}$ instead of $\nabla\Phig$; v33 had the sign of the velocity-coupling term reversed and was missing the $4\Phig\vv{g}/c^2$ Shapiro term (Section~\ref{sec:geodesic_expansion}).
  The quantities $\nabla\times\Ag$ and $\nabla\times\vv{A}$ are evaluated on-the-fly at the particle position from surrounding edge potentials via the Yee curl (Eq.~\ref{eq:yee_curl}).  Evaluate $v^2/c^2$ and $\vv{v}\cdot\vv{g}$ using $\vv{v}_{n-1/2}$ from the previous half-step.

  \item \textbf{Update momentum and position} (Boris pusher; see Section~\ref{sec:numerical}):
  \begin{equation}
    \vv{p} \leftarrow \vv{p} + \vv{F}_{\mathrm{total}}\,\dd t, \qquad
    \vv{v} = \frac{\vv{p}}{\sqrt{m^2+|\vv{p}|^2/c^2}}, \qquad
    \vv{x} \leftarrow \vv{x} + \vv{v}\,\dd t.
  \end{equation}
  Record $\vv{v}_{\mathrm{old}}$ before the update for Step~10.

  \item \textbf{Update spin.} (As in Section~\ref{sec:alg_spin}, Step~10, using updated $\vv{v}$ and $\vv{a} = (\vv{v}-\vv{v}_{\mathrm{old}})/\dd t$.)
\end{algsteps}

%====================================================================
\section{Proper Time Clocks}
\label{sec:proptime}
%====================================================================

\subsection{Overview}

Each particle carries a proper time clock $\tau_i$ accumulating the elapsed proper time along its worldline.  The clock requires no new field equations and adds negligible computational cost---one square root per particle per timestep.

\subsection{Proper time in the linearized metric}

For a non-rotating background the proper time rate is determined by the gravitoelectric potential and the particle velocity:
\begin{equation}
  \frac{\dd\tau}{\dd t} = \sqrt{\left(1+\frac{2\Phig}{c^2}\right) - \left(1-\frac{2\Phig}{c^2}\right)\frac{v^2}{c^2}}.
\end{equation}
For a rotating background the off-diagonal metric components contribute an
additional term.  To fix the cross-term coefficient consistently with the
doc's force-law convention --- the GEM Lorentz force has a factor of 4 from
the spin-2 nature of gravity, $\vv{F}/m \supset 4\vv{v}\times(\nabla\times\Ag)$
--- the linearized metric perturbation must be $h_{0i} = 4A_{g,i}/c$, not
the $2A_{g,i}/c$ written in v33 of this document (which was inconsistent with
its own force law).  See the note at the end of this subsection for the
v34 corrections.  The cross term therefore contributes
$g_{0i}\dd x^i\dd t = (4A_{g,i}/c)v^i\dd t^2$ to $\dd s^2$, and the full
linearized expression for $\dd\tau/\dd t$ becomes:
\begin{equation}
\label{eq:proper_time_full}
  \frac{\dd\tau}{\dd t} = \sqrt{\left(1+\frac{2\Phig}{c^2}\right) - \left(1-\frac{2\Phig}{c^2}\right)\frac{v^2}{c^2} - \frac{8\,\Ag\cdot\vv{v}}{c^2}}.
\end{equation}
The $-8\Ag\cdot\vv{v}/c^2$ term has opposite sign for prograde and retrograde orbits around a spinning mass (since $\vv{v}$ reverses but $\Ag$ does not), producing a proper time difference between co- and counter-rotating clocks at the same orbital radius --- the gravitomagnetic clock effect.  The negative sign means \emph{prograde} clocks run slower than retrograde clocks at the same orbital radius, consistent with the standard Cohen--Mashhoon result for Kerr~\cite{CohenMashhoon1993,Mashhoon2003}.  For a non-rotating background $\Ag = 0$ and Eq.~(\ref{eq:proper_time_full}) reduces to the simpler form.

To first order:
\begin{equation}
  \frac{\dd\tau}{\dd t} \approx 1 + \underbrace{\frac{\Phig}{c^2}}_{\text{gravitational}} - \underbrace{\frac{v^2}{2c^2}}_{\text{kinematic}} - \underbrace{\frac{4\,\Ag\cdot\vv{v}}{c^2}}_{\text{gravitomagnetic}}.
\end{equation}
The exact expression is used in the integrator; the first-order form is shown for interpretation.

\paragraph{Provenance and v34 correction.}  The linearized line element
underlying Eq.~(\ref{eq:proper_time_full}) is the one given in
\cite{AliHaimoud2019} eq.~41:
$\dd s^2 = -(1+2\phi+\ldots)\dd t^2 + 2\xi_i\,\dd t\,\dd x^i + (1-2\phi)\delta_{ij}\dd x^i\dd x^j$
in geometric units.  The lecture's gravitomagnetic potential $\xi_i$ equals
$4A_{g,i}/c$ in the doc's GEM-with-spin-2-factor-in-force convention (since
the lecture's geodesic equation eq.~37 carries coefficient~1 on
$\vv{v}\times(\nabla\times\xi)$ whereas this document's force law carries
coefficient~4 on $\vv{v}\times(\nabla\times\Ag)$).  Substituting
$\xi_i = 4A_{g,i}/c$ into $\dd\tau^2 = -\dd s^2$ and restoring $c$ gives
Eq.~(\ref{eq:proper_time_full}) directly.  For an explicit gravitomagnetic
proper-time formula in a different (factor-of-2-in-force) GEM convention
see \cite{Mashhoon2003} \S2; results agree up to an overall $\Ag$
normalization.

v33 of this document wrote $h_{0i} = 2A_{g,i}/c$ and a corresponding
$+(4\Ag\cdot\vv{v})/c^2$ inside the square root (linearized
$+(2\Ag\cdot\vv{v})/c^2$).  Both the magnitude and the sign were wrong;
the magnitude conflicted with v33's own force-law factor of~4, and the
sign would have predicted prograde clocks running \emph{faster} than
retrograde, opposite to the canonical clock-effect direction.

\subsection{Physical interpretation}

The kinematic term $-v^2/(2c^2)$ always slows the clock.  The gravitational term $\Phig/c^2$ also slows it in a well ($\Phig<0$), and speeds it up at high altitude ($\Phig$ less negative).  For GPS satellites the two terms compete: the satellite is high up (gravitational term speeds clock) but moving fast (kinematic term slows clock).  The gravitational effect wins, so GPS clocks run fast by $38\,\mu$s/day and must be corrected.

The gravitomagnetic term $-4\Ag\cdot\vv{v}/c^2$ is nonzero only around rotating masses.  It has opposite sign for co- and counter-rotating orbits because $\vv{v}$ reverses direction while $\Ag$ does not, and with the minus sign \emph{prograde} clocks accumulate proper time more slowly than retrograde clocks at the same orbital radius.  Integrating around a full orbit the difference in accumulated proper time between prograde ($+$) and retrograde ($-$) orbits is:
\begin{equation}
  \Delta\tau = \tau_+ - \tau_- = -\frac{8}{c^2}\oint\Ag\cdot\dd\vv{x} = -\frac{8}{c^2}\int(\nabla\times\Ag)\cdot\dd\vv{A} = -\frac{8}{c^2}\int\vv{B}_g\cdot\dd\vv{A},
\end{equation}
by Stokes' theorem --- the proper time difference is proportional to (minus) the gravitomagnetic flux through the orbit.  This is the gravitomagnetic clock effect, and it is a direct consequence of the same $\Ag$ that appears in the GEM force expression.  \cite{CohenMashhoon1993} gave the original Kerr derivation; \cite{Iorio2001} and \cite{Tartaglia2000} re-derive and discuss detectability.  Expressed in this document's $\Ag$ convention the prograde-minus-retrograde clock difference matches the formula above up to the overall $\Ag$ normalization (\cite{Mashhoon2003} expresses the same effect in his factor-of-2 GEM convention).

\subsection{Integration}

Using the total potential $\Phig^{\mathrm{total}} = \Phig^{\mathrm{bg}} + \Phig^{\mathrm{pert}}$ and total vector potential $\Ag^{\mathrm{total}} = \Ag^{\mathrm{bg}} + \Ag^{\mathrm{pert}}$, both already evaluated during force calculation:
\begin{equation}
  \dd\tau_i = \dd t\,\sqrt{1 + \frac{2\Phig}{c^2} - \left(1-\frac{2\Phig}{c^2}\right)\beta^2 - \frac{8\,\Ag\cdot\vv{v}}{c^2}},
  \qquad \beta^2 = v^2/c^2,
\end{equation}
\begin{equation}
  \tau_i \;\leftarrow\; \tau_i + \dd\tau_i.
\end{equation}
The quantity $\Ag\cdot\vv{v}$ is already computed as an intermediate result in the GEM force expression $4\nabla(\vv{v}\cdot\Ag) - 4(\vv{v}\cdot\nabla)\Ag$ --- specifically it is the scalar $\vv{v}\cdot\Ag$ whose gradient appears in the force.  No additional grid interpolation is required; the same interpolated value is reused.

In the leapfrog scheme, $\vv{v}$ lives at half-integer times.  A second-order accurate estimate at the integer step is $\vv{v}_n \approx \tfrac{1}{2}(\vv{v}_{n-1/2}+\vv{v}_{n+1/2})$.

\subsection{Pedagogical demonstrations}

\begin{center}
\begin{tabular}{@{}L{3.5cm}L{4cm}L{5.5cm}@{}}
\toprule
\textbf{Demonstration} & \textbf{Setup} & \textbf{What the student sees} \\
\midrule
Twin paradox & One particle stationary, one on an elliptical orbit. & $\tau_{\mathrm{traveler}} < \tau_{\mathrm{stationary}}$; traveling twin is younger. \\[4pt]
Gravitational time dilation & Two stationary particles at different radii. & Deeper clock ticks slower; purely the $\Phig/c^2$ term. \\[4pt]
GPS effect & Circular orbiter vs.\ stationary at same radius. & Two competing effects; exact cancellation at one specific radius. \\[4pt]
Clock in gravitational wave & Stationary particle, scalar wave passes through. & Proper time rate oscillates as $\Phig$ alternately deepens and shallows. \\[4pt]
Gravitomagnetic clock effect & Two particles on identical circular orbits in opposite directions around a spinning background mass. & Prograde and retrograde clocks accumulate different proper time per orbit; difference grows linearly with number of orbits and is proportional to the gravitomagnetic flux $\int\vv{B}_g\cdot\dd\vv{A}$ through the orbit. \\
\bottomrule
\end{tabular}
\end{center}

%====================================================================
\section{Background Spacetime and the Perturbation Split}
%====================================================================

\subsection{Motivation}

When one mass dominates the gravitational field, the dominant background can be precomputed once and held fixed; only the perturbation fields need to be time-stepped.  This is the background + perturbation split, used in black hole perturbation theory, restricted celestial mechanics, and QFT in curved spacetime.

\subsection{The decomposition}

\begin{equation}
  \Phig(\vv{x},t) = \Phig^{\mathrm{bg}}(\vv{x}) + \Phig^{\mathrm{pert}}(\vv{x},t),
\end{equation}
and similarly for $\Ag$, $\phi$, $\vv{A}$.  The total force on each particle is:
\begin{equation}
  \vv{F}_{\mathrm{total}} = \vv{F}(\vv{g}^{\mathrm{bg}}+\vv{g}^{\mathrm{pert}},\;\vv{B}_g^{\mathrm{bg}}+\vv{B}_g^{\mathrm{pert}},\;\vv{E}^{\mathrm{bg}}+\vv{E}^{\mathrm{pert}},\;\vv{B}^{\mathrm{bg}}+\vv{B}^{\mathrm{pert}}).
\end{equation}

\subsection{Analytic background configurations}

\begin{center}
\begin{tabular}{@{}L{3.5cm}L{4.5cm}L{5cm}@{}}
\toprule
\textbf{Type} & \textbf{Potentials} & \textbf{Notes} \\
\midrule
Schwarzschild-like & $\Phig^{\mathrm{bg}} = -GM/r$ & Pure gravitoelectric \\[4pt]
Reissner--Nordstr\"{o}m-like & $\Phig^{\mathrm{bg}} = -GM/r$,\; $\phi^{\mathrm{bg}} = Q/(4\pi\varepsilon_0 r)$ & Gravity vs.\ electrostatics \\[4pt]
Kerr-like (2D) & $\Phig^{\mathrm{bg}} = -GM/r$,\; $A_{g,\theta}^{\mathrm{bg}} = 2GJ/(c^2 r)$ & Frame dragging \\[4pt]
Kerr--Newman-like & All four potentials nonzero & Full background \\[4pt]
Uniform field & $\Phig^{\mathrm{bg}} = -g_0 y$ & Equivalence principle \\
\bottomrule
\end{tabular}
\end{center}

\subsection{Frame dragging from the Kerr-like background}

A mass $M$ with angular momentum $J$ generates in 2D:
\begin{equation}
  A_{g,\theta}^{\mathrm{bg}} = \frac{2GJ}{c^2 r}, \qquad
  B_{g,z}^{\mathrm{bg}} = -\frac{2GJ}{c^2 r^2}.
\end{equation}
A gyroscope at radius $r$ precesses at $\dd\phi/\dd t = B_{g,z}^{\mathrm{bg}}$.

\subsection{Three modes of operation}

\begin{center}
\begin{tabular}{@{}L{2.5cm}L{3cm}L{3cm}L{4cm}@{}}
\toprule
\textbf{Mode} & \textbf{Background} & \textbf{FDTD solver} & \textbf{Use case} \\
\midrule
1: Pure background & Fixed analytic & Not run & Test particles; zero FDTD cost \\
2: Background + perturbation & Fixed analytic & Perturbation fields & Few active interacting particles \\
3: Full dynamic & None fixed & All fields & Equal-mass many-body \\
\bottomrule
\end{tabular}
\end{center}

In Mode~2, the perturbation fields are sourced only by non-background particles.  The static background contributes no source to the perturbation wave equations.

\subsection{Validity conditions}

\begin{itemize}
  \item $|\Phig^{\mathrm{pert}}| \ll |\Phig^{\mathrm{bg}}|$: perturbations are small.
  \item Background approximately stationary (update adiabatically if needed).
  \item Back-reaction on background negligible.
  \item Background in linearized regime---exclude particles within a few Schwarzschild radii.
\end{itemize}

%====================================================================
\section{Numerical Implementation}
\label{sec:numerical}
%====================================================================

\subsection{Field grid}

The simulation is formulated primarily in terms of the scalar and vector \emph{potentials} rather than the derived field vectors $\vv{E}$, $\vv{B}$, $\vv{g}$, $\vv{B}_g$.  The potentials are the quantities evolved by the wave equations and stored on the grid.  The derived field vectors are gauge-invariant and are computed on demand --- at particle positions during force evaluation, and over the full grid only for visualization.  This eliminates four grid arrays ($g_x$, $g_y$, $E_x$, $E_y$, $B_{g,z}$, $B_z$ as persistent arrays) and the associated full-grid curl passes from the per-timestep inner loop.

The potentials are discretized on a staggered Yee grid~\cite{Yee1966,TafloveHagness2005}, which assigns different components to different spatial locations within each grid cell.  The staggering ensures that finite-difference gradient and curl operations on the potentials are exact --- no interpolation is needed.

\subsubsection*{The 2D Yee grid layout}

In 2D each cell is a square of side $\Delta x$.  The cell at index $(i,j)$ has its lower-left corner at position $(i\Delta x, j\Delta y)$.  Only the potentials and their sources are stored as persistent grid arrays:

\begin{center}
\begin{tabular}{@{}L{5cm}L{3cm}L{5cm}@{}}
\toprule
\textbf{Quantity} & \textbf{Location} & \textbf{Index notation} \\
\midrule
$\Phig$, $\phi$ (scalar potentials) & Cell corners & $(i,\; j)$ \\
$\rho_{\mathrm{matter}}$, $\rho_q$ (charge/mass density) & Cell corners & $(i,\; j)$ \\
$A_{g,x}$, $A_x$ ($x$-potentials) & Horizontal edge centers & $(i+\tfrac{1}{2},\; j)$ \\
$J_{\mathrm{matter},x}$, $J_{q,x}$ ($x$-currents) & Horizontal edge centers & $(i+\tfrac{1}{2},\; j)$ \\
$A_{g,y}$, $A_y$ ($y$-potentials) & Vertical edge centers & $(i,\; j+\tfrac{1}{2})$ \\
$J_{\mathrm{matter},y}$, $J_{q,y}$ ($y$-currents) & Vertical edge centers & $(i,\; j+\tfrac{1}{2})$ \\
\bottomrule
\end{tabular}
\end{center}

The derived quantities $g_x$, $g_y$, $E_x$, $E_y$ (which would live on edges) and $B_{g,z}$, $B_z$ (which would live at cell centers) are \emph{not} stored as persistent arrays.  They are computed locally at particle positions during force evaluation, and computed over the full grid only when requested for visualization.

The layout is shown schematically for one cell:

\begin{center}
\begin{tabular}{ccccc}
$(i,j+1)$ & $-$ & $A_x(i+\tfrac{1}{2},j+1)$ & $-$ & $(i+1,j+1)$ \\
$|$ & & & & $|$ \\
$A_y(i,j+\tfrac{1}{2})$ & & [cell center] & & $A_y(i+1,j+\tfrac{1}{2})$ \\
$|$ & & & & $|$ \\
$(i,j)$ & $-$ & $A_x(i+\tfrac{1}{2},j)$ & $-$ & $(i+1,j)$
\end{tabular}
\end{center}

\subsubsection*{Why this staggering: exact gradient and curl operations}

\paragraph{Gradient (for force and for visualization):}
The gravitoelectric field component $g_x = -\partial_x\Phig - \partial_t A_{g,x}$ is evaluated at a horizontal edge position $(i+\tfrac{1}{2}, j)$.  With $\Phig$ at corners, the spatial gradient $-\partial_x\Phig$ naturally lands at horizontal edges --- exactly where $A_{g,x}$ lives:
\begin{equation}
\label{eq:yee_grad}
  g_{x,i+1/2,j} = -\frac{(\Phig)_{i+1,j} - (\Phig)_{i,j}}{\Delta x}
                  - \frac{A_{g,x,i+1/2,j}^{n+1} - A_{g,x,i+1/2,j}^{n-1}}{2\Delta t},
\end{equation}
and similarly for $g_y$ at vertical edges, and for $\vv{E}$ from $\phi$ and $\vv{A}$.  This operation is performed on demand at particle positions (interpolating $\Phig$ from nearby corners and $\Ag$ from nearby edges) rather than sweeping the full grid.

\paragraph{Curl (for Boris rotation, spin precession, and visualization):}
The out-of-plane field $B_z = \partial_x A_y - \partial_y A_x$ at a cell center $(i+\tfrac{1}{2}, j+\tfrac{1}{2})$ is:
\begin{equation}
\label{eq:yee_curl}
  B_{z,i+1/2,j+1/2} = \frac{A_{y,i+1,j+1/2} - A_{y,i,j+1/2}}{\Delta x}
                     - \frac{A_{x,i+1/2,j+1} - A_{x,i+1/2,j}}{\Delta y},
\end{equation}
and identically for $B_{g,z}$ from $\Ag$.  This is evaluated at the few cell centers surrounding the particle when needed for the Boris rotation or spin precession, not over the full grid each step.  For visualization, a full-grid sweep applies Eq.~(\ref{eq:yee_curl}) at every cell center.

\paragraph{Source deposition staggering:}
Source terms must be deposited at the same locations as the potentials they source:
\begin{itemize}
  \item $\rho_{\mathrm{matter}}$, $\rho_q$: at cell corners $(i,j)$.  The CIC kernel $W_2$ is evaluated at corner positions relative to the particle.
  \item $J_{\mathrm{matter},x}$, $J_{q,x}$: at horizontal edge centers $(i+\tfrac{1}{2},j)$.  The CIC kernel is evaluated at horizontal edge positions.  The sub-cell fraction is $\alpha = (x_p - (i+\tfrac{1}{2})\Delta x)/\Delta x$, measured from the edge center rather than the corner.
  \item $J_{\mathrm{matter},y}$, $J_{q,y}$: at vertical edge centers $(i,j+\tfrac{1}{2})$.
\end{itemize}

\paragraph{Spin dipole curl deposition:}
The magnetization curl $\nabla\times[\mu\,W_3]$ deposits $J_{q,x}$ at horizontal edges and $J_{q,y}$ at vertical edges, consistent with the current staggering.  The $x$-component is $\mu\,W_3(x-x_p)\cdot(\dd W_3/\dd y)(y-y_p)$, depositing to horizontal edges via the smooth $W_3$ weight in $x$ and the two-tent $\dd W_3/\dd y$ weight in $y$, centered on horizontal edge positions $(i+\tfrac{1}{2}, j)$.

\subsection{Time stepping: the leapfrog scheme}
\label{sec:leapfrog}

\subsubsection*{The field leapfrog and its three-point stencil}

The d'Alembert equation $\partial_t^2\Phi = c^2\nabla^2\Phi + \mathrm{source}$ is discretized in time using the three-point leapfrog (Störmer-Verlet) scheme.  The superscripts $n-1$, $n$, $n+1$ denote values at three consecutive integer timesteps separated by $\dd t$:
\begin{equation}
\label{eq:leapfrog_field}
  \Phi^{n+1} = 2\Phi^n - \Phi^{n-1} + (c\,\dd t)^2\!\left[\nabla^2\Phi^n + \mathrm{source}^n\right].
\end{equation}
The term $2\Phi^n - \Phi^{n-1}$ is a linear extrapolation of $\Phi$ forward by one timestep using the last two known values.  It is the free-propagation prediction --- what $\Phi^{n+1}$ would be if the right-hand side were zero (no source, no spatial coupling).  The $(c\,\dd t)^2[\cdots]$ term is the correction from the wave equation.  This is not overrelaxation in the iterative solver sense.  Rather, it is the standard second-order centered finite difference for $\partial_t^2\Phi$:
\begin{equation}
  \frac{\Phi^{n+1} - 2\Phi^n + \Phi^{n-1}}{\dd t^2} \approx \frac{\partial^2\Phi}{\partial t^2}\bigg|_n + \order{\dd t^2},
\end{equation}
rearranged to solve for $\Phi^{n+1}$.  The scheme is second-order accurate in time and requires storing only two time levels ($\Phi^n$ and $\Phi^{n-1}$) to advance to $\Phi^{n+1}$.

The CFL stability condition requires:
\begin{equation}
  \dd t < \frac{\dd x}{c\sqrt{d}}, \qquad d = \text{spatial dimension} \quad (d=2 \text{ or } 3).
\end{equation}

\subsubsection*{The $\alpha$-family: interpolating between Euler and leapfrog}

The leapfrog can be placed in a one-parameter family that continuously interpolates between a naive Euler scheme and the full leapfrog:
\begin{equation}
  \Phi^{n+1} = (1+\alpha)\Phi^n - \alpha\Phi^{n-1} + \dd t^2\,f^n,
\end{equation}
where $f^n = c^2\nabla^2\Phi^n + \mathrm{source}^n$.  The parameter $\alpha$ controls the history weighting:
\begin{align}
  \alpha = 0: &\quad \Phi^{n+1} = \Phi^n + \dd t^2 f^n \quad\text{(forward Euler applied to }\partial_t^2; \text{ first-order, unstable for wave eq.)} \\
  \alpha = 1: &\quad \Phi^{n+1} = 2\Phi^n - \Phi^{n-1} + \dd t^2 f^n \quad\text{(standard leapfrog; second-order, conditionally stable)}
\end{align}
Values $0 < \alpha < 1$ give intermediate schemes that are first-order accurate but more stable than pure Euler.  Values $\alpha > 1$ are analogous to overrelaxation and are \emph{unstable} for the wave equation --- the leapfrog at $\alpha=1$ is the maximally stable point in the explicit family.

The analogous family for the implicit $\theta$-method (Newmark-$\beta$) is:
\begin{equation}
  \Phi^{n+1} - 2\Phi^n + \Phi^{n-1} = \dd t^2\!\left[\theta f^{n+1} + (1-2\theta)f^n + \theta f^{n-1}\right],
\end{equation}
with $\theta=0$ recovering the explicit leapfrog and $\theta=\tfrac{1}{4}$ giving an unconditionally stable implicit scheme.  The implicit schemes require solving a linear system at each step and are not used here.

\subsubsection*{The leapfrog stagger: what lives when}

The field and particle updates are interleaved on a staggered time grid.  Understanding what quantity lives at what time is essential for correctly evaluating forces, sources, and proper time.  The complete stagger is:

\begin{center}
\begin{tabular}{@{}L{6cm}L{3cm}L{4cm}@{}}
\toprule
\textbf{Quantity} & \textbf{Time level} & \textbf{Notes} \\
\midrule
Scalar potentials $\Phig$, $\phi$ & Integer $n\,\dd t$ & Two levels stored: $n$ and $n-1$ \\
Vector potentials $\Ag$, $\vv{A}$ & Integer $n\,\dd t$ & Two levels stored: $n$ and $n-1$ \\
Particle positions $\vv{x}_i$ & Integer $n\,\dd t$ & \\
Source terms $\rho$, $\vv{J}$ & Integer $n\,\dd t$ & Deposited from $\vv{x}_i^n$, $\vv{v}_i^{n-1/2}$ \\
Forces $\vv{F}_i$ & Integer $n\,\dd t$ & Computed from potentials at particle position \\
Particle momenta $\vv{p}_i$, velocities $\vv{v}_i$ & Half-integer $(n\pm\tfrac{1}{2})\dd t$ & Boris kick uses $\vv{F}^n$ \\
Proper time increment $\dd\tau_i$ & Integer $n\,\dd t$ & Uses $\bar{v}_i = \tfrac{1}{2}(v^{n-1/2}+v^{n+1/2})$; reuses $(\Ag\cdot\vv{v})_i$ from force step \\
\bottomrule
\end{tabular}
\end{center}

The derived fields $\vv{g}$, $\vv{B}_g$, $\vv{E}$, $\vv{B}$ are not stored as persistent arrays.  They are computed on demand from the potentials at particle positions during force evaluation (Step~4 below) and over the full grid only for visualization.

The per-timestep sequence is:
\begin{enumerate}
  \item Deposit sources from $\vv{x}_i^n$, $\vv{v}_i^{n-1/2}$ onto grid $\to$ $\rho^n$, $\vv{J}^n$.
  \item Step potentials: $(\Phig^{n-1}, \Phig^n, \rho^n) \to \Phig^{n+1}$, and similarly for $\phi$, $\Ag$, $\vv{A}$.
  \item Apply gauge cleaning to $\Phig^{n+1}$, $\Ag^{n+1}$ (Section~\ref{sec:gauge}).
  \item Compute force $\vv{F}_i^n$ at each particle from potentials at surrounding grid points (see Section on force from potentials below).
  \item Boris kick: $(\vv{p}_i^{n-1/2}, \vv{F}_i^n) \to \vv{p}_i^{n+1/2}$.
  \item Drift: $\vv{x}_i^{n+1} = \vv{x}_i^n + \vv{v}_i^{n+1/2}\,\dd t$.
  \item Update spin using $\vv{a}_i^n = (\vv{v}_i^{n+1/2} - \vv{v}_i^{n-1/2})/\dd t$.
  \item Update proper time using $\bar{v}_i = \tfrac{1}{2}(\vv{v}_i^{n-1/2} + \vv{v}_i^{n+1/2})$, $(\Phig^{\mathrm{total}})_i$, and $(\Ag^{\mathrm{total}}\cdot\bar{\vv{v}}_i)$ --- the last quantity already computed during the force evaluation step.
\end{enumerate}
The force $\vv{F}_i^n$ and position $\vv{x}_i^n$ are both at integer time, and the momentum $\vv{p}_i^{n+1/2}$ is at the half-integer time between them --- this stagger is what makes the scheme symplectic.

\subsubsection*{Force from potentials}
\label{sec:force_potential}

The Lorentz force in EM, written directly in terms of the scalar potential $\phi$ and vector potential $\vv{A}$, using the vector identity $\vv{v}\times(\nabla\times\vv{A}) = \nabla(\vv{v}\cdot\vv{A}) - (\vv{v}\cdot\nabla)\vv{A}$:
\begin{equation}
\label{eq:em_force_potential}
  \vv{F}_{\mathrm{EM}} = q\!\left(-\nabla\phi - \partial_t\vv{A} + \nabla(\vv{v}\cdot\vv{A}) - (\vv{v}\cdot\nabla)\vv{A}\right).
\end{equation}
The GEM analogue (with factor of 4 on the vector potential terms from the spin-2 structure):
\begin{equation}
\label{eq:gem_force_potential}
  \vv{F}_{\mathrm{GEM}} = m\!\left(-\nabla\Phig - \partial_t\Ag + 4\nabla(\vv{v}\cdot\Ag) - 4(\vv{v}\cdot\nabla)\Ag\right).
\end{equation}

\paragraph{$\vv{v}$ is a constant vector, not a field.}
In these expressions $\vv{v} = \vv{v}_i$ is the velocity of a specific particle evaluated at its position --- a single constant vector, not a spatially varying field.  The full vector identity for two arbitrary vector fields $\vv{u}(\vv{x})$ and $\vv{A}(\vv{x})$ contains two additional terms:
\begin{equation}
  \vv{u}\times(\nabla\times\vv{A}) = \nabla(\vv{u}\cdot\vv{A}) - (\vv{u}\cdot\nabla)\vv{A} - (\vv{A}\cdot\nabla)\vv{u} - \vv{A}\times(\nabla\times\vv{u}).
\end{equation}
When $\vv{u} = \vv{v}_i$ is a constant (i.e.\ $\nabla\vv{v}_i = 0$), the last two terms vanish and the identity reduces to the two-term form used here.  The spatial gradient $\nabla$ in $\nabla(\vv{v}\cdot\vv{A})$ and the directional derivative $(\vv{v}\cdot\nabla)$ act only on $\vv{A}(\vv{x})$, with $\vv{v}$ held fixed.

\paragraph{The directional derivative $(\vv{v}\cdot\nabla)\vv{A}$.}
The notation $(\vv{v}\cdot\nabla)\vv{A}$ is the directional derivative of the vector field $\vv{A}$ in the direction of $\vv{v}$.  The operator $\vv{v}\cdot\nabla = v_x\partial_x + v_y\partial_y + v_z\partial_z$ is a scalar operator applied component-wise:
\begin{equation}
  (\vv{v}\cdot\nabla)\vv{A} =
  \begin{pmatrix}
    v_x\,\partial_x A_x + v_y\,\partial_y A_x + v_z\,\partial_z A_x \\
    v_x\,\partial_x A_y + v_y\,\partial_y A_y + v_z\,\partial_z A_y \\
    v_x\,\partial_x A_z + v_y\,\partial_y A_z + v_z\,\partial_z A_z
  \end{pmatrix}.
\end{equation}
In 2D the $z$-row is absent and $v_z = 0$.  On the Yee grid the $x$-component of $(\vv{v}\cdot\nabla)\vv{A}$ at the particle position is:
\begin{equation}
  \left[(\vv{v}\cdot\nabla)A_x\right]_p = v_x\,\frac{A_{x,i+1,j} - A_{x,i,j}}{\Delta x}
                                         + v_y\,\frac{A_{x,i+1/2,j+1} - A_{x,i+1/2,j}}{\Delta y},
\end{equation}
where $v_x$ and $v_y$ are the scalar particle velocity components --- constants that multiply the spatial finite differences of $A_x$.  The $y$-component follows identically with $A_y$.

The combination $\nabla(\vv{v}\cdot\vv{A}) - (\vv{v}\cdot\nabla)\vv{A}$ encodes the magnetic cross-product force.  To verify, the $x$-component is:
\begin{equation}
  \left[\nabla(\vv{v}\cdot\vv{A}) - (\vv{v}\cdot\nabla)\vv{A}\right]_x
  = v_y(\partial_x A_y - \partial_y A_x) = v_y B_z,
\end{equation}
which is the $x$-component of $\vv{v}\times\vv{B}$, as expected.  The cross product is therefore correctly encoded without ever explicitly forming $\vv{B}$ as a grid array.

These expressions involve no derived field vectors --- only spatial and temporal differences of the potentials evaluated at and near the particle position.  They are gauge-invariant: under $\phi\to\phi-\partial_t\Lambda$, $\vv{A}\to\vv{A}+\nabla\Lambda$ the extra terms cancel identically, so gauge drift in the potentials does not contaminate the computed force.

Each term is evaluated by interpolating potentials from the surrounding grid points to the particle position using the CIC kernel, then forming the appropriate differences:
\begin{itemize}
  \item $-\nabla\phi$: finite difference of $\phi$ at surrounding corners, CIC-interpolated to the particle.
  \item $-\partial_t\vv{A}$: centered time difference $(A_\alpha^{n+1} - A_\alpha^{n-1})/(2\Delta t)$ at edge positions, CIC-interpolated to the particle.
  \item $\nabla(\vv{v}\cdot\vv{A}) = v_x\nabla A_x + v_y\nabla A_y$: spatial differences of the edge-located $A_x$ and $A_y$, with $v_x$, $v_y$ as constant multipliers.
  \item $(\vv{v}\cdot\nabla)\vv{A}$: directional derivative of $\vv{A}$ along $\vv{v}$ as defined above --- spatial differences of edge values weighted by velocity components.
\end{itemize}
The full force stencil touches the same $4$ corners and $8$ edge points as the CIC interpolation, requiring no larger neighborhood.

\subsubsection*{The Boris pusher and on-the-fly curl}

The Boris algorithm~\cite{Boris1970,BirdsallLangdon1991} requires the magnetic rotation step, which needs $B_{g,z}$ and $B_z$ at the particle position.  These are computed on the fly from the surrounding edge potentials using Eq.~(\ref{eq:yee_curl}), evaluated at the few cell centers near the particle and then CIC-interpolated to the particle position:
\begin{equation}
  B_{z}(\vv{x}_p) \approx \sum_{j,k} B_{z,j+1/2,k+1/2} \cdot W_2(\vv{x}_{j+1/2,k+1/2} - \vv{x}_p),
\end{equation}
where each $B_{z,j+1/2,k+1/2}$ is evaluated from the four surrounding edge values via Eq.~(\ref{eq:yee_curl}).  This touches at most 4 cell centers near the particle --- a negligible cost.

The effective combined fields for Boris are then:
\begin{equation}
  \vv{B}_{\mathrm{eff}} = 4\,(\nabla\times\Ag)(\vv{x}_p) + \frac{q}{m}(\nabla\times\vv{A})(\vv{x}_p),
\end{equation}
where both curls are evaluated on-the-fly at the particle position from surrounding edge potentials via Eq.~(\ref{eq:yee_curl}).  In 2D the curl fields reduce to scalars with only a $z$-component, so:
\begin{equation}
  B_{\mathrm{eff},z} = 4\,B_{g,z}(\vv{x}_p) + \frac{q}{m}B_z(\vv{x}_p) \qquad\text{(2D specialization).}
\end{equation}

The Boris algorithm splits the total force into two parts.  The \textbf{electric-like force} $\vv{F}_{\mathrm{eff,E}}$ collects all terms that do not involve the magnetic cross product --- the gradient and time-derivative potential terms:
\begin{equation}
\label{eq:feff}
  \vv{F}_{\mathrm{eff,E}} = m(-\nabla\Phig - \partial_t\Ag)
                          + q(-\nabla\phi  - \partial_t\vv{A})
                          + \vv{F}_{\mathrm{spin}} + \vv{F}_{\mathrm{rel}},
\end{equation}
where $\vv{F}_{\mathrm{spin}} = \nabla(\vv{S}\cdot\nabla\times\Ag) + \nabla(\boldsymbol{\mu}\cdot\nabla\times\vv{A})$ and $\vv{F}_{\mathrm{rel}}$ are the spin gradient and relativistic correction forces (Section~\ref{sec:relcorr}).  The velocity-dependent cross-product terms $4\nabla(\vv{v}\cdot\Ag) - 4(\vv{v}\cdot\nabla)\Ag$ and $\nabla(\vv{v}\cdot\vv{A}) - (\vv{v}\cdot\nabla)\vv{A}$ encode the magnetic rotation and are handled implicitly by the Boris rotation step using $B_{\mathrm{eff},z}$, not included in $\vv{F}_{\mathrm{eff,E}}$.

The Boris update applies $\vv{F}_{\mathrm{eff,E}}$ as two half-kicks bracketing the exact magnetic rotation:
\begin{align}
  \vv{p}^- &= \vv{p}^{n-1/2} + \tfrac{1}{2}\vv{F}_{\mathrm{eff,E}}\,\dd t, \\
  \vv{t}   &= \frac{q\,\vv{B}_{\mathrm{eff}}}{2\gamma^- m}\,\dd t, \qquad
  \vv{s} = \frac{2\vv{t}}{1+|\vv{t}|^2}, \\
  \vv{p}'  &= \vv{p}^- + \vv{p}^-\times\vv{t}, \qquad
  \vv{p}^+ = \vv{p}^- + \vv{p}'\times\vv{s}, \\
  \vv{p}^{n+1/2} &= \vv{p}^+ + \tfrac{1}{2}\vv{F}_{\mathrm{eff,E}}\,\dd t.
\end{align}
In 2D $\vv{B}_{\mathrm{eff}} = B_{\mathrm{eff},z}\hat{z}$, so $\vv{t} = t_z\hat{z}$ and $\vv{s} = s_z\hat{z}$ are purely out-of-plane, and the cross products $\vv{p}^-\times\vv{t}$ and $\vv{p}'\times\vv{s}$ reduce to in-plane rotations by the 2D scalar $t_z = q\,B_{\mathrm{eff},z}\,\dd t/(2\gamma^- m)$.  The Boris rotation step is exactly symplectic.  The $\order{v^2/c^2}$ gravitational corrections enter as additional contributions to $\vv{F}_{\mathrm{rel}}$ in Eq.~(\ref{eq:feff}).

\subsection{Symplectic structure and relativistic limitations}

The leapfrog scheme preserves the phase-space area element $\dd\vv{x}\wedge\dd\vv{p}$ exactly for linear (shear) force laws, preventing secular energy drift.  At relativistic speeds the electric kick becomes a curved shear whose Jacobian depends on $\vv{p}$, and the Boris approximation is symplectic only to $\order{\dd t^2}$.  The symplectic error grows as $(\gamma v/c)^2(\dd t/T)^2$ and can be mitigated by subcycling fast particles.  Photon-like particles ($v\approx c$) should use null geodesic ray tracing in the $c_{\mathrm{local}}$ field rather than the Boris pusher.

\subsection{Numerical dispersion and anisotropy}
\label{sec:dispersion}

\subsubsection*{The exact numerical dispersion relation}

The continuous wave equation has the dispersion relation $\omega^2 = c^2k^2$, meaning all wavelengths propagate at exactly $c$.  The discrete leapfrog FDTD scheme on a uniform grid with spacing $\Delta x$ and timestep $\Delta t$ has a different, exact dispersion relation.  Substituting a plane wave ansatz $\Phi^n_{j,k} = \hat\Phi\,e^{i(\omega n\Delta t - k_x j\Delta x - k_y k\Delta x)}$ into the leapfrog update gives:
\begin{equation}
\label{eq:dispersion_exact}
  \sin^2\!\left(\frac{\omega\Delta t}{2}\right)
  = c_{\mathrm{CFL}}^2\!\left[\sin^2\!\left(\frac{k_x\Delta x}{2}\right) + \sin^2\!\left(\frac{k_y\Delta x}{2}\right)\right],
\end{equation}
where $c_{\mathrm{CFL}} = c\Delta t/\Delta x$ is the CFL number.  This is not an approximation --- it is the exact dispersion relation of the discrete scheme.

\subsubsection*{Dimensionless variables for the expansion}

To expand Eq.~(\ref{eq:dispersion_exact}) cleanly it is convenient to work entirely in dimensionless phase variables:
\begin{equation}
  \xi = \frac{\omega\Delta t}{2}, \qquad
  \beta_x = \frac{k_x\Delta x}{2}, \qquad
  \beta_y = \frac{k_y\Delta x}{2},
\end{equation}
so that Eq.~(\ref{eq:dispersion_exact}) becomes simply:
\begin{equation}
  \sin^2\xi = c_{\mathrm{CFL}}^2(\sin^2\beta_x + \sin^2\beta_y).
\end{equation}
All quantities are dimensionless.  The continuous limit is $\xi = c_{\mathrm{CFL}}\sqrt{\beta_x^2+\beta_y^2}$, i.e.\ $\omega = ck$.

\subsubsection*{Taylor expansion}

Using $\sin^2 u = u^2 - u^4/3 + \order{u^6}$:
\begin{align}
  \text{Left side:}&\quad \sin^2\xi = \xi^2 - \frac{\xi^4}{3} + \ldots \label{eq:lhs_expand}\\
  \text{Right side:}&\quad c_{\mathrm{CFL}}^2\!\left[(\beta_x^2+\beta_y^2) - \frac{\beta_x^4+\beta_y^4}{3} + \ldots\right] \label{eq:rhs_expand}
\end{align}
Equating and collecting terms, with $\beta^2 \equiv \beta_x^2+\beta_y^2$:
\begin{equation}
  \xi^2\!\left(1 - \frac{\xi^2}{3}\right) = c_{\mathrm{CFL}}^2\beta^2\!\left(1 - \frac{\beta_x^4+\beta_y^4}{3\beta^2}\right).
\end{equation}
At leading order $\xi^2 = c_{\mathrm{CFL}}^2\beta^2$, confirming the continuous dispersion relation.  Substituting this into the $\order{\beta^4}$ correction terms and solving for $\xi^2$:
\begin{equation}
\label{eq:dispersion_expanded}
  \xi^2 \approx c_{\mathrm{CFL}}^2\beta^2\!\left[1 + \frac{\beta^2}{3}\!\left(c_{\mathrm{CFL}}^2 - \frac{\beta_x^4+\beta_y^4}{\beta^4}\right)\right].
\end{equation}
All terms are dimensionless throughout.  Converting back to physical variables via $\xi = \omega\Delta t/2$ and $\beta_\alpha = k_\alpha\Delta x/2$:
\begin{equation}
\label{eq:dispersion_physical}
  \omega^2 \approx c^2k^2\!\left[1 + \frac{k^2\Delta x^2}{12}\!\left(c_{\mathrm{CFL}}^2 - \frac{k_x^4+k_y^4}{k^4}\right)\right].
\end{equation}
Every term now has units $[\mathrm{rad}^2/\mathrm{s}^2]$, with the correction being a dimensionless factor.

\subsubsection*{Structure of the dispersion error}

The correction factor in Eq.~(\ref{eq:dispersion_physical}) has two contributions that partially oppose each other:
\begin{itemize}
  \item $+c_{\mathrm{CFL}}^2\,k^2\Delta x^2/12$: from the \textbf{temporal} discretization.  Isotropic, positive --- makes $\omega$ too large, so the numerical phase velocity is slightly above $c$.
  \item $-(k_x^4+k_y^4)/k^4\cdot k^2\Delta x^2/12$: from the \textbf{spatial} discretization.  Anisotropic, negative --- makes $\omega$ too small, with a magnitude that depends on the propagation angle $\theta$.
\end{itemize}
The numerical phase velocity along a general direction $\theta$ (where $k_x = k\cos\theta$, $k_y = k\sin\theta$) is:
\begin{equation}
  \frac{v_\phi(\theta)}{c} \approx 1 + \frac{k^2\Delta x^2}{24}\!\left(c_{\mathrm{CFL}}^2 - \frac{\cos^4\theta + \sin^4\theta}{1}\right) + \order{(k\Delta x)^4}.
\end{equation}
Using $\cos^4\theta + \sin^4\theta = 1 - \tfrac{1}{2}\sin^2(2\theta)$, this becomes:
\begin{equation}
\label{eq:phase_vel}
  \frac{v_\phi(\theta)}{c} \approx 1 + \frac{k^2\Delta x^2}{24}\!\left(c_{\mathrm{CFL}}^2 - 1 + \frac{\sin^2(2\theta)}{2}\right).
\end{equation}

\subsubsection*{Anisotropy}

The anisotropy --- the variation of phase velocity with propagation angle --- comes from the $\sin^2(2\theta)/2$ term in Eq.~(\ref{eq:phase_vel}).  This term is independent of $c_{\mathrm{CFL}}$: the anisotropy is a fixed property of the square grid and the 5-point Laplacian stencil, not of the timestep.  The maximum anisotropy occurs between the axis directions ($\theta = 0°$, $90°$) and the diagonals ($\theta = 45°$):
\begin{align}
  \frac{v_\phi(0°)}{c}  &\approx 1 + \frac{k^2\Delta x^2}{24}(c_{\mathrm{CFL}}^2 - 1), \\
  \frac{v_\phi(45°)}{c} &\approx 1 + \frac{k^2\Delta x^2}{24}\!\left(c_{\mathrm{CFL}}^2 - \frac{1}{2}\right).
\end{align}
The difference between the two:
\begin{equation}
  \frac{v_\phi(45°) - v_\phi(0°)}{c} \approx \frac{k^2\Delta x^2}{48}
\end{equation}
is independent of $c_{\mathrm{CFL}}$, confirming that the CFL number does not affect the anisotropy.  A circular wavefront generated by a point source will appear as a slightly rounded square due to this anisotropy, with the diagonal directions propagating faster than the axis directions by a fractional amount $(k\Delta x)^2/48$.

\subsubsection*{Overall dispersion magnitude and the role of $c_{\mathrm{CFL}}$}

While the anisotropy is fixed, the overall magnitude of the dispersion error does depend on $c_{\mathrm{CFL}}$ through the $c_{\mathrm{CFL}}^2 - 1$ term.  Along the grid axes:
\begin{itemize}
  \item At $c_{\mathrm{CFL}} = 1$ (1D stability limit): $v_\phi(0°) = c$ exactly --- zero dispersion along the axis in 1D.
  \item At $c_{\mathrm{CFL}} = 1/\sqrt{2}$ (2D stability limit): $v_\phi(0°)/c \approx 1 - k^2\Delta x^2/48$ (slightly slow), $v_\phi(45°)/c \approx 1$ (exact at the diagonal).
  \item At $c_{\mathrm{CFL}} \ll 1$: $v_\phi(0°)/c \approx 1 - k^2\Delta x^2/24$ (maximally slow along axes).
\end{itemize}
So running well below the stability limit does worsen the overall dispersion (larger negative error along the axes) while the anisotropy remains the same.  Running at the stability limit minimizes the frequency-dependent part of the dispersion error by maximally cancelling the temporal and spatial contributions, though at no $c_{\mathrm{CFL}}$ value is the 2D dispersion simultaneously zero in all directions.

\subsubsection*{Practical implications for the sandbox}

\begin{itemize}
  \item \textbf{Run at the CFL limit.} $c_{\mathrm{CFL}} = 1/\sqrt{2}$ in 2D minimizes the overall dispersion error while remaining stable.  The conventional advice to run ``safely below'' the CFL limit worsens wave propagation accuracy without improving stability.
  \item \textbf{Ensure adequate spatial resolution.}  The dispersion error scales as $(k\Delta x)^2$.  A wavelength spanning $N$ grid cells has $k\Delta x = 2\pi/N$, giving a dispersion error of order $\pi^2/(6N^2)$.  For $N=10$ cells per wavelength this is $\approx 1.6\%$; for $N=20$ it is $\approx 0.4\%$.  For the sandbox, 10--20 cells per wavelength is adequate for qualitative demonstrations.
  \item \textbf{Anisotropy is a fixed grid artifact.}  A circularly-emitted wavefront will be slightly non-circular due to the $(k\Delta x)^2/48$ anisotropy, appearing as a rounded square with the diagonals slightly ahead of the axes.  This can be reduced by replacing the standard 5-point Laplacian with the 9-point isotropic stencil:
  \begin{equation}
    \nabla^2\Phi \approx \frac{1}{\Delta x^2}\!\left[\frac{1}{6}\sum_{\mathrm{diag}}\Phi + \frac{2}{3}\sum_{\mathrm{axis}}\Phi - \frac{10}{3}\Phi_{i,j}\right],
  \end{equation}
  which replaces $k_x^4+k_y^4$ with $k^4$ in the spatial error term, eliminating the angular dependence and reducing the anisotropy coefficient by a factor of $\sim\!4$ at the cost of involving 8 neighbours instead of 4.
  \item \textbf{Higher-order stencils.}  Using a 4th-order accurate spatial stencil replaces the $\order{(k\Delta x)^2}$ error with $\order{(k\Delta x)^4}$, dramatically reducing dispersion for well-resolved wavelengths at the cost of a wider stencil (5 points per dimension rather than 3).
\end{itemize}

\subsection{Source deposition: approximating the delta function}
\label{sec:deposition}

The source-deposition strategy follows standard particle-in-cell (PIC) practice~\cite{HockneyEastwood1988,BirdsallLangdon1991}.  The physical sources involve $\delta(\vv{x}-\vv{x}_i)$ and for spin sources $\nabla\times[\boldsymbol{\mu}\,\delta(\vv{x}-\vv{x}_i)]$, neither of which can exist on a discrete grid.  They must be replaced by smooth kernel functions $W(\vv{x}-\vv{x}_i)$ that approximate the delta function on the grid scale.  Three competing requirements must be balanced:

\begin{enumerate}
  \item \textbf{Sub-cell position accuracy:} the kernel must encode the particle's position within a cell, not just which cell it occupies.
  \item \textbf{Smoothness for curl deposition:} taking $\nabla\times[W\,\boldsymbol{\mu}]$ requires $W$ to have well-defined spatial derivatives.
  \item \textbf{Compactness:} a wide kernel smears the source over many cells, reducing spatial resolution.
\end{enumerate}

\subsubsection*{The B-spline hierarchy}

There is a well-established family of kernels that interpolates between a box function (maximally compact, discontinuous) and a Gaussian (infinitely smooth, infinite support).  These are the B-spline kernels $W_n$ of order $n$, each being the convolution of its predecessor with a box function:

\begin{center}
\begin{tabular}{@{}lllll@{}}
\toprule
\textbf{Order} & \textbf{Common name} & \textbf{Smoothness} & \textbf{Support (cells)} & \textbf{Curl kernel} \\
\midrule
$n=1$ & Nearest-grid-point (NGP) & Discontinuous & 1 & Undefined \\
$n=2$ & Tent / CIC               & $C^0$         & 2 & Box (discontinuous) \\
$n=3$ & Quadratic B-spline        & $C^1$         & 3 & Tent ($C^0$) \\
$n=4$ & Cubic B-spline            & $C^2$         & 4 & Quadratic ($C^1$) \\
$n\to\infty$ & Gaussian          & $C^\infty$    & $\infty$ & Hermite--Gaussian \\
\bottomrule
\end{tabular}
\end{center}

\subsubsection*{Monopole and current deposition: the CIC kernel and bilinear interpolation}

For $\rho_{\mathrm{matter}}$, $\vv{J}_{\mathrm{matter}}$, $\rho_q$, $\vv{J}_q$ the cloud-in-cell (CIC) tent kernel $W_2$ is the standard choice.  In 1D it is:
\begin{equation}
  W_2(x) = \begin{cases}
    1 - |x|/\Delta x & |x| < \Delta x \\
    0                & \text{otherwise,}
  \end{cases}
\end{equation}
and in 2D it is the tensor product:
\begin{equation}
  W_2(x,y) = W_2(x)\,W_2(y)
            = \left(1-\frac{|x|}{\Delta x}\right)\!\left(1-\frac{|y|}{\Delta y}\right),
  \quad |x|<\Delta x,\;|y|<\Delta y.
\end{equation}
Deposition of charge $q$ at sub-cell position $(\alpha,\beta)$ with $\alpha=(x_p-x_i)/\Delta x$, $\beta=(y_p-y_j)/\Delta y$ distributes to the four surrounding grid points:
\begin{equation}
\label{eq:cic_deposit}
\begin{aligned}
  \rho_{i,j}     &\mathrel{+}= q\,(1-\alpha)(1-\beta)/\Delta x^2, \\
  \rho_{i+1,j}   &\mathrel{+}= q\,\alpha(1-\beta)/\Delta x^2, \\
  \rho_{i,j+1}   &\mathrel{+}= q\,(1-\alpha)\beta/\Delta x^2, \\
  \rho_{i+1,j+1} &\mathrel{+}= q\,\alpha\beta/\Delta x^2.
\end{aligned}
\end{equation}
The weights $(1-\alpha)(1-\beta)$, $\alpha(1-\beta)$, $(1-\alpha)\beta$, $\alpha\beta$ are exactly the bilinear interpolation weights.  Field interpolation back to the particle position uses the same four weights:
\begin{equation}
\label{eq:cic_interp}
  E_p = (1-\alpha)(1-\beta)\,E_{i,j}
       + \alpha(1-\beta)\,E_{i+1,j}
       + (1-\alpha)\beta\,E_{i,j+1}
       + \alpha\beta\,E_{i+1,j+1}.
\end{equation}
The deposition weights in Eq.~(\ref{eq:cic_deposit}) and the interpolation weights in Eq.~(\ref{eq:cic_interp}) are identical.  This is the adjoint condition, discussed in detail below.

\subsubsection*{Dipole deposition: the quadratic B-spline and its derivative as two tents}

The quadratic B-spline $W_3$ is obtained by convolving $W_2$ with a box function of width $\Delta x$.  In 1D:
\begin{equation}
  W_3(x) = \begin{cases}
    \tfrac{3}{4} - (x/\Delta x)^2                             & |x| < \tfrac{1}{2}\Delta x \\
    \tfrac{1}{2}\!\left(\tfrac{3}{2} - |x|/\Delta x\right)^2 & \tfrac{1}{2}\Delta x \leq |x| < \tfrac{3}{2}\Delta x \\
    0                                                          & \text{otherwise.}
  \end{cases}
\end{equation}
$W_3$ is $C^1$, spans 3 cells, and has an inflection point at $|x|=\tfrac{1}{2}\Delta x$---concave down in the central region and concave up in the two outer regions.  Its derivative is:
\begin{equation}
\label{eq:dW3}
  \frac{\dd W_3}{\dd x}(x) = \begin{cases}
    +\!\left(\tfrac{3}{2} - |x|/\Delta x\right)/\Delta x & -\tfrac{3}{2}\Delta x \leq x < -\tfrac{1}{2}\Delta x \\
    -2x/\Delta x^2                                        & |x| < \tfrac{1}{2}\Delta x \\
    -\!\left(\tfrac{3}{2} - |x|/\Delta x\right)/\Delta x & \phantom{-}\tfrac{1}{2}\Delta x \leq x < \phantom{-}\tfrac{3}{2}\Delta x \\
    0                                                     & \text{otherwise.}
  \end{cases}
\end{equation}

The key geometric insight is that $\dd W_3/\dd x$ consists of \textbf{two tent functions placed side by side}---one positive and one negative---straddling the particle position at $x=0$:
\begin{itemize}
  \item A \textbf{positive tent} centered at $x=-\tfrac{1}{2}\Delta x$, rising from 0 at $x=-\tfrac{3}{2}\Delta x$, peaking at $+1/\Delta x$ at $x=-\tfrac{1}{2}\Delta x$, and falling back to 0 at $x=0$.
  \item A \textbf{negative tent} centered at $x=+\tfrac{1}{2}\Delta x$, falling from 0 at $x=0$, reaching $-1/\Delta x$ at $x=+\tfrac{1}{2}\Delta x$, and returning to 0 at $x=+\tfrac{3}{2}\Delta x$.
\end{itemize}
The two tents meet at zero at the particle position.  This antisymmetric structure is a direct consequence of $W_3$ being symmetric: the derivative of a symmetric function is antisymmetric and must pass through zero at the center, positive on the rising (left) side and negative on the falling (right) side.

When deposited onto the grid, the three cells receive:
\begin{itemize}
  \item \textbf{Cell $i-1$} (left): only the positive tail of the positive tent.
  \item \textbf{Cell $i$} (containing particle): the right half of the positive tent (positive) plus the left half of the negative tent (negative), partially cancelling with net value depending on sub-cell position $\alpha$.
  \item \textbf{Cell $i+1$} (right): only the negative tail of the negative tent.
\end{itemize}
The total weight across all three cells sums to zero, since $W_3(\pm\tfrac{3}{2}\Delta x)=0$.

For comparison, the derivative of the CIC kernel $W_2$ is a discontinuous box function:
\begin{equation}
  \frac{\dd W_2}{\dd x}(x) = \begin{cases}
    -\mathrm{sgn}(x)/\Delta x & |x| < \Delta x \\
    0                         & \text{otherwise,}
  \end{cases}
\end{equation}
which produces Gibbs-like ringing when used as a grid source.  The $W_3$ two-tent derivative avoids this discontinuity.

For reference, the Gaussian $W_G = (2\pi\sigma^2)^{-1/2}\exp(-x^2/2\sigma^2)$ gives a Hermite--Gaussian derivative:
\begin{equation}
  \frac{\dd W_G}{\dd x} = -\frac{x}{\sigma^2}\,W_G(x),
\end{equation}
which is $C^\infty$ but requires truncation at $\pm 3\sigma$ ($\approx 6$ cells).

\subsubsection*{The 2D curl: sign flip in the direction of differentiation}

In 2D, the spin dipole moment is a scalar $\mu$ (the out-of-plane component), and the magnetization current is the 2D curl of $\mu\,W_3(x-x_p)\,W_3(y-y_p)$:
\begin{align}
  J_{q,x}(\vv{x}) &= \mu\,W_3(x-x_p)\cdot\frac{\dd W_3}{\dd y}(y-y_p), \label{eq:curl2d_x}\\
  J_{q,y}(\vv{x}) &= -\mu\,\frac{\dd W_3}{\dd x}(x-x_p)\cdot W_3(y-y_p). \label{eq:curl2d_y}
\end{align}
Each component is a tensor product of one smooth $W_3$ factor (all positive, $C^1$) and one two-tent $\dd W_3/\dd x_\alpha$ factor ($C^0$, sign flip).  The sign flip occurs only in the direction of differentiation:
\begin{itemize}
  \item $J_{q,x}$: smooth in $x$, two-tent in $y$.  Positive above the particle ($y>y_p$), negative below.
  \item $J_{q,y}$: two-tent in $x$ with overall minus sign, smooth in $y$.  Negative to the left of the particle, positive to the right.
\end{itemize}
Together these form a counter-clockwise current loop in the $xy$-plane---the correct physical current distribution for a magnetic dipole $\mu>0$ pointing out of the plane.  Each component spans $3\times 3 = 9$ cells in 2D.

\subsubsection*{Generalization to 3D}

All formulas above are written in 2D for clarity.  The 3D generalization adds the third tensor product factor straightforwardly.

\paragraph{CIC (monopole/current):} add factor $W_2(z-z_p)$ with $\gamma=(z_p-z_k)/\Delta z$.  The 8 surrounding cells receive weights that are all products of factors from $\{(1-\alpha),\alpha\}\times\{(1-\beta),\beta\}\times\{(1-\gamma),\gamma\}$.

\paragraph{Quadratic B-spline (dipole):} the 3D kernel is $W_3(x-x_p)\,W_3(y-y_p)\,W_3(z-z_p)$, spanning $3^3=27$ cells.

\paragraph{3D curl for spin sources:} with $\boldsymbol{\mu}$ now a 3-vector, the three components of $\nabla\times[\boldsymbol{\mu}\,W_3]$ are:
\begin{align}
  [\nabla\times(\boldsymbol{\mu}\,W_3)]_x &= \mu_z\,W_3(x)\,\tfrac{\dd W_3}{\dd y}(y)\,W_3(z)
                                           - \mu_y\,W_3(x)\,W_3(y)\,\tfrac{\dd W_3}{\dd z}(z), \\
  [\nabla\times(\boldsymbol{\mu}\,W_3)]_y &= \mu_x\,W_3(x)\,W_3(y)\,\tfrac{\dd W_3}{\dd z}(z)
                                           - \mu_z\,\tfrac{\dd W_3}{\dd x}(x)\,W_3(y)\,W_3(z), \\
  [\nabla\times(\boldsymbol{\mu}\,W_3)]_z &= \mu_y\,\tfrac{\dd W_3}{\dd x}(x)\,W_3(y)\,W_3(z)
                                           - \mu_x\,W_3(x)\,\tfrac{\dd W_3}{\dd y}(y)\,W_3(z),
\end{align}
where arguments $(x-x_p)$ etc.\ are suppressed.  Each term has exactly one two-tent $\dd W_3/\dd x_\alpha$ factor in the direction of differentiation and smooth $W_3$ factors in the other two directions, spanning 27 cells and summing to zero over the grid.  In 2D only $\mu_z$ is nonzero and only the $x$ and $y$ components of the curl exist (Eqs.~\ref{eq:curl2d_x}--\ref{eq:curl2d_y}).

\subsubsection*{The adjoint condition and self-force cancellation}

The deposition operator $D$ maps a particle quantity at position $\vv{x}_p$ to the grid:
\begin{equation}
  \rho(\vv{x}_j) = [D\,q](\vv{x}_j) = q\,W(\vv{x}_j - \vv{x}_p),
\end{equation}
and the interpolation operator $I$ maps a grid field back to the particle:
\begin{equation}
  E_p = [I\,E](\vv{x}_p) = \sum_j E(\vv{x}_j)\,\tilde{W}(\vv{x}_j - \vv{x}_p).
\end{equation}
The \textbf{adjoint condition}~\cite{HockneyEastwood1988} requires $\tilde{W} = W$---the same kernel is used for both operations.

To see why this cancels the static self-force, write the force of the particle's own field on itself:
\begin{equation}
  \vv{F}_{\mathrm{self}} = q\sum_j \vv{E}_{\mathrm{self}}(\vv{x}_j)\,W(\vv{x}_j-\vv{x}_p).
\end{equation}
For a stationary particle, $\vv{E}_{\mathrm{self}}$ is produced by the symmetric charge distribution $\rho(\vv{x}_j) = q\,W(\vv{x}_j-\vv{x}_p)$.  On a uniform grid with a symmetric kernel $W$, the field $\vv{E}_{\mathrm{self}}$ is antisymmetric about the particle position---it points away from the particle equally in all directions.  Interpolating this antisymmetric field back with the same symmetric kernel $W$ gives exactly zero:
\begin{equation}
  \vv{F}_{\mathrm{self}} = q\sum_j \vv{E}_{\mathrm{self}}(\vv{x}_j)\,W(\vv{x}_j-\vv{x}_p) = 0 \qquad \text{(stationary particle).}
\end{equation}
If different kernels are used for deposition and interpolation, this symmetry is broken and the self-force is nonzero even for a stationary particle, producing unphysical heating.

For the CIC/bilinear case, the adjoint condition is immediately visible: the deposition weights $(1-\alpha)(1-\beta)$, $\alpha(1-\beta)$, $(1-\alpha)\beta$, $\alpha\beta$ in Eq.~(\ref{eq:cic_deposit}) are identical to the interpolation weights in Eq.~(\ref{eq:cic_interp}).  Bilinear interpolation \emph{is} CIC deposition---they are the same operation, so the adjoint condition is automatically satisfied.

For the gradient fields needed by the spin forces, the adjoint condition has a critical implication that must be stated precisely.  The dipole source was deposited using the kernel $\dd W_3/\dd x_\alpha \cdot W_3$ over 9 cells in 2D (27 in 3D).  The adjoint of that deposition operator is the interpolation that uses those \emph{exact same weights} to map grid field values back to the particle.  Concretely, the $x$-gradient of $B_{g,z}$ at the particle is:
\begin{align}
  (\partial_x B_{g,z})_p &= \sum_{j,k} B_{g,z}(\vv{x}_{j,k})\cdot
    \frac{\dd W_3}{\dd x}(x_j - x_p)\cdot W_3(y_k - y_p), \label{eq:grad_adjoint_x}\\
  (\partial_y B_{g,z})_p &= \sum_{j,k} B_{g,z}(\vv{x}_{j,k})\cdot
    W_3(x_j - x_p)\cdot\frac{\dd W_3}{\dd y}(y_k - y_p), \label{eq:grad_adjoint_y}
\end{align}
where the sums run over the same $3\times 3 = 9$ cells used during dipole deposition, and $\dd W_3/\dd x$ is the two-tent function of Eq.~(\ref{eq:dW3}).  The identical formulas apply to $\nabla B_z$ for the magnetic gradient force, and generalize to $3\times 3\times 3 = 27$ cells in 3D.

It is important to note what is \emph{not} correct here.  The box-function derivative $-\mathrm{sgn}(x)/\Delta x$ is the derivative of the CIC kernel $W_2$, not of $W_3$.  It would be the correct adjoint gradient interpolation only if the dipole had been deposited with $W_2$---but $W_2$ was deliberately avoided for dipole deposition because its discontinuous derivative causes Gibbs ringing.  Using $\mathrm{sgn}(x)/\Delta x$ for gradient interpolation while using $W_3$ for dipole deposition breaks the adjoint pairing and introduces a residual dipole self-force: a spinning stationary particle would experience a spurious gradient force from its own field.  The two-tent weights of Eqs.~(\ref{eq:grad_adjoint_x})--(\ref{eq:grad_adjoint_y}) are the unique correct choice.

Similarly, computing the gradient by finite-differencing the bilinearly interpolated field (another common approach) is not adjoint to $W_3$ dipole deposition and also introduces a self-force.  The rule is simple: \textbf{use the same kernel weights for interpolation that were used for deposition, applied to the same set of cells}.

\subsubsection*{Self-interaction, radiation reaction, and their limitations}

The self-force cancellation described above is exact only for a \emph{stationary} particle.  For an \emph{accelerating} particle the retarded self-field is asymmetric---the field the particle emitted in the past arrives back at the particle's current position from a different direction---and the cancellation is incomplete.  The residual is the \textbf{radiation reaction force}, which in classical electrodynamics is the Abraham--Lorentz force:
\begin{equation}
  \vv{F}_{\mathrm{rad}} = \frac{q^2}{6\pi\varepsilon_0 c^3}\dddot{\vv{x}},
\end{equation}
involving the third time derivative of position (the jerk).  The gravitational analogue for GEM is structurally identical with the appropriate coupling constants.

In a grid-based simulation this radiation reaction force is captured only approximately.  The particle deposits its field onto the grid via the CIC kernel---effectively giving the particle a finite size of order $\Delta x$---and the retarded self-field propagates back to the particle's position via the wave equation.  The resulting self-force approximates but does not exactly reproduce the Abraham--Lorentz force, for two reasons:

\begin{itemize}
  \item The effective particle size $\sim\Delta x$ regularizes the $1/r$ self-field divergence at a scale chosen for numerical convenience rather than physical accuracy, making the radiation reaction force grid-spacing-dependent.
  \item The exact Abraham--Lorentz force involves the $r\to 0$ limit of the retarded self-field, which the grid cannot resolve below $\Delta x$.
\end{itemize}

The practical consequence is that a radiating particle in the simulation will correctly emit energy into the field (the outgoing wave is accurately computed), but the back-reaction force on the particle that causes it to slow down may be slightly inconsistent with the emitted power, leading to small violations of energy conservation at the particle level over long integration times.

For qualitative demonstrations---watching waves propagate outward from an oscillating mass, observing the radiation pattern, verifying the $1/\sqrt{r}$ field fall-off in 2D---the simulation is fully adequate.  For quantitative long-time tracking of radiation-induced orbital decay, a hybrid approach is recommended: add an explicit phenomenological radiation damping term computed analytically from the particle's current acceleration, rather than relying entirely on the grid self-interaction.  This maintains energy conservation to much higher accuracy without the numerical stiffness of the exact Abraham--Lorentz--Dirac equation.

\subsubsection*{Practical recommendation}

\begin{center}
\begin{tabular}{@{}L{5.5cm}L{3.5cm}L{3.5cm}@{}}
\toprule
\textbf{Source / Force} & \textbf{Kernel} & \textbf{Cells (2D / 3D)} \\
\midrule
$\rho$, $\vv{J}$ monopole/current deposition & CIC $W_2$ & 4 / 8 \\
Potential interpolation at particle ($\phi$, $\Phig$) & Bilinear $W_2$ at corners & 4 / 8 \\
Potential interpolation at particle ($\vv{A}$, $\Ag$) & Bilinear $W_2$ at edges & 4 / 8 \\
On-the-fly $B_{g,z}$, $B_z$ at particle & Yee curl + bilinear $W_2$ & 4 cell centers \\
Spin dipole $\nabla\times[\boldsymbol{\mu}\,\delta]$ deposition & $W_3$ curl & 9 / 27 \\
Gradient fields $\nabla B_{g,z}$, $\nabla B_z$ at particle & Two-tent $\dd W_3/\dd x_\alpha$ & 9 / 27 \\
\bottomrule
\end{tabular}
\end{center}

\subsection{Gauge condition enforcement}
\label{sec:gauge}

\subsubsection*{The Lorenz gauge and why it drifts}

The wave equations for the GEM and EM potentials were derived in the Lorenz gauge, which for GEM requires:
\begin{equation}
\label{eq:lorenz_gem}
  \mathcal{G}_{\mathrm{GEM}} \equiv \frac{1}{c^2}\frac{\partial\Phig}{\partial t} + \nabla\cdot\Ag = 0,
\end{equation}
and for EM:
\begin{equation}
\label{eq:lorenz_em}
  \mathcal{G}_{\mathrm{EM}} \equiv \frac{1}{c^2}\frac{\partial\phi}{\partial t} + \nabla\cdot\vv{A} = 0.
\end{equation}
In the continuous theory, if these conditions hold at $t=0$ and the sources satisfy the continuity equations $\partial_t\rho + \nabla\cdot\vv{J} = 0$, then the gauge conditions are preserved exactly by the wave equation evolution --- the gauge violation satisfies the homogeneous wave equation $\Box\mathcal{G} = 0$ and simply propagates outward and exits through the absorbing boundary.  There is no exponential amplification mechanism.

In the discrete FDTD scheme the gauge condition drifts for two reasons.  First, the finite-difference operators do not satisfy the same algebraic identities as the continuous operators.  Second, the discretely deposited sources satisfy the continuity equation only approximately.  The drift is sourced by truncation errors of order $(k\Delta x)^2$ per step --- small, bounded, and not self-amplifying.  It accumulates at a slow linear rate, not exponentially.

\subsubsection*{Why gauge drift does not affect particle dynamics}

A crucial observation is that the force expressions Eqs.~(\ref{eq:em_force_potential})--(\ref{eq:gem_force_potential}) are \emph{gauge-invariant}.  Under a gauge transformation $\phi\to\phi-\partial_t\Lambda$, $\vv{A}\to\vv{A}+\nabla\Lambda$, the extra terms introduced into the force cancel exactly --- the physical force is independent of the gauge function $\Lambda$.  Therefore gauge drift in the stored potentials has \emph{no effect} on the computed forces or particle trajectories.

Gauge drift matters only when the potentials themselves are used directly: for visualization of $\phi$, $\Phig$, $\Ag$, $\vv{A}$; for the gauge condition monitoring; and for the EM effective GEM source $\rho_{\mathrm{EM}} = u_{\mathrm{EM}}/c^2$, which uses $\vv{E}$ and $\vv{B}$ derived from the potentials.  For a purely qualitative pedagogical simulation gauge correction can be omitted entirely without affecting the particle dynamics.

\subsubsection*{The discrete continuity equation: the primary defense}

The most important gauge-preserving measure is to ensure that the deposited sources satisfy the \emph{discrete} continuity equation exactly.  On the Yee grid with $\rho$ at corners and $\vv{J}$ at edges, the discrete continuity equation is:
\begin{equation}
  \frac{\rho^{n+1}_{i,j} - \rho^{n-1}_{i,j}}{2\Delta t}
  + \frac{J_{x,i+1/2,j}^n - J_{x,i-1/2,j}^n}{\Delta x}
  + \frac{J_{y,i,j+1/2}^n - J_{y,i,j-1/2}^n}{\Delta y} = 0.
\end{equation}
Using a consistent CIC kernel for both $\rho$ and $\vv{J}$ deposition satisfies this condition to the accuracy of the kernel, bounding gauge drift at the truncation error level automatically.

\subsubsection*{Parabolic cleaning: optional refinement for long integrations}

For very long simulations where the slow linear gauge drift eventually becomes visible in potential visualizations, a parabolic cleaning step can be applied after each field update:
\begin{equation}
\label{eq:parabolic_clean}
  \Phig \;\leftarrow\; \Phig - R\,c^2\,\Delta t\,(\nabla\cdot\Ag), \qquad
  \phi  \;\leftarrow\; \phi  - R\,c^2\,\Delta t\,(\nabla\cdot\vv{A}),
\end{equation}
where $R$ is a small positive dimensionless constant.  This drives the scalar potentials toward the Lorenz gauge condition, damping the gauge violation on a timescale $\sim 1/(Rc)$.  A practical choice is $R \sim c_{\mathrm{CFL}}^2/4$, giving a damping timescale of a few domain-crossing times.

The cleaning is a single extra multiplication per cell and is applied only to the perturbation fields.  Since gauge drift does not affect particle dynamics, this correction is purely cosmetic --- it keeps the potential visualizations free of slowly drifting background artifacts.  It is not needed for physical correctness of the simulation.

\subsubsection*{Hyperbolic cleaning: stronger alternative}

If parabolic cleaning is insufficient (e.g.\ for very long integrations or highly dynamic configurations), hyperbolic divergence cleaning can be used.  An auxiliary field $\psi_g$ (for GEM; similarly $\psi$ for EM) is introduced:
\begin{align}
  \frac{\partial\Ag}{\partial t} &\;\leftarrow\; \frac{\partial\Ag}{\partial t} - \nabla\psi_g, \\
  \frac{\partial\psi_g}{\partial t} &= -c_h^2\nabla\cdot\Ag - \kappa\psi_g,
\end{align}
making the gauge violation satisfy a damped wave equation that propagates violations outward at speed $c_h = c$ and damps them at rate $\kappa = c/L$.  This is more robust than parabolic cleaning but requires an additional scalar field per wave equation.

\subsubsection*{Monitoring}

The discrete gauge residual:
\begin{equation}
  \mathcal{G}_{i,j}^n = \frac{(\Phig^{\mathrm{pert}})_{i,j}^{n+1} - (\Phig^{\mathrm{pert}})_{i,j}^{n-1}}{2c^2\Delta t}
  + \frac{A_{g,x,i+1/2,j}^n - A_{g,x,i-1/2,j}^n}{\Delta x}
  + \frac{A_{g,y,i,j+1/2}^n - A_{g,y,i,j-1/2}^n}{\Delta y}
\end{equation}
should remain at the level of the truncation error $\sim(k\Delta x)^2$.  Rapid growth indicates inconsistent source deposition.  Since gauge drift does not affect dynamics, monitoring is primarily useful for diagnosing deposition bugs rather than correcting physical errors.

\subsection{Absorbing boundary conditions}
\label{sec:abc}

Without special treatment at the grid edges the leapfrog FDTD scheme imposes reflecting boundary conditions by default --- the field values outside the grid are implicitly zero, which acts as a perfect reflector for any outgoing wave.  Reflected waves re-enter the simulation domain and interfere with the physics being demonstrated.  For a GR sandbox this is particularly problematic: a gravitational wave emitted by an oscillating mass would bounce off the grid boundary, return to the source region, and create standing wave patterns that have no physical meaning and would confuse any student watching.

\subsubsection*{The damping layer}

The recommended approach for the sandbox is a \textbf{damping layer}: a border region of $N_d$ cells on each side of the grid where the field is multiplied by a spatially varying attenuation factor at each timestep.  This is a simplification of the perfectly matched layer (PML)~\cite{Berenger1994,TafloveHagness2005}; the quadratic profile used here is a degenerate matched-impedance case suitable for the qualitative pedagogical purpose of the sandbox.  Outgoing waves entering the layer are attenuated gradually and the small residual that reaches the hard boundary is attenuated again on the return trip, making double-pass reflections negligible.

The damping factor at position $x$ within the layer (measured from the inner edge of the layer, $0 \leq x \leq L = N_d\Delta x$) is:
\begin{equation}
  f(x) = \exp\!\left(-\sigma(x)\,\dd t\right) \approx 1 - \sigma(x)\,\dd t,
\end{equation}
where $\sigma(x)$ is the conductivity profile.  A smooth polynomial ramp is recommended:
\begin{equation}
\label{eq:damp_profile}
  \sigma(x) = \sigma_{\max}\!\left(\frac{x}{L}\right)^2,
\end{equation}
which rises from zero at the interior edge to $\sigma_{\max}$ at the outer edge.  The quadratic profile ensures the attenuation starts gently, minimizing the impedance mismatch at the interior-damping interface that would otherwise cause partial reflection.

The round-trip reflection coefficient for a wave traversing the layer twice is:
\begin{equation}
  R \approx \exp\!\left(-\frac{2}{c}\int_0^L \sigma(x)\,\dd x\right)
           = \exp\!\left(-\frac{2\sigma_{\max}L}{3c}\right).
\end{equation}
Setting $2\sigma_{\max}L/3c = 7$ gives $R \approx e^{-7} \approx 10^{-3}$, which is adequate for all qualitative demonstrations.  For a layer of $N_d = 16$ cells this requires:
\begin{equation}
  \sigma_{\max} = \frac{21c}{2L} = \frac{21c}{2 N_d \Delta x}.
\end{equation}

\subsubsection*{Implementation}

The damping factor array is precomputed once at initialization.  For a grid of size $N\times N$ with damping layer of thickness $N_d$ cells on each side, define:
\begin{equation}
  d_{i,j} = \max\!\left(d_x(i),\; d_y(j)\right),
\end{equation}
where:
\begin{align}
  d_x(i) &= \begin{cases}
    \sigma_{\max}\!\left(\dfrac{N_d - i}{N_d}\right)^2 \dd t & i < N_d \quad\text{(left layer)} \\[6pt]
    \sigma_{\max}\!\left(\dfrac{i - (N-N_d)}{N_d}\right)^2 \dd t & i > N - N_d \quad\text{(right layer)} \\
    0 & \text{otherwise,}
  \end{cases}
\end{align}
and similarly for $d_y(j)$.  The per-timestep application after each field update is then:

\begin{equation}
  \Phi^{n+1}_{i,j} \;\leftarrow\; \Phi^{n+1}_{i,j} \cdot (1 - d_{i,j})
\end{equation}

applied to all six perturbation field arrays $(\Phig^{\mathrm{pert}},\; A_{g,x}^{\mathrm{pert}},\; A_{g,y}^{\mathrm{pert}},\; \phi^{\mathrm{pert}},\; A_x^{\mathrm{pert}},\; A_y^{\mathrm{pert}})$.  The same damping array $d_{i,j}$ applies to all six because they all satisfy the same wave equation with the same propagation speed.

\subsubsection*{What is and is not damped}

Only the perturbation fields are damped.  The background fields $\Phig^{\mathrm{bg}}$, $A_{g,\alpha}^{\mathrm{bg}}$, $\phi^{\mathrm{bg}}$, $A_\alpha^{\mathrm{bg}}$ are evaluated analytically at particle positions and are never stored on the grid, so they are automatically exempt from the damping operation.  This is a further advantage of the background + perturbation split: the static $1/r$ gravitational field, which is not a propagating wave and should not be attenuated, requires no special treatment at the boundary.

Particles should be confined to the interior region and prevented from entering the damping layer.  A simple reflecting or absorbing boundary at the inner edge of the damping layer (at $i = N_d$ and $i = N - N_d$, and similarly in $y$) is sufficient.  A particle that reaches the inner boundary can be reflected elastically, absorbed and removed from the simulation, or wrapped with a warning — the choice depends on the scenario being demonstrated.

\subsubsection*{Practical parameters}

\begin{center}
\begin{tabular}{@{}ll@{}}
\toprule
\textbf{Parameter} & \textbf{Recommended value} \\
\midrule
Layer thickness $N_d$ & 16 cells \\
Profile exponent & 2 (quadratic) \\
$\sigma_{\max}$ & $21c / (2 N_d \Delta x)$ \\
Grid overhead & $(N + 2N_d)^2 / N^2 \approx 1.27$ for $N=256$, $N_d=16$ \\
Residual reflection $R$ & $\approx 10^{-3}$ \\
\bottomrule
\end{tabular}
\end{center}

For scenarios where lower reflection is needed --- for example, measuring radiated power or verifying the $1/\sqrt{r}$ amplitude fall-off --- the layer thickness can be increased to $N_d = 32$ cells and $\sigma_{\max}$ adjusted accordingly, reducing $R$ to $\approx 10^{-6}$.

\subsubsection*{Realized reflection in practice (added in v33)}

The $R\approx10^{-3}$ figure above is the \emph{plane-wave normal-incidence} limit of the continuum derivation.  Two effects degrade the realized in-interior residual for a 2D circular wavefront launched from a centered Gaussian initial condition:

\begin{itemize}
  \item \textbf{Linearization error}.  The continuum derivation uses $f(x)=\exp(-\sigma(x)\,\dd t)$, while the implementation block above prescribes the linearized form $(1-d_{i,j})$.  For the default parameters ($N_d=16$, $\sigma_{\max}=21c/(2N_d\Delta x)$, $c_{\mathrm{CFL}}=1/\sqrt{2}$), $\sigma_{\max}\,\dd t \approx 0.46$, where the two forms differ by $\sim15\%$.  The linearization \emph{over-attenuates} per step, but distorts the discrete dispersion of the absorbing region and introduces a small impedance-mismatch reflection at the inner edge of the layer.
  \item \textbf{Oblique-incidence corner trapping}.  The axis-aligned profile $d_{i,j}=\max(d_x,d_y)$ damps perpendicular incidence strongly but is weak for waves heading into the grid \emph{corners} at $\approx45^\circ$ from both walls.  Such waves arrive at the corner having spent comparable depth along both axes; the $\max$ operator does not accumulate their depths.  A perfectly matched layer (PML) addresses this by introducing complex coordinate stretching that absorbs in each direction independently.
\end{itemize}

Empirically (measured in Stage~2 of the implementation plan, §\ref{sec:impl}): with $N_d=16$ and a $\sigma=4\Delta x$ Gaussian pulse at the center of a $256\times256$ grid, after 800 timesteps the field at the cell just inside the wall is suppressed by $\sim 3000\times$ versus the no-damping case, while at the domain center it is suppressed by $\sim 20\times$ --- the difference being the corner-trapped energy slowly refocusing through the interior over multiple round-trips.  The damping layer is therefore fully adequate for qualitative pedagogical demonstrations where outgoing waves visibly exit the domain, but does \emph{not} achieve $10^{-3}$ residual at the source for a 2D circular wavefront without moving to PML.

\subsection{Units and non-dimensionalization}
\label{sec:units}

The physical values of $G$ and $c$ make gravitational effects completely invisible at any computationally tractable scale.  For example, the Schwarzschild radius of a solar mass is $\sim 3\,\mathrm{km}$, which would require astronomical grid extents to simulate at physical scales.  The simulation therefore uses a rescaled set of units in which all quantities are of order unity.

Choose three base scales: a length scale $\ell_0$ (the grid cell size in physical units, or equivalently the domain size), a time scale $t_0 = \ell_0/c_{\mathrm{eff}}$, and a mass scale $m_0$.  All simulation quantities are then expressed in terms of these:

\begin{center}
\begin{tabular}{@{}L{4cm}L{3cm}L{3cm}L{3cm}@{}}
\toprule
\textbf{Quantity} & \textbf{Physical} & \textbf{Simulation} & \textbf{Relation} \\
\midrule
Length & $x$ & $\tilde{x} = x/\ell_0$ & $\tilde{x} \in [0, N]$ \\
Time & $t$ & $\tilde{t} = t/t_0$ & $\dd\tilde{t} = 1$ per step \\
Speed of light & $c$ & $\tilde{c} = c_{\mathrm{eff}}$ & chosen so $\tilde{c}\,\dd\tilde{t} = c_{\mathrm{CFL}}\,\dd\tilde{x}$ \\
Mass & $m$ & $\tilde{m} = m/m_0$ & typically $\tilde{m} \sim 1$ \\
Charge & $q$ & $\tilde{q} = q/q_0$ & $q_0$ chosen freely \\
Momentum & $p$ & $\tilde{p} = p/(m_0 c_{\mathrm{eff}})$ & \\
Grav.\ potential & $\Phig$ & $\tilde{\Phi}_g = \Phig/c_{\mathrm{eff}}^2$ & dimensionless \\
Grav.\ constant & $G$ & $\tilde{G} = G_{\mathrm{eff}}$ & chosen for visibility \\
\bottomrule
\end{tabular}
\end{center}

The key free parameter is $G_{\mathrm{eff}}$ — the effective gravitational constant in simulation units.  It should be chosen so that a typical particle-particle gravitational interaction produces a visible deflection over a few timesteps.  A useful guideline: for two particles of mass $\tilde{m}$ separated by distance $\tilde{r}$, the gravitational acceleration is $\tilde{G}\tilde{m}/\tilde{r}^2$ in simulation units.  Setting this to be a fraction of $\tilde{c}^2/\tilde{r}$ (the speed squared over the separation, a natural relativistic scale) gives:

\begin{equation}
  G_{\mathrm{eff}} \sim \frac{c_{\mathrm{eff}}^2\,\tilde{r}}{\tilde{m}} \cdot \epsilon,
\end{equation}

where $\epsilon \ll 1$ keeps the system in the weak-field linearized regime.  Typical values in practice are $G_{\mathrm{eff}} \sim 10^{-3}$ to $10^{-1}$ depending on the desired strength of gravitational effects.

The charge-to-mass ratio $q/m$ for test particles can be set independently of any physical value, allowing electromagnetic and gravitational effects to be tuned relative to each other.  Setting $q/m \gg \sqrt{G_{\mathrm{eff}}}$ makes EM effects dominate; setting $q/m \ll \sqrt{G_{\mathrm{eff}}}$ makes gravitational effects dominate.  The coupling of EM field energy to the GEM equations through $\rho_{\mathrm{EM}} = u_{\mathrm{EM}}/c^2$ is automatically consistent with whatever $G_{\mathrm{eff}}$ and $c_{\mathrm{eff}}$ are chosen, since it follows directly from the field equations.

\subsection{Field initialization}
\label{sec:init}

At $t=0$ the perturbation fields must be initialized to the static solution corresponding to the initial particle configuration.  Starting with zero fields and nonzero particle sources creates a spurious transient wavefront that propagates outward at $c$ --- a pulse with no physical meaning that takes several domain crossing times to damp out in the absorbing boundary layer.

\subsubsection*{Convergence iteration}

The cleanest initialization approach is to iterate the FDTD field update with all \textbf{positions} held fixed at their initial values and all \textbf{velocities} held fixed at their initial values, until the fields converge.  The particles are not moved during this phase --- only the field arrays are updated each step.  With fixed positions and velocities:
\begin{itemize}
  \item Charge/mass density sources $\rho_{\mathrm{matter}}$, $\rho_q$ are time-independent (positions fixed).
  \item Current sources $\vv{J}_{\mathrm{matter}} = \gamma m\vv{v}$ and $\vv{J}_q = q\vv{v}$ are also time-independent (velocities fixed).
  \item All six source terms are constant, so the wave equations iterate toward their static solutions --- the Poisson equation solutions for $\Phig$ and $\phi$, and the vector potential solutions for $\Ag$ and $\vv{A}$ sourced by the constant currents.
\end{itemize}
This is important: velocities must \emph{not} be set to zero during initialization unless the particles are actually stationary.  The vector potentials $\Ag$ and $\vv{A}$ are sourced by $\vv{J}_{\mathrm{matter}}$ and $\vv{J}_q$ respectively, which depend on velocity.  A particle with nonzero initial velocity generates gravitomagnetic and magnetic fields that must be present in the initial field configuration.  Setting velocities to zero during initialization would converge to the wrong field --- one with no vector potentials --- which is only correct for truly stationary particles.

For the special case of all particles stationary ($\vv{v}_i = 0$), the currents vanish and $\Ag \to 0$, $\vv{A} \to 0$ as a consequence rather than as an assumption.

The transient oscillations during convergence are absorbed by the damping layer.  The convergence timescale is a few domain-crossing times $\sim N\sqrt{2}$ timesteps (at $c_{\mathrm{CFL}} = 1/\sqrt{2}$), roughly 1000--3000 steps for $N=256$.  This can be done invisibly before the simulation begins displaying to the user.

\subsubsection*{Leapfrog initialization}

The leapfrog requires field values at two consecutive time levels $\Phi^{-1}$ and $\Phi^0$ to start.  With zero initial time derivative ($\partial_t\Phi = 0$ for a static initial configuration), the correct initialization is:
\begin{equation}
  \Phi^{-1} = \Phi^0 = 0,
\end{equation}
which makes the free-propagation extrapolation $2\Phi^0 - \Phi^{-1} = 0$ correct.  As the convergence iteration proceeds both levels are updated normally.

For the particle momenta, what is needed is $\vv{p}_i^{-1/2}$ --- the momentum at a half timestep before $t=0$.  The standard approach is a single half-step using the force from the converged initial fields:
\begin{equation}
  \vv{p}_i^{-1/2} = \vv{p}_i^0 - \vv{F}_i^0\,\frac{\dd t}{2},
\end{equation}
where $\vv{F}_i^0$ is the total force on particle $i$ from the converged initial fields at $t=0$, and $\vv{p}_i^0$ is the specified initial momentum.  For a stationary initial configuration $\vv{p}_i^0 = 0$ and the half-step momentum is $-\vv{F}_i^0\,\dd t/2$ --- a small quantity representing the velocity the particle would have acquired over a half timestep from the static initial field.  This initialization is first-order accurate; the error affects only the very first full leapfrog step, after which the scheme is self-consistent and second-order accurate.

The particle position $\vv{x}_i^0$ lives at integer time and is specified directly by the initial conditions --- no back-extrapolation of position is needed.

\subsubsection*{Interactive particle placement}

When a particle is added mid-simulation the sudden deposition of its charge and mass creates a field disturbance that propagates outward at $c$.  This is physically correct: in classical field theory a source that suddenly appears radiates a spherical pulse, with the new static field inside the pulse and the pre-existing field outside it.  For the pedagogical sandbox this is a feature --- a student who places a mass and watches a gravitational wavefront propagate outward is seeing real physics.

If a smooth turn-on is preferred, ramp the particle's effective mass and charge from zero to their final values over $\sim N\sqrt{2}$ timesteps, suppressing the pulse at the cost of physical honesty.

The convergence iteration used for startup cannot be applied to mid-simulation particle addition without erasing all existing propagating waves in the domain --- it would reset the entire field to the new static configuration.  It is only appropriate before any dynamics have begun.

Particle removal is symmetric: a particle that disappears leaves a void in its Coulomb/Newton field that also propagates outward as a pulse.  If a particle exits through the domain boundary the most physically honest treatment is to remove it and accept the outgoing pulse, which will be absorbed by the damping layer.  The fundamental ambiguity in both creation and removal is that a particle appearing or disappearing has no prior history to specify what the surrounding field configuration should have been, making any treatment somewhat arbitrary; the pulse-based approach at least has a clear physical interpretation.

\subsection{Field-particle energy accounting}
\label{sec:energy}

The linearized GEM theory is not strictly energy-conserving in the full GR sense: the gravitational field carries energy in the nonlinear theory through the Landau-Lifshitz pseudotensor, which does not appear in the linearized equations.  What the linearized theory does conserve, in the same structural sense that EM conserves energy, is the combined energy of the matter, GEM field, and EM field, with fluxes through the boundary accounting for radiation losses.

This section describes the conserved quantities and their use as \emph{consistency diagnostics} rather than strict conservation laws.  They verify that the numerical scheme is correctly implementing the linearized equations, not that the linearized equations are a complete description of reality.

\subsubsection*{Energy}

The total energy in the simulation domain is:
\begin{equation}
  E_{\mathrm{total}} = E_{\mathrm{particles}} + E_{\mathrm{GEM}} + E_{\mathrm{EM}},
\end{equation}
where:
\begin{align}
  E_{\mathrm{particles}} &= \sum_i \gamma_i m_i c^2, \\
  E_{\mathrm{GEM}}       &= \frac{1}{8\pi G}\int\!\!\left(|\vv{g}|^2 + c^2|\vv{B}_g|^2\right)\dd^3x, \\
  E_{\mathrm{EM}}        &= \frac{\varepsilon_0}{2}\int\!\!\left(|\vv{E}|^2 + c^2|\vv{B}|^2\right)\dd^3x.
\end{align}
In 2D the volume integrals become area integrals ($\dd^3x \to \dd^2x$), $\vv{B}_g \to B_{g,z}\hat{z}$ and $\vv{B} \to B_z\hat{z}$.

The rate of change of $E_{\mathrm{total}}$ should equal the energy flux leaving through the damping layer boundary:
\begin{equation}
  \frac{\dd E_{\mathrm{total}}}{\dd t} = -\oint_{\partial\Omega}(\vv{S}_{\mathrm{GEM}} + \vv{S}_{\mathrm{EM}})\cdot\hat{n}\,\dd A,
\end{equation}
where $\vv{S}_{\mathrm{GEM}} = (\vv{g}\times\vv{B}_g)/(4\pi G)$ and $\vv{S}_{\mathrm{EM}} = (\vv{E}\times\vv{B})/\mu_0$ are the GEM and EM Poynting vectors (in 2D, $\dd A \to \dd\ell$).  In the presence of an absorbing boundary layer the boundary flux is not easily computed directly, so in practice one monitors $E_{\mathrm{total}}$ and checks that it changes smoothly and monotonically (decreasing as radiation exits through the boundary), with no sudden jumps that would indicate a numerical instability.

\subsubsection*{Momentum}

The total canonical momentum is:
\begin{equation}
  \vv{P}_{\mathrm{total}} = \sum_i\vv{p}_i + \frac{1}{c^2}\int\!\!\left(\vv{S}_{\mathrm{GEM}} + \vv{S}_{\mathrm{EM}}\right)\dd^3x.
\end{equation}
In the absence of external forces and with symmetric boundary conditions, this should be conserved.  In practice the damping layer introduces a small asymmetry if waves exit preferentially in one direction, but for a symmetric simulation $\vv{P}_{\mathrm{total}}$ should remain near zero.  A systematic drift in one direction indicates a bug in the force calculation or source deposition.

\subsubsection*{Angular momentum}

The total angular momentum is:
\begin{equation}
  \vv{L}_{\mathrm{total}} = \sum_i\!\left(\vv{r}_i\times\vv{p}_i + \vv{S}_i\right) + \frac{1}{c^2}\int\vv{r}\times(\vv{S}_{\mathrm{GEM}}+\vv{S}_{\mathrm{EM}})\,\dd^3x,
\end{equation}
where $\vv{S}_i$ is the spin angular momentum of particle $i$.  In 2D only the $z$-component survives, with $\vv{S}_i \to s_i\hat{z}$.  For isolated systems this should be conserved; systematic drift indicates an error in the spin precession or the spin-field coupling.

\subsubsection*{EM effective sources and the background field}

When a charged background (Kerr-Newman-like) is present, its EM field has energy density $u_{\mathrm{EM}}^{\mathrm{bg}}$ that gravitates.  In the exact Kerr-Newman solution this is already incorporated into the background metric --- the gravitational potentials $\Phig^{\mathrm{bg}}$, $\Ag^{\mathrm{bg}}$ reflect the mass $M$ plus the EM field energy, not $M$ alone.  Feeding the background EM energy density into the perturbation GEM solver as $\rho_{\mathrm{EM}}$ would therefore double-count its gravitational effect.

The correct rule is to use only the perturbation potentials in the effective GEM sources (Eqs.~\ref{eq:rho_em_pot}--\ref{eq:J_em_pot}):
\begin{align}
  \rho_{\mathrm{EM}} &= \frac{\varepsilon_0}{2c^2}\!\left(|\nabla\phi^{\mathrm{pert}} + \partial_t\vv{A}^{\mathrm{pert}}|^2
    + c^2|\nabla\times\vv{A}^{\mathrm{pert}}|^2\right).
\end{align}
Cross terms involving background potentials are second-order small for weakly charged backgrounds and are neglected.  For the pedagogical scenarios the sandbox is designed for --- a Kerr-Newman background precessing gyroscopes, curving charged particle paths, and producing Larmor precession alongside gravitomagnetic precession --- this approximation is entirely adequate.

\subsubsection*{Practical monitoring}

Computing the full field integrals at every timestep is an $O(N^3)$ operation in 3D ($O(N^2)$ in 2D).  A practical compromise is to compute them every $M$ steps (say $M=100$) and store the time series.  The field integrals can be accumulated in the same WASM loop as the field update with minimal overhead since field values are already in cache.

The most useful single diagnostic is the total canonical momentum $\vv{P}_{\mathrm{total}}$: it should be small and non-drifting for any symmetric simulation.  A non-zero systematic drift indicates either a force asymmetry (bug) or genuine momentum carried away by asymmetric radiation (physics).

\subsubsection*{Benchmark validation}

The most reliable check that the simulation is physically correct is comparison against known analytic results:
\begin{itemize}
  \item A circular orbit should have period $T = 2\pi r/v$ with $v$ determined by the force law.
  \item A gyroscope in the Kerr-like background should precess at rate $\dd\phi/\dd t = -2GJ/(c^2r^2)$.
  \item The Shapiro delay for a test signal past the background mass should match $\Delta t = -2\int\Phig\,\dd\ell/c^3$.
  \item The proper time accumulated along a circular orbit should match the analytic expression involving both kinematic and gravitational time dilation.
\end{itemize}
These provide rigorous point checks against exact results and are more diagnostic than energy monitoring for catching subtle physics errors.

\subsection{Artificial constants}

The physical value of $G/c^4$ makes gravitational effects invisible at any reasonable simulation scale.  An effective $G_{\mathrm{eff}}$ is chosen to make effects visible, as described in Section~\ref{sec:units}.  The CFL constraint $c_{\mathrm{eff}}\,\dd t < \dd x/\sqrt{2}$ in 2D determines the maximum simulation speed given the chosen $c_{\mathrm{eff}}$ and $\dd x$.

%====================================================================
\section{Two-Dimensional Reduction and Pedagogical Goals}
%====================================================================

\subsection{Motivation}

The primary goal of the sandbox is to serve as a teaching tool for general relativity.  A 2D simulation ($2+1$ dimensions) scales as $N^2$ rather than the $N^3$ of 3D, runs at interactive frame rates in JavaScript or WebAssembly, and retains essentially all physically interesting GR effects.  The 3D extension is a natural future direction once the 2D browser version is established.

\subsection{Reduction of fields and spin}

In $2+1$ dimensions:
\begin{itemize}
  \item $\vv{g}$ and $\vv{E}$ become 2-vectors ($x$ and $y$ components only).
  \item $\vv{B}_g$ and $\vv{B}$ reduce to scalars $B_{g,z}$ and $B_z$ (the only non-vanishing curl components).
  \item Cross products reduce to scalar multiplications: $(v\times B_g)_x = v_y B_{g,z}$, $(v\times B_g)_y = -v_x B_{g,z}$.
  \item The spin 3-vector $\vv{S}$ reduces to a scalar angle $\phi_{\mathrm{spin}}$ with precession:
  \begin{equation}
    \frac{\dd\phi_{\mathrm{spin}}}{\dd t} = B_{g,z} + \frac{g_s q}{2m}B_z + \Omega_{\mathrm{Thomas},z}.
  \end{equation}
  The Tulczyjew--Dixon constraint is automatically satisfied in 2D---no correction step is needed.
\end{itemize}

\subsection{The 2D Green's function}

In strict 2D the static Green's function is logarithmic:
\begin{equation}
  \text{3D:}\; \Phi\sim 1/r \;\Rightarrow\; |\vv{g}|\sim 1/r^2; \qquad
  \text{2D:}\; \Phi\sim \ln r \;\Rightarrow\; |\vv{g}|\sim 1/r.
\end{equation}
This produces flat rotation curves.  If 3D-like $1/r^2$ gravity is preferred, prescribe $\Phig^{\mathrm{bg}} = -GM/r$ analytically for the background---the FDTD perturbation solver is unchanged.

\subsection{Gravitational waves in 2D}

In pure $2+1$ GR the vacuum Einstein equations have no propagating tensor degrees of freedom---the Riemann tensor is fully determined by the Ricci tensor, which vanishes in vacuum.  However, this does not mean masses affect each other instantaneously.  The scalar and vector potential wave equations in the GEM approximation do propagate at $c$, carrying causal influence between masses.  A mass that accelerates sends a wavefront outward at $c$; particles beyond the wavefront are not affected until it reaches them.  The difference from 3D is the \emph{polarization structure}: in 2D only the scalar (compression/dilation) wave mode propagates; the tensor $+$ and $\times$ polarizations require two independent transverse directions and do not exist.

\subsection{Effects demonstrable in the 2D sandbox}

\begin{center}
\begin{tabular}{@{}L{4.5cm}L{2cm}L{6.5cm}@{}}
\toprule
\textbf{Effect} & \textbf{In 2D?} & \textbf{Notes} \\
\midrule
Newtonian gravity          & Yes & Logarithmic; background can use $1/r$ \\
Gravitational wave propagation & Yes & Scalar mode only \\
Frame dragging             & Yes & $B_{g,z}$ from spinning/moving masses \\
Perihelion precession      & Yes & $\order{v^2/c^2}$ corrections \\
Shapiro delay              & Yes & Via $c_{\mathrm{local}}$ \\
Gravitational lensing      & Yes & Via $c_{\mathrm{local}}$ \\
Spin precession            & Yes & Scalar $\phi_{\mathrm{spin}}$; directly visible \\
Proper time dilation       & Yes & Per-particle clock \\
Dipole radiation           & Yes & Time-varying spin source \\
EM--gravity coupling       & Yes & Via $\rho_{\mathrm{EM}}$, $\vv{J}_{\mathrm{EM}}$ \\
Tensor GW polarizations    & No  & Require $h^{ij}$, absent in 2D GEM \\
\bottomrule
\end{tabular}
\end{center}

\subsection{Implementation in JavaScript and WebAssembly}

For a $256\times 256$ grid with 6 scalar field equations and $\sim$10 operations per cell per step:
\begin{equation}
  N_{\mathrm{ops}} \approx 6 \times 10 \times 256^2 \approx 4\times 10^6 \text{ per timestep.}
\end{equation}
Achievable at 60\,fps in pure JavaScript.  Recommended architecture:
\begin{itemize}
  \item \textbf{WASM module:} inner FDTD loop and particle integration (performance-critical); WASM SIMD processes 4 \texttt{float32} cells per instruction.
  \item \textbf{JavaScript:} user interaction and outer simulation loop.
  \item \textbf{WebGL / WebGPU:} field visualization; WebGPU compute shaders extend practical grid size to $2048\times 2048$.
\end{itemize}

\subsection{Complete 2D algorithm: all effects}
\label{sec:alg_2d}

This is the definitive per-timestep algorithm for the 2D sandbox incorporating the background spacetime split, all spin effects, Tier~3 relativistic force corrections, proper time integration, and the correct deposition kernels.

\subsubsection*{Initialization (once before simulation begins)}

\begin{algsteps}
  \item \textbf{Set background object parameters:} mass $M$, charge $Q$, angular momentum $J$, position.

  \item \textbf{Define analytic background field functions:}
  \begin{align}
    \Phig^{\mathrm{bg}}(x,y) &= -GM/r \quad\text{(or }-GM\ln r\text{ for strict 2D)}, \\
    A_{g,x}^{\mathrm{bg}},\; A_{g,y}^{\mathrm{bg}} &\quad\text{from angular momentum }J, \\
    \phi^{\mathrm{bg}}(x,y)  &= Q/(2\pi\varepsilon_0\ln r) \quad\text{(2D Coulomb)}, \\
    \vv{g}^{\mathrm{bg}}     &= -\nabla\Phig^{\mathrm{bg}}, \quad
    B_{g,z}^{\mathrm{bg}}  = \partial_x A_{g,y}^{\mathrm{bg}} - \partial_y A_{g,x}^{\mathrm{bg}}, \\
    \vv{E}^{\mathrm{bg}}     &= -\nabla\phi^{\mathrm{bg}}, \quad
    B_z^{\mathrm{bg}} = 0 \quad\text{(static background)}, \\
    \nabla B_{g,z}^{\mathrm{bg}},\; \nabla B_z^{\mathrm{bg}} &\quad\text{(analytic gradients for spin forces).}
  \end{align}

  \item \textbf{Initialize perturbation field arrays} $\Phig^{\mathrm{pert}}$, $A_{g,x}^{\mathrm{pert}}$, $A_{g,y}^{\mathrm{pert}}$, $\phi^{\mathrm{pert}}$, $A_x^{\mathrm{pert}}$, $A_y^{\mathrm{pert}}$ and their previous-timestep copies to zero (or to specified initial conditions).

  \item \textbf{Initialize particles:} for each particle $i$ set $\vv{x}_i$, $\vv{p}_i$, $m_i$, $q_i$, $g_{s,i}$, $\phi_{\mathrm{spin},i}$, $\tau_i = 0$.
\end{algsteps}

\subsubsection*{Per-timestep iteration}

\begin{algsteps}
  \item \textbf{Deposit matter and charge sources from non-background particles} using the CIC kernel $W_2$:
  \begin{align}
    \rho_{\mathrm{matter}}(\vv{x}) &= \sum_i \gamma_i m_i\,W_2(\vv{x}-\vv{x}_i), \\
    \vv{J}_{\mathrm{matter}}(\vv{x}) &= \sum_i \gamma_i m_i\vv{v}_i\,W_2(\vv{x}-\vv{x}_i), \\
    \rho_q(\vv{x}) &= \sum_i q_i\,W_2(\vv{x}-\vv{x}_i), \\
    \vv{J}_q(\vv{x}) &= \sum_i q_i\vv{v}_i\,W_2(\vv{x}-\vv{x}_i).
  \end{align}

  \item \textbf{Deposit spin dipole sources} using the quadratic B-spline kernel $W_3$:
  \begin{align}
    \mu_i    &= \frac{g_{s,i} q_i}{2m_i} s_i, \qquad \sigma_i = 2 s_i, \\
    \vv{J}_q(\vv{x}) &\mathrel{+}= \nabla_{2D}\times\!\left[\sum_i \mu_i\,W_3(\vv{x}-\vv{x}_i)\right], \\
    \vv{J}_{\mathrm{matter}}(\vv{x}) &\mathrel{+}= \nabla_{2D}\times\!\left[\sum_i \sigma_i\,W_3(\vv{x}-\vv{x}_i)\right],
  \end{align}
  where $\nabla_{2D}\times f = (\partial_y f,\,-\partial_x f)$ is the 2D curl of a scalar field, and $s_i$ is the scalar spin.

  \item \textbf{Compute EM effective GEM sources} from perturbation potentials (2D specialization of Eqs.~\ref{eq:rho_em_pot}--\ref{eq:J_em_pot}, with $\nabla\times\vv{A}^{\mathrm{pert}} \to F_{12}^{\mathrm{pert}}\hat{z}$ where $F_{12}^{\mathrm{pert}} = \partial_x A_y^{\mathrm{pert}} - \partial_y A_x^{\mathrm{pert}}$):
  \begin{align}
    \rho_{\mathrm{EM}} &= \frac{\varepsilon_0}{2c^2}\!\left(|\nabla\phi^{\mathrm{pert}} + \partial_t\vv{A}^{\mathrm{pert}}|^2
      + c^2 (F_{12}^{\mathrm{pert}})^2\right), \\
    \vv{J}_{\mathrm{EM}} &= -\frac{F_{12}^{\mathrm{pert}}}{\mu_0 c^2}
      \begin{pmatrix} \partial_y\phi^{\mathrm{pert}} + \partial_t A_y^{\mathrm{pert}} \\ -(\partial_x\phi^{\mathrm{pert}} + \partial_t A_x^{\mathrm{pert}}) \end{pmatrix}.
  \end{align}
  (Background potentials are excluded to avoid double-counting; Section~\ref{sec:energy}.)

  \item \textbf{Sum effective GEM sources:}
  \begin{equation}
    \rho_{\mathrm{eff}} = \rho_{\mathrm{matter}} + \rho_{\mathrm{EM}}, \qquad
    \vv{J}_{\mathrm{eff}} = \vv{J}_{\mathrm{matter}} + \vv{J}_{\mathrm{EM}}.
  \end{equation}

  \item \textbf{Step perturbation GEM wave equations} (leapfrog; 3 scalar equations):
  \begin{align}
    (\Phig^{\mathrm{pert}})^{n+1} &= 2(\Phig^{\mathrm{pert}})^n - (\Phig^{\mathrm{pert}})^{n-1}
       + (c\,\dd t)^2\!\left[\nabla_{2D}^2(\Phig^{\mathrm{pert}})^n - 4\pi G\rho_{\mathrm{eff}}^n\right], \\
    (A_{g,\alpha}^{\mathrm{pert}})^{n+1} &= 2(A_{g,\alpha}^{\mathrm{pert}})^n - (A_{g,\alpha}^{\mathrm{pert}})^{n-1}
       + (c\,\dd t)^2\!\left[\nabla_{2D}^2(A_{g,\alpha}^{\mathrm{pert}})^n
       - \frac{4\pi G}{c^2}(J_{\mathrm{eff},\alpha})^n\right], \;\alpha\in\{x,y\}.
  \end{align}

  \item \textbf{Evaluate total gravitational potential} at each grid point for $c_{\mathrm{local}}$:
  \begin{equation}
    \Phig^{\mathrm{total}}(\vv{x}) = \Phig^{\mathrm{bg}}(\vv{x}) + (\Phig^{\mathrm{pert}})^{n+1}(\vv{x}),
  \end{equation}
  \begin{equation}
    c_{\mathrm{local}}(\vv{x}) = c\!\left(1 + \frac{2\Phig^{\mathrm{total}}(\vv{x})}{c^2}\right).
  \end{equation}

  \item \textbf{Step perturbation EM wave equations} (3 scalar equations) using $c_{\mathrm{local}}$:
  \begin{align}
    (\phi^{\mathrm{pert}})^{n+1} &= 2(\phi^{\mathrm{pert}})^n - (\phi^{\mathrm{pert}})^{n-1}
       + (c_{\mathrm{local}}\,\dd t)^2\!\left[\nabla_{2D}^2(\phi^{\mathrm{pert}})^n + (\rho_q^n)/\varepsilon_0\right], \\
    (A_\alpha^{\mathrm{pert}})^{n+1} &= 2(A_\alpha^{\mathrm{pert}})^n - (A_\alpha^{\mathrm{pert}})^{n-1}
       + (c_{\mathrm{local}}\,\dd t)^2\!\left[\nabla_{2D}^2(A_\alpha^{\mathrm{pert}})^n
       - \mu_0 (J_{q,\alpha})^n\right], \;\alpha\in\{x,y\}.
  \end{align}

  \item \textbf{Evaluate potentials and on-the-fly curls at each particle position,} combining analytic background and grid perturbation:
  \begin{align}
    (\Phig)_{i}^{\mathrm{total}}    &= \Phig^{\mathrm{bg}}(\vv{x}_i) + (\Phig^{\mathrm{pert}})^n_{W_2}(\vv{x}_i), \\
    (A_{g,\alpha})_i            &= A_{g,\alpha}^{\mathrm{bg}}(\vv{x}_i) + (A_{g,\alpha}^{\mathrm{pert}})^n_{W_2}(\vv{x}_i), \quad\alpha\in\{x,y\}, \\
    (\partial_t A_{g,\alpha})_i &= \frac{(A_{g,\alpha}^{\mathrm{pert}})^{n+1} - (A_{g,\alpha}^{\mathrm{pert}})^{n-1}}{2\dd t}\bigg|_{W_2}(\vv{x}_i), \\
    B_{g,z,i}                   &= B_{g,z}^{\mathrm{bg}}(\vv{x}_i) + (B_{g,z}^{\mathrm{pert}})_{\mathrm{curl},W_2}(\vv{x}_i),
  \end{align}
  and identically for the EM potentials $\phi_i$, $(A_\alpha)_i$, $(\partial_t A_\alpha)_i$, $B_{z,i}$.  The on-the-fly curl values $B_{g,z,i}$ and $B_{z,i}$ are computed from the four cell centers surrounding the particle via Eq.~(\ref{eq:yee_curl}) and bilinearly interpolated.  The background curl $B_{g,z}^{\mathrm{bg}}$ is evaluated analytically.  Also evaluate the background-referenced gradient vectors:
  \begin{align}
    \vv{g}_i^{\mathrm{bg}}    &= -\nabla\Phig^{\mathrm{bg}}(\vv{x}_i), \\
    \nabla B_{g,z,i}          &= \nabla B_{g,z}^{\mathrm{bg}}(\vv{x}_i) + (\nabla B_{g,z}^{\mathrm{pert}})_{W_3}(\vv{x}_i), \\
    \nabla B_{z,i}            &= \nabla B_z^{\mathrm{bg}}(\vv{x}_i) + (\nabla B_z^{\mathrm{pert}})_{W_3}(\vv{x}_i).
  \end{align}

  \item \textbf{Compute total force on each particle} (2D, Tier~3, potential form):
  \begin{eqbox}
  \begin{align*}
    g_{x,i} &= -\partial_x(\Phig)_{i} - (\partial_t A_{g,x})_i, \quad g_{y,i} = -\partial_y(\Phig)_{i} - (\partial_t A_{g,y})_i, \\[4pt]
    F_{x,i} &= m_i\!\left[g_{x,i}
                         + 4v_{y,i}B_{g,z,i}
                         + \frac{v_i^2}{c^2}g_{x,i}
                         + \frac{4(\Phig)_i}{c^2}g_{x,i}
                         - \frac{4(\vv{v}_i\cdot\vv{g}_i)v_{x,i}}{c^2}\right. \\
            &\qquad\quad \left.+ 4\partial_x(\vv{v}_i\cdot(\Ag)_{i}) - 4(\vv{v}_i\cdot\nabla)A_{g,x,i}\right] \\
            &\quad + q_i\!\left[E_{x,i} + v_{y,i}B_{z,i}
                         + \partial_x(\vv{v}_i\cdot\vv{A}_i) - (\vv{v}_i\cdot\nabla)A_{x,i}\right] \\
            &\quad + s_i\,\partial_x B_{g,z,i} + \mu_i\,\partial_x B_{z,i}, \\[6pt]
    F_{y,i} &= m_i\!\left[g_{y,i}
                         - 4v_{x,i}B_{g,z,i}
                         + \frac{v_i^2}{c^2}g_{y,i}
                         + \frac{4(\Phig)_i}{c^2}g_{y,i}
                         - \frac{4(\vv{v}_i\cdot\vv{g}_i)v_{y,i}}{c^2}\right. \\
            &\qquad\quad \left.+ 4\partial_y(\vv{v}_i\cdot(\Ag)_{i}) - 4(\vv{v}_i\cdot\nabla)A_{g,y,i}\right] \\
            &\quad + q_i\!\left[E_{y,i} - v_{x,i}B_{z,i}
                         + \partial_y(\vv{v}_i\cdot\vv{A}_i) - (\vv{v}_i\cdot\nabla)A_{y,i}\right] \\
            &\quad + s_i\,\partial_y B_{g,z,i} + \mu_i\,\partial_y B_{z,i},
  \end{align*}
  \end{eqbox}
  where $E_{x,i} = -\partial_x\phi_i - (\partial_t A_x)_i$ and $E_{y,i} = -\partial_y\phi_i - (\partial_t A_y)_i$ are computed on-the-fly, and $\vv{v}_i\cdot\vv{g}_i = v_{x,i}g_{x,i} + v_{y,i}g_{y,i}$.  Evaluate $v_i^2/c^2$ and $\vv{v}_i\cdot\vv{g}_i$ using $\vv{v}_i^{n-1/2}$. The vector potential terms $\partial_\alpha(\vv{v}\cdot\Ag)$ and $(\vv{v}\cdot\nabla)A_{g,\alpha}$ encode the gravitomagnetic cross-product force and are computed from spatial differences of edge potentials at the particle's CIC stencil.

  \item \textbf{Update momentum and position} (Boris pusher, 2D):
  \begin{align}
    \vv{v}_i^{\mathrm{old}} &= \vv{v}_i^{n-1/2}, \\
    \vv{p}_i^{n+1/2} &\leftarrow \text{Boris}(\vv{p}_i^{n-1/2},\; F_{x,i},\; F_{y,i},\; B_{g,z,i},\; B_{z,i},\; \dd t), \\
    \vv{v}_i^{n+1/2} &= \vv{p}_i^{n+1/2} / \sqrt{m_i^2 + |\vv{p}_i^{n+1/2}|^2/c^2}, \\
    \vv{x}_i^{n+1}   &\leftarrow \vv{x}_i^n + \vv{v}_i^{n+1/2}\,\dd t.
  \end{align}

  \item \textbf{Update spin} (scalar in 2D):
  \begin{align}
    \vv{a}_i &= (\vv{v}_i^{n+1/2} - \vv{v}_i^{\mathrm{old}})/\dd t, \\
    \Omega_{\mathrm{Thomas},z,i} &= -\frac{\gamma_i^2}{\gamma_i+1}\,\frac{v_{x,i}a_{y,i} - v_{y,i}a_{x,i}}{c^2}, \\
    \Omega_{z,i} &= B_{g,z,i} + \frac{g_{s,i}q_i}{2m_i}B_{z,i} + \Omega_{\mathrm{Thomas},z,i}, \\
    \phi_{\mathrm{spin},i} &\leftarrow \phi_{\mathrm{spin},i} + \Omega_{z,i}\,\dd t.
  \end{align}
  (No constraint correction step is needed in 2D.)

  \item \textbf{Update proper time:}
  \begin{equation}
    \beta_i^2 = \tfrac{1}{2}\!\left(|\vv{v}_i^{\mathrm{old}}|^2 + |\vv{v}_i^{n+1/2}|^2\right)/c^2, \qquad
    \dd\tau_i = \dd t\,\sqrt{1 + \frac{2(\Phig^{\mathrm{total}})_i}{c^2}
                             - \left(1-\frac{2(\Phig^{\mathrm{total}})_i}{c^2}\right)\beta_i^2},
  \end{equation}
  \begin{equation}
    \tau_i \leftarrow \tau_i + \dd\tau_i.
  \end{equation}
  $(\Phig^{\mathrm{total}})_i = \Phig^{\mathrm{bg}}(\vv{x}_i^{n+1}) + (\Phig^{\mathrm{pert}})^{n+1}(\vv{x}_i^{n+1})$, evaluated analytically and by grid interpolation respectively.
\end{algsteps}

\subsection{Pedagogical design notes}

\begin{itemize}
  \item \textbf{Separate field visualizations:} display $\Phig$, $B_{g,z}$, $\vv{E}$, $B_z$ independently.
  \item \textbf{Slow-motion and pause:} continuously adjustable speed; single-frame stepping.
  \item \textbf{Particle display:} velocity vector, spin angle indicator, GEM force vector, EM force vector, proper time readout.
  \item \textbf{Preloaded scenarios:} binary orbit, spinning mass with test gyroscopes, gravitational lens, EM wave scattering.
  \item \textbf{Adjustable constants:} sliders for $G_{\mathrm{eff}}$, $c_{\mathrm{eff}}$, $q/m$ allow students to explore the Newtonian limit ($c\to\infty$), strong-field limit, and EM-dominated regime.
\end{itemize}

%====================================================================
\section{Visualization}
%====================================================================

The visualization layer translates the numerical state of the simulation into something a student can see and understand intuitively.  It is as important as the physics for a pedagogical tool: a correct simulation that displays nothing meaningful teaches nothing.  This section describes the recommended particle and field visualizations, organized by what physical quantity they represent.  All visualizations should be independently toggleable so a student can isolate the effect they are studying.

\subsection{Particle visualizations}

Each particle is displayed as a small icon at its current position.  The icon serves as an anchor for several overlaid indicators, each encoding a different aspect of the particle's state.

\subsubsection*{Trajectory tracks}

A fading trail of the particle's past positions is drawn, length corresponding to a user-adjustable number of recent timesteps.  For orbital motion this produces a visible arc or closed ellipse.  The trail should fade in opacity with age so the most recent position is brightest and the oldest is nearly invisible, giving an immediate sense of the particle's direction of travel and speed.

At regularly spaced intervals along the track, small tick marks are drawn with an associated time readout.  Two modes are useful:
\begin{itemize}
  \item \textbf{Coordinate time ticks}: equally spaced in simulation time $t$, showing when the particle was at each point in coordinate time.  Tick spacing is uniform along the track for a particle at constant speed.
  \item \textbf{Proper time ticks}: equally spaced in proper time $\tau_i$.  For a fast-moving or deep-potential particle, proper time ticks are closer together (more time elapses per tick in coordinate time) than coordinate time ticks, making relativistic time dilation directly visible as a non-uniform spacing along the track.
\end{itemize}
Displaying both simultaneously with different symbols (e.g.\ circles for coordinate time, triangles for proper time) makes the twin paradox and gravitational time dilation immediately readable from the track geometry.

\subsubsection*{Velocity arrow}

A short arrow from the particle center in the direction of $\vv{v}_i$, with length proportional to $|\vv{v}_i|/c_{\mathrm{eff}}$ (so the arrow reaches a standard reference length at $v = c_{\mathrm{eff}}$).  Color can encode $\gamma_i$ --- blue for nonrelativistic, red for relativistic.

\subsubsection*{Spin phase indicator}

A small icon resembling a spinning top or gyroscope overlaid on the particle, with an arrow pointing in the direction $(\cos\phi_{\mathrm{spin}}, \sin\phi_{\mathrm{spin}})$ encoding the current spin phase $\phi_{\mathrm{spin}}$.  As the particle moves through a gravitomagnetic or magnetic field the arrow rotates, directly visualizing Lense--Thirring precession, geodetic precession, Larmor precession, and Thomas precession as a continuous rotation of the indicator.

The rotation rate of the arrow is $\Omega_{z,i} = B_{g,z,i} + (g_s q_i/2m_i)B_{z,i} + \Omega_{\mathrm{Thomas},z,i}$, so regions of strong gravitomagnetic or magnetic field produce visible rapid precession and field-free regions produce no rotation.  This is one of the most pedagogically direct visualizations in the sandbox --- frame dragging becomes literally visible as a spinning arrow.

\subsubsection*{Magnetic dipole indicator}

A separate icon resembling a small bar magnet (a rectangle with N and S poles indicated) that rotates with the magnetic dipole moment direction $\boldsymbol{\mu}_i = g_{s,i}(q_i/2m_i)\vv{S}_i$.  In 2D this is the same angle as the spin indicator but with a different icon to emphasize the physical distinction between the spin angular momentum (gyroscope icon) and the magnetic dipole moment (magnet icon).  The magnet icon makes it visually clear that the dipole moment is what couples to $\vv{B}$ for the Larmor force, while the gyroscope icon is what couples to $\vv{B}_g$ for the gravitomagnetic precession.

The magnitude $|\boldsymbol{\mu}_i|$ is constant (spin magnitude is conserved), which can be indicated by fixing the length of the bar magnet icon.

\subsubsection*{Force arrows}

One or more arrows showing the force components acting on each particle:
\begin{itemize}
  \item \textbf{Gravitational force} $\vv{F}_{\mathrm{grav}}$: the GEM Lorentz force plus relativistic corrections.  Color: e.g.\ orange.
  \item \textbf{Electromagnetic force} $\vv{F}_{\mathrm{EM}} = q(\vv{E}+\vv{v}\times\vv{B})$.  Color: e.g.\ blue.
  \item \textbf{Total force} $\vv{F}_{\mathrm{total}}$.  Color: e.g.\ white.
  \item \textbf{Spin gradient force} $\vv{F}_{\mathrm{spin}} = \nabla(s\,B_{g,z}) + \mu\nabla B_z$.  Color: e.g.\ green.
\end{itemize}
Arrow length is proportional to force magnitude, normalized to a reference force scale.  Individual components can be toggled, allowing a student to isolate which force is responsible for a given orbital feature.

\subsubsection*{Proper time clock}

A small analog clock face next to each particle, with a hand rotating at rate $\dd\tau_i/\dd t$ relative to coordinate time.  When $\dd\tau_i/\dd t < 1$ (relativistic motion or deep potential well) the hand rotates visibly more slowly than a reference clock displayed at a fixed point on the screen.  The running total $\tau_i$ can also be displayed as a digital readout beside the clock face.

\subsection{Field visualizations}

The perturbation field quantities available for visualization are listed below.  The scalar and vector potentials are stored on the grid and available at all times.  The derived fields $\vv{g}$, $\vv{E}$, $B_{g,z}$, $B_z$ are computed from the potentials on demand at display time --- they are not persistent grid arrays.  The per-step derivations are:
\begin{align}
  \vv{g}^{\mathrm{pert}}    &= -\nabla\Phig^{\mathrm{pert}} - \partial_t\Ag^{\mathrm{pert}}, &
  B_{g,z}^{\mathrm{pert}}   &= \partial_x A_{g,y}^{\mathrm{pert}} - \partial_y A_{g,x}^{\mathrm{pert}}, \\
  \vv{E}^{\mathrm{pert}}    &= -\nabla\phi^{\mathrm{pert}} - \partial_t\vv{A}^{\mathrm{pert}},  &
  B_z^{\mathrm{pert}}       &= \partial_x A_y^{\mathrm{pert}} - \partial_y A_x^{\mathrm{pert}}.
\end{align}

\begin{center}
\begin{tabular}{@{}L{4cm}L{4.5cm}L{4cm}@{}}
\toprule
\textbf{Field} & \textbf{Symbol} & \textbf{Physical meaning} \\
\midrule
Gravitational scalar potential & $\Phig^{\mathrm{pert}}$ (stored) & Depth of gravitational well \\
Gravitomagnetic vector potential & $\Ag^{\mathrm{pert}}$ (stored) & Source of frame dragging \\
Gravitoelectric field & $\vv{g}^{\mathrm{pert}}$ (computed) & Gravitational force per unit mass \\
Gravitomagnetic field & $B_{g,z}^{\mathrm{pert}}$ (computed) & Frame dragging / precession rate \\
Electrostatic potential & $\phi^{\mathrm{pert}}$ (stored) & Electric potential per unit charge \\
EM vector potential & $\vv{A}^{\mathrm{pert}}$ (stored) & Source of magnetic effects \\
Electric field & $\vv{E}^{\mathrm{pert}}$ (computed) & Electric force per unit charge \\
Magnetic field & $B_z^{\mathrm{pert}}$ (computed) & Larmor precession rate \\
\bottomrule
\end{tabular}
\end{center}

The background contributions $\Phig^{\mathrm{bg}}$, $B_{g,z}^{\mathrm{bg}}$, $\vv{E}^{\mathrm{bg}}$, $B_z^{\mathrm{bg}}$ can optionally be added to give the total field, or displayed separately.

\subsubsection*{Scalar field display modes}

Scalar fields ($\Phig$, $\phi$, $B_{g,z}$, $B_z$) can be displayed in three modes:

\begin{itemize}
  \item \textbf{Grayscale or false color magnitude}: the field value is mapped to a color via a diverging colormap (e.g.\ blue for negative, white for zero, red for positive).  This gives an immediate spatial impression of the field structure.  The colormap range should be auto-scaled to the current field maximum, or fixed to a user-specified range.

  \item \textbf{Contour lines}: isosurfaces of constant field value drawn as lines.  Contours of $\Phig$ are the gravitational equipotential surfaces; for a static mass they are concentric circles.  Contours of $B_{g,z}$ show the spatial structure of the frame-dragging field.  Spacing between contours encodes field gradient --- closely spaced contours indicate strong fields.

  \item \textbf{Animated phase}: for oscillating or wave fields, the field value can be mapped to a slowly cycling colormap so that wavefronts appear to propagate visually even when the field amplitude is small.  This makes gravitational and EM waves visible at low amplitude that would be invisible in a static colormap.
\end{itemize}

\subsubsection*{Vector field display modes}

Vector fields ($\vv{g}$, $\vv{E}$, $\Ag$, $\vv{A}$) can be displayed in two modes:

\begin{itemize}
  \item \textbf{Quiver plot}: arrows at a subsampled grid of points (e.g.\ every 8 cells), each pointing in the field direction with length proportional to field magnitude.  Length is typically logarithmically scaled to handle the large dynamic range near particle positions.  Color can encode magnitude independently of arrow length for additional clarity.

  \item \textbf{Streamlines / field lines}: curves everywhere tangent to the field direction, integrated from seed points.  For $\vv{g}$ these are the gravitational field lines (analogous to electric field lines), which originate at mass sources and point toward them.  For $\Ag$ the field lines show the curl structure that generates $\vv{B}_g$.  Streamlines give a cleaner global picture than quiver plots for smooth fields, but are less informative near particle positions where the field changes rapidly.
\end{itemize}

\subsubsection*{Dynamic range and normalization}

Near a particle the fields diverge as $1/r$ or $1/r^2$, making naive display of the full field dominated by the near-field singularity and showing nothing useful in the rest of the domain.  The following normalizations are recommended:

\begin{itemize}
  \item \textbf{Clamp}: cap the displayed value at a user-adjustable maximum, saturating the color near particles.  Simple but loses information near the source.
  \item \textbf{Logarithmic scale}: display $\mathrm{sgn}(\Phi)\log(1 + |\Phi|/\Phi_0)$ where $\Phi_0$ is a reference scale.  Shows both the near-field and far-field structure simultaneously.
  \item \textbf{Show perturbation only}: display $\Phig^{\mathrm{pert}}$ rather than $\Phig^{\mathrm{total}}$, removing the large static background.  This makes propagating waves visible even in the presence of a dominant static field.  This is the most useful mode for demonstrating gravitational wave emission.
\end{itemize}

\subsubsection*{Recommended display combinations for specific demonstrations}

\begin{center}
\begin{tabular}{@{}L{4cm}L{8cm}@{}}
\toprule
\textbf{Demonstration} & \textbf{Recommended display} \\
\midrule
Newtonian orbits & Contours of $\Phig^{\mathrm{total}}$ + particle tracks with coordinate time ticks \\
Frame dragging & False color $B_{g,z}$ + gyroscope spin indicators on test particles \\
Gravitational waves & Animated false color $\Phig^{\mathrm{pert}}$ (perturbation only) + particle tracks \\
Gravitational lensing & Streamlines of $\vv{E}^{\mathrm{pert}}$ (EM wave) on false color $\Phig^{\mathrm{bg}}$ \\
Larmor precession vs Lense--Thirring & False color $B_{g,z}$ and $B_z$ side by side + magnet and gyroscope icons \\
Twin paradox & Particle tracks with both coordinate and proper time ticks \\
EM--gravity coupling & False color $\Phig^{\mathrm{pert}}$ with EM source active, showing EM energy gravitating \\
Charged particle orbits & Quiver $\vv{g}$ and $\vv{E}$ + force arrows split by GEM and EM components \\
\bottomrule
\end{tabular}
\end{center}

\subsubsection*{Implementation notes}

All field arrays are already on the GPU as WebGL textures (one float per cell for scalar fields, two floats per cell for 2D vector fields).  Rendering is a fragment shader that maps the texture value to a color via the chosen colormap --- a single draw call per displayed field.  Quiver plots are rendered as a second pass drawing arrow geometry at subsampled positions, reading field values from the texture.  Particle icons, spin indicators, and force arrows are rendered as sprite geometry in a final pass.  The separation into three render passes (fields, quivers, sprites) allows each to be toggled independently with no performance cost when disabled.

%====================================================================
\section{Implementation Plan}
\label{sec:impl}
%====================================================================

This section describes a staged implementation strategy.  Each stage adds one capability, is independently testable against a known analytic result, and must pass its tests before the next stage begins.  This prevents bugs introduced early from manifesting as wrong answers late in the stack when the cause is difficult to trace.

The governing principle throughout: \textbf{never proceed to the next stage until the current stage passes its quantitative tests.}  A simulation that produces plausible-looking output is not the same as a correct simulation.

\subsection{Stage 1: Scalar wave equation, free propagation}

Implement the leapfrog update for a single scalar potential $\Phi$ in 2D on a uniform grid, with no sources and no damping boundary.

\textbf{Test:} place a narrow Gaussian initial condition at the center of the domain.  The pulse should propagate outward as a circular ring at speed $c_{\mathrm{eff}}$.  Verify:
\begin{itemize}
  \item The wavefront reaches radius $r$ at time $t = r/c_{\mathrm{eff}}$ to within one timestep.
  \item The amplitude falls off as $1/\sqrt{r}$ in 2D.
  \item The scheme is stable at $c_{\mathrm{CFL}} = 1/\sqrt{2}$ and unstable at $c_{\mathrm{CFL}} = 1/\sqrt{2} + \epsilon$.
  \item Measured numerical dispersion matches the analytic prediction of Section~\ref{sec:dispersion}.
\end{itemize}

\subsection{Stage 2: Absorbing boundary layer}

Add the damping layer (Section~\ref{sec:abc}) to Stage~1.

\textbf{Test:} the Gaussian pulse from Stage~1 should exit the domain cleanly.  Record the residual field amplitude at the domain center after the pulse should have fully exited.  Per the realized-reflection paragraph added to §\ref{sec:abc} in v33, for a 2D circular wavefront the in-interior residual is degraded from the plane-wave $R\approx10^{-3}$ limit by linearization and corner-trapping effects; a few-percent central residual with $\sim20\times$ suppression versus the no-damping case is the realistic target.

\subsection{Stage 3: Static CIC source, Poisson convergence}

Add a fixed point source using CIC deposition at the domain center.  Run the convergence iteration (positions and velocities held fixed) until fields converge.

\textbf{Test:} the scalar potential should converge to the 2D Green's function.  Plot $\Phi(r)$ vs $\ln r$ and verify linearity.  Verify convergence takes $\sim N\sqrt{2}$ timesteps.  Verify the discrete Poisson equation $\nabla^2_{h}\Phi = -\mathrm{source}$ is satisfied to machine precision at convergence.

\subsection{Stage 4: All six wave equations, gauge monitoring}

Extend to all six potentials ($\Phig$, $A_{g,x}$, $A_{g,y}$, $\phi$, $A_x$, $A_y$) simultaneously.  Add the discrete Lorenz gauge residual monitor.

\textbf{Test:} a static point source with zero velocity should produce $\Ag = 0$ and $\vv{A} = 0$ automatically (zero current sources).  Verify the gauge residual $\langle\mathcal{G}^2\rangle^{1/2}$ remains at the truncation error level $(k\Delta x)^2$.  Enable parabolic cleaning and verify it reduces the residual without destabilizing.

\subsection{Stage 5: Moving source, vector potential, Yee curl}

Initialize a source with nonzero velocity $\vv{v}$ during the convergence phase.  Verify that $\Ag$ and $\vv{A}$ develop the correct boosted Coulomb values.

\textbf{Test:} for a source moving at velocity $v$ the vector potential should satisfy $|\vv{A}| \approx v|\phi|/c^2$ to leading order.  Compute $B_z = \partial_x A_y - \partial_y A_x$ using the Yee curl at several points and compare against the analytic result for a uniformly moving charge.

\subsection{Stage 6: Sampled background field arrays}
\label{sec:stage_bg}

Introduce a parallel set of read-only grid arrays for the background contributions to each of the six potentials:
\[
  \Phig^{\mathrm{bg}}, \quad A_{g,x}^{\mathrm{bg}}, \quad A_{g,y}^{\mathrm{bg}}, \quad \phi^{\mathrm{bg}}, \quad A_x^{\mathrm{bg}}, \quad A_y^{\mathrm{bg}}.
\]
These arrays live on the same staggered Yee grid as the perturbation potentials (§\ref{sec:numerical}), are filled \emph{once} at scenario setup, and are not touched by the leapfrog or the damping layer (§\ref{sec:abc}).  Force evaluations from Stage~7 onward compute totals as
\[
  \Phig^{\mathrm{total}}(\vv{x}) = \Phig^{\mathrm{bg}}(\vv{x}) + \Phig^{\mathrm{pert}}(\vv{x}),
\]
and similarly for $\Ag$, $\phi$, $\vv{A}$, through a single combined CIC interpolation per quantity per particle.

\subsubsection*{Motivation: testability decoupled from source-coupling}

Particle-dynamics tests (orbital period, time dilation, precession, frame dragging, cyclotron motion) require a field for the particle to move in.  In v32 those tests presumed an analytic background available at the particle position via closed-form evaluation; the alternative --- creating the field with a second body and the perturbation FDTD --- would require all of the source-deposition machinery (Stages~3--5) to be working perfectly before any single-particle test could run.  By introducing sampled background arrays as their own stage, every later particle-dynamics test runs against a prescribed field, isolating the particle integrator from the source path.

\subsubsection*{Implementation}

Add the six pointers to the simulation struct, allocated lazily on first use.  Provide a small registry of background generators that fill the arrays at scenario setup time.  The minimum required for this stage is the softened point mass:
\begin{equation}
\label{eq:bg_softened_point_mass}
  \Phig^{\mathrm{bg}}(\vv{x}) = -\frac{G_{\mathrm{eff}}M}{\sqrt{|\vv{x}-\vv{x}_0|^2 + \epsilon^2}},
\end{equation}
sampled onto the grid at scenario load, with a smoothing length $\epsilon$ of a few cells to avoid the $1/r$ singularity at the source.  Later stages add further generators as needed: uniform gravitational field ($\Phig^{\mathrm{bg}}=-g_0 y$), gravitomagnetic dipole ($A_{g,\theta}^{\mathrm{bg}}\propto J/r$ as in §\ref{sec:gem_metric}), point charge, and the full mode list collected in Stage~15 below.

\textbf{Test:}
\begin{itemize}
  \item Sampled values match Eq.~(\ref{eq:bg_softened_point_mass}) at interior cells to within floating-point precision.
  \item Finite-difference gradient $-\nabla_h\Phig^{\mathrm{bg}}$ at sample radii $r\gg\epsilon$ matches the analytic Newtonian field $G_{\mathrm{eff}}M\,\hat r/r^2$ to within the CIC interpolation tolerance.
  \item Re-running the Stage~1 wave pulse with a non-trivial background installed produces a perturbation field bit-identical to the no-background case (perturbation FDTD does not read backgrounds).
  \item Toggling the damping layer with a background present does not alter the background array.
\end{itemize}

\subsubsection*{Caveat: discretization noise vs.\ small effects}

Sampling a closed-form background onto a discrete grid carries an $\order{(\Delta x/L)^2}$ representation error where $L$ is the local feature scale.  For \emph{coarse} tests (Keplerian period, gravitomagnetic precession at moderate $r$) this error is negligible.  For \emph{small} effects whose signature competes with the discretization noise --- most notably the post-Newtonian perihelion precession of Stage~8, where the per-orbit advance is $\sim G_{\mathrm{eff}}M/(c_{\mathrm{eff}}^2 a)$ and may approach the smoothing- and finite-difference-induced error --- an analytic-form background evaluator is provided as an optional code path that bypasses the grid for those specific tests.  See Section~\ref{sec:bg_mode}.

\subsubsection*{Analytic-background evaluation (added in v34)}
\label{sec:bg_mode}

In addition to the sampled-grid path described above (which is the default
and the only path supporting arbitrary user-installed background fields), the
simulator provides an \emph{analytic-mode} evaluation in which a small registry
of closed-form generators is queried directly at the particle's exact
position, with no grid sampling or CIC$+$FD chain.  The generator parameters
$(\vv{x}_0, G_{\mathrm{eff}}M, \epsilon, q, \vv{S}, \ldots)$ are stored in the
simulation struct at scenario setup, and the evaluator returns
$\Phig^{\mathrm{bg}}(\vv{x})$ and $\nabla\Phig^{\mathrm{bg}}(\vv{x})$
(and the corresponding vector-potential quantities for spinning / charged
variants) directly.

\textit{Why this matters.}  The sampled-grid path introduces a small
tangential force error at $\order{(\Delta x/r)^2}$ for any non-grid-aligned
particle position --- the same kernel that lets the W$_2$ deposition smear a
point source out also breaks the perfect spherical symmetry of the sampled
$\Phig^{\mathrm{bg}}$ a few cells away.  In the test-particle limit this
manifests as a slow drift in angular momentum, which over many orbits
produces a small spurious precession that contaminates the Stage~8
measurement.  The analytic path eliminates this systematic at the cost of
restricting the available backgrounds to the registered generator catalogue.

\textit{Currently registered.}
\begin{itemize}
  \item \texttt{POINT\_MASS}: softened Newtonian scalar
    $\Phig^{\mathrm{bg}}(\vv{x}) = -G_{\mathrm{eff}}M/\sqrt{|\vv{x}-\vv{x}_0|^2+\epsilon^2}$,
    with closed-form gradient
    $\nabla\Phig^{\mathrm{bg}} = G_{\mathrm{eff}}M(\vv{x}-\vv{x}_0)/(|\vv{x}-\vv{x}_0|^2+\epsilon^2)^{3/2}$.
\end{itemize}

\textit{Reserved (deferred to later stages).}
\begin{itemize}
  \item \texttt{SPINNING\_POINT\_MASS}: adds a gravitomagnetic $\Ag^{\mathrm{bg}}$
    from a point angular momentum $\vv{S}$ (Lense--Thirring; Stage~13/14).
  \item \texttt{CHARGED\_POINT\_MASS}: adds Coulomb $\phi^{\mathrm{bg}}$
    (Stage~11/14).
  \item \texttt{KERR\_NEWMAN}: combined spinning charged mass (Stage~14).
\end{itemize}

The mode is a runtime switch independent of which generator was last
installed.  Stage~7 and Stage~8 scenarios opt into the analytic path; Stage~6
explicitly exercises the sampled path so its tests remain meaningful for
arbitrary user-supplied backgrounds.

\subsection{Stage 7: Single test particle, Boris pusher, Keplerian orbit}

Add a single test particle with mass $m$ and no charge, placed in the softened-point-mass background of Stage~6 (Eq.~\ref{eq:bg_softened_point_mass}).  Implement the potential-form force (Eq.~\ref{eq:force_gem_pot}) and the Boris pusher.  No perturbation FDTD is exercised in this stage --- the particle only reads the background array.

\textbf{Test:} a circular orbit should have period $T = 2\pi r^{3/2}/\sqrt{G_{\mathrm{eff}}M}$.  Verify the total energy $\gamma mc^2 + m\Phig^{\mathrm{bg}}$ is approximately conserved over many orbits with no secular drift (Boris symplectic structure).  Verify that taking the smoothing length $\epsilon\to0$ recovers the textbook $1/r^2$ force law at $r\gg\epsilon$.

\subsection{Stage 8: Relativistic corrections, perihelion precession}
\label{sec:stage_precession}

Enable Tier~2 force corrections (Section~\ref{sec:relcorr}) --- the
test-particle limit of the Einstein--Infeld--Hoffmann 1PN equation:
\begin{equation}
\label{eq:tier2_force}
  \vv{F}_{\mathrm{grav}}/m \;=\; \vv{g}\!\left(1 + \frac{v^2}{c^2} + \frac{4\Phig}{c^2}\right)
  \;-\; \frac{4(\vv{v}\cdot\vv{g})\vv{v}}{c^2},
\end{equation}
with $\vv{g}=-\nabla\Phig$ (no $\Ag$ in the static-background case).  This is
the same expression as in the Tier~3 algorithm of
Section~\ref{sec:geodesic_expansion}; the EIH equation is correct through
$\order{(v/c)^2}$ in the coordinate acceleration, and Lense--Thirring
enters at the same 1PN order via the $\Ag$ terms (Section~\ref{sec:relcorr}).

\textbf{Test:} a moderately eccentric orbit on an analytic
softened-point-mass background should precess at the GR rate:
\begin{equation}
\label{eq:dphi_prec}
  \Delta\phi_{\mathrm{prec}} = \frac{6\pi G_{\mathrm{eff}} M}{c_{\mathrm{eff}}^2\,a(1-e^2)}
\end{equation}
per radial period, where $a$ is the semi-major axis and $e$ the eccentricity
of the measured orbit (i.e.\ $a' = (r_{\mathrm{peri}}+r_{\mathrm{apo}})/2$,
$e' = (r_{\mathrm{apo}}-r_{\mathrm{peri}})/(r_{\mathrm{apo}}+r_{\mathrm{peri}})$,
which generally differ from the Newtonian-derived inputs at 1PN).  Use the
analytic-background path (Section~\ref{sec:bg_mode}) for this test ---
the CIC$+$FD-of-sampled-background chain introduces an
$\order{(\Delta x/r)^2}$ tangential force error that breaks angular-momentum
conservation and contaminates the precession measurement.

A small ``kinematic'' baseline contribution arises from running the
relativistic 3-momentum dynamics with a Tier-0 (Newtonian) force law:
this is not Schwarzschild and does precess, by $\sim 0.1\,(v/c)^2$ rad/orbit
at our parameter scale.  Subtract it (run Tier-0 first as a baseline);
the residual EIH contribution should match Eq.~(\ref{eq:dphi_prec}) to
within the residual $\order{(v/c)^4}$ 2PN truncation.  At
$v/c \approx 0.04$ in the implementation the agreement is better than
$1\%$.

\subsection{Stage 9: Proper time accumulation}
\label{sec:stage_proper_time}

Add the proper time clock $\tau_i$ to each particle (Section~\ref{sec:proptime}).
Each timestep accumulate
\begin{equation}
\label{eq:dtau_step}
  \dd\tau = \dd t\,\sqrt{1 + \frac{2\Phig}{c^2} - \left(1-\frac{2\Phig}{c^2}\right)\frac{v^2}{c^2} - \frac{8\,\Ag\!\cdot\!\vv{v}}{c^2}},
\end{equation}
the full Eq.~(\ref{eq:proper_time_full}) with the v34 sign$+$factor
corrections.  $\Phig$ and $\Ag$ are evaluated at the particle's current
position; for static-background scenarios the analytic-mode path
(Section~\ref{sec:bg_mode}) gives a clean signal free of CIC$+$FD tangential
artifact.

\textbf{Two-phase test.}

\emph{Phase A — non-spinning point mass, several orbital radii.}  Place a
test particle on the relativistic-Newtonian circular orbit (as in Stage~7)
at $r \in \{10, 20, 40, 80\}$ with $G_{\mathrm{eff}}M = 1$, $\epsilon = 1$,
$c_{\mathrm{eff}} = 1$.  After one orbital period verify
\begin{equation}
  \left.\frac{\dd\tau}{\dd t}\right|_{\rm measured}
  \approx \sqrt{1 + \frac{2\Phig}{c^2} - \left(1-\frac{2\Phig}{c^2}\right)\frac{v^2}{c^2}}
\end{equation}
to within a relative error consistent with the leapfrog truncation
($\sim 10^{-5}$ at $r \gtrsim 10$).  For a Newtonian circular orbit with
$v^2 = G_{\mathrm{eff}}M/r$ the gravitational and kinematic contributions
to the leading-order $\dd\tau/\dd t - 1$ stand in a $2:1$ ratio
($\Phig/c^2 : -v^2/(2c^2)$); as $r$ increases the SR contribution drops
faster than the gravitational one (both in absolute value).  The
implementation recovers this ratio: at $r=10$ the measured ratio is
$2.12$, converging to $2.01$ at $r=80$.

\emph{Phase B — spinning point mass, prograde vs.\ retrograde.}  Install a
spinning-point-mass background (Section~\ref{sec:bg_mode}) with central
spin $J_z$ and place a test particle on the same circular orbit, once
prograde ($\vv{v} = +v_{\mathrm{circ}}\hat\phi$) and once retrograde
($\vv{v} = -v_{\mathrm{circ}}\hat\phi$).  Run each for the same coordinate
time $T_{\mathrm{coord}}$ (one orbital period at the non-spinning
$v_{\mathrm{circ}}$).  The proper-time difference between the two clocks is
\begin{equation}
\label{eq:clock_effect_linear}
  \Delta\tau \;=\; \tau_+ - \tau_- \;\approx\; -\frac{8\pi G_{\mathrm{eff}} J_z}{c_{\mathrm{eff}}^4\,r}
\end{equation}
per orbit at leading order (\cite{CohenMashhoon1993}, in this document's
factor-of-4-in-force GEM convention).  For exact agreement with the
implementation use
$\Delta\tau = T_{\mathrm{coord}}\!\left[(\dd\tau/\dd t)_+ - (\dd\tau/\dd
t)_-\right]$ with Eq.~(\ref{eq:dtau_step}) evaluated at the prograde and
retrograde signs of $\Ag\!\cdot\!\vv{v}$ — the linearized
Eq.~(\ref{eq:clock_effect_linear}) drops a higher-order cross-coupling
between $\Phig$, $v^2$, and $\Ag\!\cdot\!\vv{v}$ that becomes visible at
near-extremal Kerr in toy units.  In the implementation at
$r=20,\;J_z=2$ the linearized prediction is $8.4\%$ off the measured
$\Delta\tau$ while the exact-sqrt prediction agrees to better than
$0.2\%$ — the latter being the pure leapfrog integration residual.

Both phases use the analytic-background path; the perturbation FDTD is
skipped via \texttt{gr\_sim\_set\_field\_evolution(sim,~0)} since no
sources are active.

\paragraph{v35 addendum --- Phase B clock-only isolation.}  Phase B's
analytic prediction Eq.~(\ref{eq:clock_effect_linear}) is derived under
the assumption that the prograde and retrograde test particles remain on
\emph{the same} Keplerian circle, so that only the
$-8(\Ag\!\cdot\!\vv{v})/c^2$ cross-term in the $\dd\tau/\dd t$ integrand
differs between them.  Once the Tier-1 gravitomagnetic Lorentz piece
$\vv{F}_{\mathrm{gm}} = +4m\,\vv{v}\!\times\!\vv{B}_g$
(Section~\ref{sec:geodesic_expansion}, Eq.~\ref{eq:geodesic_expansion})
is wired into the pusher, that assumption no longer holds --- the same
$\vv{B}_g$ that produces the clock asymmetry also perturbs the orbit
shape oppositely for prograde and retrograde, a first-order Lense--Thirring
frame-dragging response on the dynamics.  The two effects combined would
add another $\order{J_z}$ contribution to $\Delta\tau$ on top of
Eq.~(\ref{eq:clock_effect_linear}), spoiling the comparison.  To preserve
the clock-only isolation here, the gravitomagnetic Lorentz piece is
gated by the \texttt{gr\_sim\_set\_gravitomagnetic\_force\_enabled(sim,
0)} flag for the duration of Phase B; Stage~9b
(Section~\ref{sec:stage_gm_cyclotron}) verifies that piece on its own
with $\Phig = 0$.

\subsection{Stage 9b: Gravitomagnetic Lorentz force --- unit isolation}
\label{sec:stage_gm_cyclotron}

A unit verification of the Tier-1 gravitomagnetic Lorentz force
$\vv{F}_{\mathrm{gm}}/m = +4\,\vv{v}\!\times\!\vv{B}_g$
(Eq.~\ref{eq:geodesic_expansion} at $\order{v/c}$; algorithmic form
in Section~\ref{sec:alg_rel}'s Tier-3 eqbox).  The factor of $4$ is the spin-2
enhancement of the GEM Lorentz force relative to its EM analogue
(Section~\ref{sec:gem_metric}); the sign was corrected from v33 in v34
(empirical confirmation via Stage~8 perihelion precession; the v33 sign
gave retrograde precession, the EIH sign gives prograde, matching the
canonical Schwarzschild direction).  Stage~9b verifies the coefficient
and sign \emph{directly}, isolated from the scalar gravity dynamics, by
setting $\Phig = 0$ everywhere and prescribing a spatially uniform
gravitomagnetic field.

\paragraph{Construction.}  A new analytic background generator
\texttt{gr\_sim\_set\_background\_uniform\_gravitomagnetic} installs the
symmetric-gauge vector potential
\begin{equation}
\label{eq:uniform_bg_symmetric_gauge}
  \Phig^{\mathrm{bg}} = 0,\qquad
  A_{g,x}^{\mathrm{bg}}(x,y) = -\tfrac{1}{2}B_0\,(y - y_0),\qquad
  A_{g,y}^{\mathrm{bg}}(x,y) = +\tfrac{1}{2}B_0\,(x - x_0),
\end{equation}
yielding $B_{g,z} = \partial_x A_{g,y}^{\mathrm{bg}} - \partial_y
A_{g,x}^{\mathrm{bg}} = B_0$ identically.  The gauge origin $(x_0, y_0)$
sets the zero of $\Ag$ and is otherwise physically irrelevant.

\paragraph{Analytic prediction.}  With $\Phig = 0$ the force on a test
particle reduces to $\vv{F} = m(4\,\vv{v}\!\times\!B_0\hat z)$; the
particle traces a circle (cyclotron orbit) at
\begin{equation}
\label{eq:gm_cyclotron}
  \omega_{\mathrm{gm}} = \frac{4\,|B_0|}{\gamma},\qquad
  T = \frac{2\pi}{\omega_{\mathrm{gm}}} = \frac{\pi\gamma}{2\,|B_0|},\qquad
  r_L = \frac{v_0}{\omega_{\mathrm{gm}}} = \frac{v_0\,\gamma}{4\,|B_0|}.
\end{equation}
For $B_0 > 0$ and initial velocity $\vv{v}(0) = v_0\hat x$, the initial
force is $\vv{F}_y = -4m v_0 B_0 < 0$, so the particle gyrates
clockwise (viewed from $+\hat z$).  Reversing the sign of $B_0$ reverses
the gyration direction.

\paragraph{Test (\texttt{stage20\_uniform\_bg\_gyration.c}).}  At
$v_0 = 0.1c$, $B_0 = 10^{-3}$ (so $\gamma \approx 1.005$,
$\omega_{\mathrm{gm}} \approx 3.98\times 10^{-3}$,
$T \approx 1579$, $r_L \approx 25.13$ cells on a $128\times 128$ grid):
\begin{itemize}
  \item Closure (particle returns to within $1\%$ of $r_L$ of its
    starting point after $T$): \textbf{$0.10\%$} (well under threshold;
    residual is the $\dd t \cdot \omega_{\mathrm{gm}}$ phase-rounding from
    integerizing the period in timesteps).
  \item $|\vv{v}|^2$ drift (the Lorentz force does no work):
    \textbf{$3\times 10^{-6}$}, essentially round-off.
  \item Measured Larmor radius vs.\
    Eq.~(\ref{eq:gm_cyclotron}): \textbf{$3\times 10^{-6}$}
    fractional error.
  \item Gyration direction flips under $B_0 \to -B_0$: \textbf{verified
    both signs}.
\end{itemize}
The Larmor-radius agreement to a part in $\sim 10^6$ directly fixes the
coefficient $+4$ on $\vv{v}\!\times\!\vv{B}_g$: an erroneous coefficient
$\alpha$ would give $r_L = v_0\gamma/(\alpha|B_0|)$, shifting the
measured radius linearly.

\paragraph{Notation note --- factor of 1 vs.\ 4 between
Eq.~(\ref{eq:eih_full}) and the Tier table.}
Eq.~(\ref{eq:eih_full}) is written in the textbook (lecture-25)
convention where the gravitomagnetic 4-potential carries coefficient~$1$
on $\vv{v}\!\times\!(\nabla\!\times\!\Ag)$.  The Tier table at
Section~\ref{sec:alg_rel}, Eq.~(\ref{eq:geodesic_expansion}), and the
Tier-3 implementation eqbox use the doc's own convention where the
same physical quantity carries coefficient~$4$.  The two conventions
are related by the identification
$\xi_i^{\text{(lecture)}} = 4\,A_{g,i}^{\text{(doc)}}/c$ (see the
provenance discussion immediately after Eq.~\ref{eq:proper_time_full}).
The factor-of-4 form is consistent with the spin-2 metric
$h_{0i} = 4A_{g,i}/c$, the proper-time formula
$\dd\tau/\dd t \supset -8(\Ag\!\cdot\!\vv{v})/c^2$, and the
implementation; the lecture-1 form is preserved at
Eq.~(\ref{eq:eih_full}) only for textual consistency with the cited
derivation.  All implementations use the factor-of-4 form.

\paragraph{Cross-references.}  This stage isolates the
$\order{v/c}$ Lense--Thirring force piece in a \emph{uniform} $\vv{B}_g$;
Stage~9c (Section~\ref{sec:stage_lt_period}) extends the verification
to a \emph{spatially-varying} $\vv{B}_g(r)$ around a spinning point mass.
Together with Section~\ref{sec:stage_proper_time} Phase~B (clock effect)
they exhaust the leading-order $\order{J_z}$ tests of the GEM coupling.

\subsection{Stage 9c: Lense--Thirring orbital period asymmetry}
\label{sec:stage_lt_period}

Stage~9b verified the $+4\,\vv{v}\!\times\!\vv{B}_g$ Lorentz piece against
a spatially uniform $\vv{B}_g$.  Stage~9c upgrades the verification to a
\emph{spatially varying} $\vv{B}_g(r)$ around a spinning point mass --- the
configuration of physical interest for Kerr-like dynamics --- and measures
the leading-order Lense--Thirring response: the prograde vs.\ retrograde
orbital-period asymmetry.

\paragraph{Setup.}  A non-spinning test particle on a circular orbit of
radius $r$ around the spinning-point-mass background of
Section~\ref{sec:bg_mode}.  At the orbit's radius the gravitomagnetic
field is the equatorial value of the dipole:
\begin{equation}
\label{eq:Bgz_dipole_equatorial}
  B_{g,z}(r) \;=\; \frac{G_{\mathrm{eff}}\,J_z}{2\,c^2}\,
                  \frac{2\epsilon^2 - r^2}{(r^2 + \epsilon^2)^{5/2}}
              \;\xrightarrow{r \gg \epsilon}\;
              -\frac{G_{\mathrm{eff}}\,J_z}{2\,c^2\,r^3}.
\end{equation}
The negative sign is the equatorial-plane value of a magnetic-type
dipole (antiparallel to the dipole moment by a factor 1/2 of the polar
value).  The Lorentz force $\vv{F} = 4m\,\vv{v}\!\times\!\vv{B}_g$ on
a tangential velocity is purely radial: inward for prograde
($\vv{v}\!\times\!B_{g,z}\hat z = +v\,B_{g,z}\,\hat r$, with $B_{g,z}<0$
making this $-|B_{g,z}|\hat r$), outward for retrograde.

\paragraph{Analytic prediction.}  The steady-state circular velocity
satisfies the relativistic centripetal balance
$\gamma v^2/r = a(r) \pm 4 v|B_{g,z}|$
where $a(r) = G_{\mathrm{eff}} M r/(r^2+\epsilon^2)^{3/2}$ and
$\gamma = (1-v^2/c^2)^{-1/2}$.  Linearising about the relativistic
\emph{scalar} circular velocity $v_{\mathrm{rel}}$ (defined by
$\gamma_{\mathrm{rel}} v_{\mathrm{rel}}^2 = a(r) r$; the same
$v_{\mathrm{circ}}$ used in the Stage~7 kepler\_orbit scenario) gives
\begin{equation}
\label{eq:delta_v_lt}
  \delta v \;=\; \frac{4\,|B_{g,z}|\,r}
                      {\gamma_{\mathrm{rel}}\bigl(\,2 + \gamma_{\mathrm{rel}}^2\,v_{\mathrm{rel}}^2/c^2\,\bigr)},
\end{equation}
so $v_{\mathrm{pro}} = v_{\mathrm{rel}} + \delta v$ and
$v_{\mathrm{ret}} = v_{\mathrm{rel}} - \delta v$, giving the period
asymmetry
\begin{equation}
\label{eq:dT_lt}
  T_{\mathrm{ret}} - T_{\mathrm{pro}}
  \;\approx\;  \frac{4\pi r\,\delta v}{v_{\mathrm{rel}}^2}.
\end{equation}
Using the relativistic $v_{\mathrm{rel}}$ (not Newtonian) as the
linearisation point is essential: the Newtonian
$v^2 = a(r) r$ over-predicts $v$ by $\order{v^2/c^2}$ (1.3\% at our
test parameters), which would put the integrator on a small ellipse
rather than a circle and add an $\order{v^2/c^2}$ contamination to the
measured period.

\paragraph{Test (\texttt{stage21\_lense\_thirring\_period.c}).}
At $GM=1$, $r=20$, $\epsilon=1$, $J_z=2$, $c=1$:
$B_{g,z} = -1.236\times 10^{-4}$, $v_{\mathrm{rel}} = 0.2204$,
$\gamma_{\mathrm{rel}} = 1.025$, $\delta v = 4.70\times 10^{-3}$,
analytic $T_{\mathrm{ret}} - T_{\mathrm{pro}} = +24.33$
(4.27\% of $T_{\mathrm{rel}} = 570.09$).  Measured over four orbits:
\begin{itemize}
  \item Prograde:   $T_{\mathrm{pro}} = 557.55$  ($r \in [19.98, 20.00]$).
  \item Retrograde: $T_{\mathrm{ret}} = 582.13$  ($r \in [19.99, 20.00]$).
  \item $T_{\mathrm{ret}} - T_{\mathrm{pro}} = +24.57$, agreeing with
    Eq.~(\ref{eq:dT_lt}) to \textbf{0.97\%}.
  \item Baseline (GM force disabled): prograde and retrograde periods
    agree at $570.10$ vs.\ $569.93$, matching the relativistic
    scalar prediction; the residual is the dt$\,$-rounded period-detection
    error.
\end{itemize}
The orbits stay nearly perfectly circular ($\Delta r/r \lesssim 10^{-3}$),
confirming that the perturbative $\delta v$ correction
Eq.~(\ref{eq:delta_v_lt}) captures the steady-state circular condition
in the GEM-corrected dynamics.

\paragraph{Interpretation.}  This is the leading-order Lense--Thirring
\emph{orbital} response: the spinning background's frame-dragging
torques prograde vs.\ retrograde orbits differently, shifting their
steady-state radii (and hence their periods) by $\order{J_z/(Mc^2)}$.
It is distinct from (i) the proper-time clock effect of
Section~\ref{sec:stage_proper_time} Phase~B, which is the
$\order{J_z}$ contribution to $\dd\tau/\dd t$ at fixed orbit, and
(ii) the Schwarzschild perihelion precession from the Tier-2 EIH terms,
which is direction-symmetric.  The three signals --- clock effect, GM
Lorentz force on the orbit, EIH precession --- are independent at
$\order{J_z}$ and $\order{(v/c)^2}$, and the verification stages
isolate each in turn.

\subsection{Stage 9d: Sampled-mode curl of A\_g}
\label{sec:stage_sampled_curl}

Stages~9b and 9c verified the gravitomagnetic Lorentz force with
\emph{analytic-mode} B$_g$ — closed-form curls of the uniform and
spinning-point-mass backgrounds.  Stage~9d extends the verification to
the \emph{sampled-mode} path: a 2-point Yee curl of the
\textsc{X\_EDGE} and \textsc{Y\_EDGE} potential arrays.  This is the
machinery the pusher will use once moving masses generate a
\emph{self-consistent} perturbation $\Ag^{\mathrm{pert}}$ that has no
closed form — the same code path stages 9b/9c exercise indirectly via
the background route.

\paragraph{Kernel.}  $B_{g,z}$ lives naturally on the cell-center
sublattice $(i+\tfrac12, j+\tfrac12)\Delta x$.  The discrete curl is
\begin{equation}
\label{eq:Bgz_yee_curl}
  B_{g,z}\bigl[i+\tfrac12, j+\tfrac12\bigr]
  = \frac{A_{g,y}[i+1, j] - A_{g,y}[i, j]}{\Delta x}
  - \frac{A_{g,x}[i, j+1] - A_{g,x}[i, j]}{\Delta x},
\end{equation}
using the X\_EDGE indexing of $A_{g,x}$ (sample at $(i+\tfrac12, j)\Delta x$)
and the Y\_EDGE indexing of $A_{g,y}$ (sample at $(i, j+\tfrac12)\Delta x$).
The particle's $B_{g,z}$ is the bilinear interpolation of the four
surrounding cell-center values.  The total $B_{g,z}$ at the particle
is the sum of (i) the background contribution
(analytic via $\texttt{gr\_bg\_eval\_B\_g}$ when \texttt{bg\_mode} is
\textsc{analytic}, sampled via Eq.~\ref{eq:Bgz_yee_curl} when
\textsc{sampled}), and (ii) the perturbation contribution (always
sampled, reading the \texttt{.prev} buffer when field evolution is
active to match the deposit half-step time --- the same read-prev
convention as for $\Phig$ in Stage 18's energy-conservation fix).

\paragraph{Test (\texttt{stage22\_sampled\_bg\_curl.c}).}
\begin{itemize}
  \item \emph{Phase A — uniform $\vv{B}_g$:}  Symmetric-gauge
    $\Ag$ is linear in $(x, y)$, so the 2-point Yee forward
    difference is \emph{exact}.  The sampled $B_{g,z}$ matches
    $B_0$ at every interior sample point to floating-point round-off
    ($\sim 10^{-10}$ relative).
  \item \emph{Phase B — dipole $\vv{B}_g(r)$:}  At
    $J_z = 2$, $\epsilon = 1$, $\Delta x = 1$, the relative error
    against the analytic curl is $0.62\%$ at $r=20$, $0.16\%$ at $r=40$,
    $0.04\%$ at $r=80$.  Convergence ratios: $3.97\times$ and $4.00\times$
    --- exactly the expected 2nd-order behaviour
    ($\mathrm{err} \propto (\Delta x/r)^2$).
  \item \emph{Phase C — Stage 9b cyclotron repeated with sampled
    $\vv{B}_g$:}  Larmor-radius error is $5\times 10^{-6}$, identical
    to the analytic-mode result --- linear $\Ag$ gives bit-clean
    agreement.
\end{itemize}

\paragraph{Implication.}  The same curl kernel will compute $B_{g,z}$
from the perturbation $\Ag^{\mathrm{pert}}$ that arises when a moving
mass deposits the GEM current $\vv{J}_m$ and the wave equation
propagates the response.  Stage~22 confirms that path is correct
\emph{before} the dynamics get self-consistent, so any later
discrepancy at the PIC level can be localized to deposition, smoothing,
or the wave-equation evolution, rather than to the curl evaluator
itself.

\subsection{Stage 9e: EM Lorentz force --- unit isolation}
\label{sec:stage_em_cyclotron}

EM analog of Stage~9b/9d.  The Tier-3 force-law eqbox of
Section~\ref{sec:alg_rel} contains the EM Lorentz piece
\begin{equation}
\label{eq:Fem_lorentz}
  \vv{F}_{\mathrm{em}}
  = q\bigl(-\nabla\phi - \partial_t\vv{A} + \vv{v}\!\times\!\vv{B}\bigr),
  \qquad \vv{B} = \nabla\!\times\!\vv{A},
\end{equation}
with coefficient $+1$ on $\vv{v}\!\times\!\vv{B}$ --- the EM analogue of the
GEM coefficient $+4$, but without the spin-2 enhancement.  Until the
scalar $-q\nabla\phi$ and inductive $-q\partial_t\vv{A}$ pieces land
(planned with the static-Coulomb stage), the implementation realizes the
$q\,\vv{v}\!\times\!\vv{B}$ piece alone.  Stage~9e verifies that piece in
isolation, mirroring Stage~9b but for the EM sector.

\paragraph{Construction.}  The new analytic background generator
\texttt{gr\_sim\_set\_background\_uniform\_magnetic} installs the same
symmetric-gauge pattern as the GEM case, applied to the EM potential
arrays:
\begin{equation}
\label{eq:em_uniform_bg_symmetric_gauge}
  \phi^{\mathrm{bg}} = 0,\qquad
  A_{x}^{\mathrm{bg}}(x,y) = -\tfrac{1}{2}B_0\,(y - y_0),\qquad
  A_{y}^{\mathrm{bg}}(x,y) = +\tfrac{1}{2}B_0\,(x - x_0),
\end{equation}
producing $B_z = B_0$ everywhere.  The same Yee-curl evaluator
introduced in Stage~9d serves the EM side identically --- the X\_EDGE /
Y\_EDGE staggering of $A_x, A_y$ is structurally identical to the GEM
case, so the curl kernel in
\texttt{particle.c:curl\_Agz\_sampled\_at} is shared between both.

\paragraph{Analytic prediction.}  With $\phi = 0$ and $\partial_t\vv{A}
= 0$ the force reduces to $\vv{F} = q\,(\vv{v}\!\times\!B_0\hat z)$, giving
the standard EM cyclotron:
\begin{equation}
  \omega_c = \frac{q\,|B_0|}{\gamma\,m},\qquad
  T       = \frac{2\pi}{\omega_c},\qquad
  r_L     = \frac{v_0}{\omega_c} = \frac{\gamma\,m\,v_0}{q\,|B_0|}.
\end{equation}
For $q > 0$, $B_0 > 0$, and $\vv{v}(0) = v_0\hat x$, the particle
gyrates clockwise (viewed from $+\hat z$); sign reversal of either
$q$ or $B_0$ flips the direction.

\paragraph{Test (\texttt{stage23\_em\_cyclotron.c}).}  At
$q = 1$, $m = 1$, $v_0 = 0.1c$, $B_0 = 4\times 10^{-3}$ (chosen so
$r_L \approx 25$ cells, matching the Stage~9b/9d orbit geometry on a
$128\times 128$ grid):
\begin{itemize}
  \item ANALYTIC bg: $r_L$ error $3\times 10^{-6}$,
    $|\vv{v}|^2$ drift $3\times 10^{-6}$,
    closure 0.10\% (the same $\dd t \cdot \omega_c$ phase-rounding
    floor as Stage 9b).  Gyration direction is correct.
  \item SAMPLED bg: $r_L$ error $5\times 10^{-6}$, agreeing with
    ANALYTIC to $\sim 5\times 10^{-6}$ in absolute $r_L$.  Linear
    $\vv{A}^{\mathrm{bg}}$ -> Yee curl is exact, as Stage 9d showed.
  \item Sign reversal: $B_0 \to -B_0$ flips gyration from CW to CCW
    correctly.
  \item Neutral particle ($q = 0$): feels no force (drifts linearly),
    confirming the EM Lorentz piece is correctly gated by charge.
\end{itemize}

\paragraph{Status of the rest of $\vv{F}_{\mathrm{em}}$.}
Eq.~(\ref{eq:Fem_lorentz}) has three pieces.  Stage~9e covers the
$q\,\vv{v}\!\times\!\vv{B}$ term.  Stage~9f
(Section~\ref{sec:stage_em_estatic}) adds the $-q\nabla\phi$ scalar
(electrostatic) term.  Only the $-q\,\partial_t\vv{A}$ inductive piece
remains unwired; it will land with the radiative-EM stage (where a
time-varying perturbation $\vv{A}$ is the primary observable).

\subsection{Stage 9f: Electrostatic force --- unit isolation}
\label{sec:stage_em_estatic}

Wires the $-q\nabla\phi$ piece of Eq.~(\ref{eq:Fem_lorentz}) into the
pusher.  Stage~9e covers $q\vv{v}\!\times\!\vv{B}$; Stage~9f covers
$-q\nabla\phi$; together they handle the full EM Lorentz force
\emph{for static $\vv{A}$ backgrounds} (the $-q\,\partial_t\vv{A}$
inductive piece is still pending and only matters when $\vv{A}$ is
time-varying).

\paragraph{Construction.}  A new analytic background generator
\texttt{gr\_sim\_set\_background\_uniform\_electric} installs the linear
EM scalar potential
\begin{equation}
\label{eq:em_uniform_e_bg}
  \phi^{\mathrm{bg}}(x, y) = -\bigl(E_x(x - x_0) + E_y(y - y_0)\bigr),
\end{equation}
producing $-\nabla\phi^{\mathrm{bg}} = (E_x, E_y)$ uniformly.  No
$\vv{A}^{\mathrm{bg}}$ is installed; the magnetic Lorentz piece
contributes nothing.

In the pusher, the gradient $\nabla\phi_{\mathrm{em}}^{\mathrm{total}}$
at the particle is computed by
\texttt{particle.c:phi\_em\_grad\_at\_total}, which mirrors the
grav\_grad\_at structure exactly: analytic background via
\texttt{gr\_bg\_eval\_phi\_em} (when \texttt{bg\_mode} is
\textsc{analytic}), centered finite-difference of $\phi^{\mathrm{bg}}$
on the CORNER sublattice plus CIC interpolation (when \textsc{sampled}),
and the perturbation contribution from \texttt{fields[PHI\_EM]} read
with the same \texttt{.prev} half-step convention as $\Phig$.  The
TSC and Lewis--Birdsall variants of the EM gradient kernel are not yet
ported; they will mirror the GEM-side
(Section~\ref{sec:stage_perturb_fdtd}) when self-consistent PIC
dynamics arrive on the EM side.

\paragraph{Analytic prediction.}  With a uniform $\vv{E} = E_0\hat x$
and no $\vv{B}$ or scalar gravity, a particle of charge $q$ and rest
mass $m$ starting at rest follows
\begin{align}
\label{eq:em_uniform_e_dynamics}
  p_x(t)        &= q E_0\,t, \qquad p_y(t) = 0, \\
  x(t) - x(0)   &= \frac{m c^2}{q E_0}\!\left[\sqrt{1 + \!\Bigl(\tfrac{q E_0 t}{m c}\Bigr)^{\!2}} - 1\right],
\end{align}
the exact relativistic kinematic integral.  The Boris half-step
convention stores $p^{N-1/2}$ at simulation step $N$, so the recorded
momentum at step $N$ is $p_x^{N-1/2} = q E_0 (N - \tfrac{1}{2})\Delta t$.

\paragraph{Test (\texttt{stage24\_uniform\_e\_field.c}).}  At $q = 1$,
$m = 1$, $E_0 = 10^{-3}$, $c = 1$, $\Delta t = 1/\sqrt{2}$, $N = 200$:
\begin{itemize}
  \item ANALYTIC bg:  $p_x$ matches Eq.~(\ref{eq:em_uniform_e_dynamics})
    to a relative error of $1.1\times 10^{-6}$ (float32 precision floor
    of the kinematic integral); $p_y = 0$ exactly; position drift
    matches the exact relativistic formula to $2.5\times 10^{-6}$.
  \item SAMPLED bg:  $p_x$ \emph{bit-identical} to ANALYTIC mode (linear
    $\phi^{\mathrm{bg}} \Rightarrow$ centered FD is exact, as established
    by Stage~9d for the GEM curl analogue).
  \item Sign reversal: $|p_x(E_0 > 0) + p_x(E_0 < 0)| < 10^{-8}$ ---
    exact antisymmetry to round-off.
  \item Neutral particle ($q = 0$): zero momentum and zero drift after
    200 steps, confirming the electrostatic force is correctly gated
    by charge.
\end{itemize}

\subsection{Stage 9g: Inductive force --- unit isolation}
\label{sec:stage_em_inductive}

Wires the $-q\,\partial_t\vv{A}$ piece of Eq.~(\ref{eq:Fem_lorentz})
into the pusher.  With Stages 9e ($q\,\vv{v}\!\times\!\vv{B}$) and 9f
($-q\nabla\phi$) already in place, Stage 9g completes the full EM
Lorentz force
\begin{equation*}
  \vv{F}_{\mathrm{em}} = q\bigl(-\nabla\phi - \partial_t\vv{A} + \vv{v}\!\times\!\vv{B}\bigr)
\end{equation*}
on the pusher side.

\paragraph{Temporal derivative on the leapfrog buffer rotation.}  The
field leapfrog in \texttt{sim.c:gr\_sim\_step} cycles three time
buffers per field.  After each step's rotation,
\[
  \texttt{fields[*].prev} = \vv{A}^n,\qquad
  \texttt{fields[*].curr} = \vv{A}^{n+1},\qquad
  \texttt{fields[*].next} = \vv{A}^{n-1}.
\]
The recycled \texttt{.next} slot (which the next leapfrog will overwrite)
therefore holds the \emph{previous} time slice.  This gives a clean
2nd-order centered difference at $t^n$:
\begin{equation}
\label{eq:dt_A_centered}
  \partial_t\vv{A}^n
   \;=\; \frac{\vv{A}^{n+1} - \vv{A}^{n-1}}{2\,\Delta t}
   \;=\; \frac{\texttt{curr} - \texttt{next}}{2\,\Delta t},
\end{equation}
evaluated at the same time as $\vv{B}^n = \nabla\!\times\!\vv{A}^n$
(\texttt{.prev}) and $\Phi_{\mathrm{em}}^n$ (\texttt{.prev}).  Spatial
interpolation uses the same sublattice-aware kernels that serve the
magnetic curl path: $\vv{A}_x$ on X\_EDGE, $\vv{A}_y$ on Y\_EDGE.

\paragraph{Test (\texttt{stage25\_inductive\_e.c}).}  The cleanest
isolation manually pre-populates the time buffers with a known temporal
profile.  With \texttt{field\_evolution} OFF (so the buffers don't
rotate), set
\[
  \texttt{fields[A\_X].curr}[k] = +A_0\,\Delta t,\qquad
  \texttt{fields[A\_X].next}[k] = -A_0\,\Delta t,
\]
uniformly over the grid.  The centered difference is then $A_0$ at every
position, mimicking a linear-in-time $A_x(t) = A_0\,t$.  No scalar EM
background, no magnetic field, so the only force is the inductive piece:
$F_x = -q\,\partial_t A_x = -q A_0$.  A charged particle at rest then
accelerates uniformly along $-x$ at $|q A_0|/(\gamma m)$; under the
Boris half-step convention $p_x^{N-1/2} = -q A_0 (N-\tfrac{1}{2})\Delta t$.
This is exactly the kinematics of Stage 9f (electrostatic, $-q\nabla\phi$),
but driven by the inductive force-law piece instead.

Results at $q=1$, $m=1$, $A_0 = 10^{-3}$, $N = 200$:
\begin{itemize}
  \item $p_x$ matches the analytic prediction to $1.1\times 10^{-6}$
    relative (the float32 precision floor, identical to Stage 9f's
    measurement at the same kinematic point).
  \item $p_y = 0$ exactly (no transverse force).
  \item Sign reversal $A_0 \to -A_0$: $|p_x(+) + p_x(-)| < 10^{-8}$ ---
    exact antisymmetry to round-off.
  \item Neutral particle ($q = 0$): zero force, as expected.
\end{itemize}

\paragraph{Status of the EM Lorentz force after Stages 9e/9f/9g.}
\begin{center}
\begin{tabular}{@{}lccc@{}}
\toprule
\textbf{Piece of $\vv{F}_{\mathrm{em}}$} & \textbf{Implemented} & \textbf{Verified by} & \textbf{Evaluator} \\
\midrule
$+ q\,\vv{v}\!\times\!\vv{B}$ (magnetic Lorentz) & yes & Stage 9e & \texttt{B\_em\_z\_at\_total} \\
$- q\nabla\phi$ (electrostatic)                  & yes & Stage 9f & \texttt{phi\_em\_grad\_at\_total} \\
$- q\,\partial_t\vv{A}$ (inductive)              & yes & Stage 9g & \texttt{dt\_A\_em\_at\_total} \\
\bottomrule
\end{tabular}
\end{center}
All three pieces share the same three-path structure: \textsc{analytic}
background (closed form via \texttt{gr\_bg\_eval\_*\_em}), \textsc{sampled}
background (Yee curl or centered FD of the \texttt{*\_bg} arrays), and
perturbation (always sampled, using the leapfrog buffer rotation for the
centered-time-difference in the inductive case).  A charged particle in
any combination of these three sources --- static analytic background,
sampled-grid background, and a time-evolving PIC-deposited perturbation
--- now feels the complete EM Lorentz force.  The closed PIC loop
(\emph{moving charge $\to$ Esirkepov current $\to$ wave-equation
$\to$ $\vv{A}^{\mathrm{pert}}$, $\phi^{\mathrm{pert}}$ $\to$ Lorentz force
on a probe charge}) is now end-to-end functional and yields a
CFL-convergent physical radiation-reaction inspiral
(Section~\ref{sec:stage_em_pic_orbit}, with CFL convergence verified in
Section~\ref{sec:stage_cfl_convergence}).  Closed-loop stability
required reinterpreting the $\vv{A}$ buffers as half-integer time
slices (the v36 architecture fix); the wave-equation code itself did
not change.

\subsection{Stage 9h: EM Kepler analog and PIC orbit (closed-loop)}
\label{sec:stage_em_pic_orbit}

The EM analog of Stage~10 (Section~\ref{sec:stage_perturb_fdtd}): a
charged test particle orbits a softened analytic point-charge background,
with the perturbation EM fields evolving via the wave equation under
sources deposited by the particle itself.  This is the closed-loop PIC
verification on the EM side, exposing any feedback instabilities in the
combined deposit$\,$+wave-equation$\,$+force-interp chain.

\paragraph{EM Kepler baseline (Stage~26).}  Before the closed loop, a
basic verification: charged particle in an analytic softened-Coulomb
background $\phi^{\mathrm{bg}} = +k_e Q / \sqrt{r^2+\epsilon^2}$ with
opposite-sign $Q$ and $q_{\mathrm{test}}$ (attractive).  With
$|q_{\mathrm{test}} Q| k_e / m_{\mathrm{test}} = G_{\mathrm{eff}} M$, the
dynamics are mathematically identical to gravity's Stage~7 Kepler test.
The implementation reproduces the analytic orbital period to $0.03\%$
and the orbit stays circular to $4\times 10^{-5}$ over one period
(\texttt{stage26\_em\_kepler.c}).  This verifies the EM force chain
\textsc{point\_charge bg} $\to$ \texttt{gr\_bg\_eval\_phi\_em} $\to$
\texttt{phi\_em\_grad\_at\_total} $\to$ \texttt{em\_force\_at} on a
closed-orbit observable, without any perturbation FDTD.

\paragraph{Closed-loop PIC orbit (Stage~27).}  Same setup, but
\texttt{field\_evolution} and \texttt{particle\_source\_deposition} both
on.  At $Q = +1$, $q_{\mathrm{test}} = -10^{-3}$, $m_{\mathrm{test}} =
10^{-3}$ (same coupling as gravity's Phase E sweet spot), $r = 20$,
CIC$+$LEGACY$+$smoothing $=4$, over $4$ orbital periods, the radius
drifts:
\begin{center}
\begin{tabular}{@{}lcc@{}}
\toprule
\textbf{Variant} & \textbf{Drift after 4 orbits} & \textbf{Mechanism} \\
\midrule
Baseline (\texttt{field\_evolution} OFF)         & $+0.003\%$         & --- (analytic bg only) \\
PIC, all three EM Lorentz pieces, centered $\partial_t\vv{A}$ & $\mathbf{+18.6\%}$ & PIC heating dominates \\
PIC, inductive ($-q\,\partial_t\vv{A}$) disabled  & $\mathbf{-7.0\%}$  & radiation-reaction-style inspiral \\
PIC, backward-diff $\partial_t\vv{A}$            & \textbf{unbound by orbit 4} & worse heating (see below) \\
\bottomrule
\end{tabular}
\end{center}

\emph{The classical-atom-collapse inspiral is present in this sandbox
when the inductive piece is disabled.}  The $-7.0\%/4\,\mathrm{orbits}$
inward drift (no inductive) is the same order as gravity's Phase E
result at the same coupling ($-2.4\%$ per orbit), consistent with
radiation-reaction physics: an orbiting charge radiates EM waves into
the absorbing boundary, loses orbital energy, and spirals in.

But \emph{with the inductive piece enabled, PIC numerical heating
overwhelms radiation reaction by an order of magnitude}, driving an
outward $+18.6\%/4\,\mathrm{orbits}$ instead of an inspiral.  The
mechanism is specific to the inductive piece: the centered-time
derivative
\[
  \partial_t\vv{A}^n \;=\; (\vv{A}^{n+1} - \vv{A}^{n-1})/(2\,\Delta t)
\]
reads $\vv{A}^{n+1}$, which contains the particle's own current
deposit at this step.  Although the time-CENTER is at $t^n$ (matching
the Stage~18 read-prev fix's timing principle for the spatial
gradient), the discrete propagation of the self-field into $\vv{A}^{n+1}$
appears to inject a net positive-feedback torque on the orbit.

\paragraph{Investigation and resolution.}  The initial closed-loop
result (with all three EM Lorentz pieces enabled, $A$ stored at integer
time) showed $+18.6\%$ outward drift over 4 orbits at $q = m = 10^{-3}$
--- runaway PIC heating, not the expected radiation-reaction inspiral.
A symmetric gravity-side diagnostic (Stage 28, wiring the analogous
$-m\,\partial_t\vv{A}_g$ piece into the gravity Lorentz force) exhibited
the SAME heating ($+23\%/4\,\mathrm{orbits}$), confirming that the
phenomenon was structural to the discretization rather than EM-specific.
Three further diagnostic experiments --- sign-flip on the inductive
term, backward-difference discretization, and source-time correction
$J^{n-1/2} \to J^n$ via linear extrapolation --- all merely rotated
the direction of the closed-loop bias without taming its magnitude,
ruling out simple sign-convention or single-step source-timing bugs.

The resolution was a structural insight: in the integer-time storage
convention, $\partial_t\vv{A}^n = (\vv{A}^{n+1} - \vv{A}^{n-1})/(2\Delta t)$
spans a 2-step stencil and necessarily reads $\vv{A}^{n+1}$, which
contains the particle's own current-step deposit (asymmetric self-
deposit visibility relative to the spatial-derivative reads at
$\vv{A}^n$).  Furthermore $J^{n-1/2}$ from Esirkepov is at a half-step
offset from the operator's integer-time-center, an independent but
related inconsistency.

\paragraph{Half-step $\vv{A}$ architecture (v36 fix).}  Reinterpret the
$\vv{A}$ buffers (both $\vv{A}_g$ and $\vv{A}_{\mathrm{em}}$) as living
at half-integer times: \texttt{.prev} = $\vv{A}^{n-1/2}$, \texttt{.curr}
= $\vv{A}^{n+1/2}$.  The wave-equation leapfrog code does not change ---
only the time-interpretation shifts.  This collapses BOTH issues at
once:
\begin{itemize}
  \item $J^{n-1/2}$ from Esirkepov now naturally matches $\vv{A}^{n-1/2}$
    in the source/operator pairing of the wave-equation update for
    $\vv{A}^{n+1/2}$.
  \item The pusher's $\partial_t\vv{A}^n$ becomes
    $(\texttt{.curr} - \texttt{.prev})/\Delta t = (\vv{A}^{n+1/2} -
    \vv{A}^{n-1/2})/\Delta t$ --- a tight 1-step centered difference
    at $t^n$, with no 2-step asymmetric self-deposit reach.
  \item The pusher's $\vv{B}^n$ becomes the average of curl on the two
    surrounding half-steps: $\vv{B}^n \approx \tfrac{1}{2}(\nabla\!\times\!
    \vv{A}^{n+1/2} + \nabla\!\times\!\vv{A}^{n-1/2})$.
\end{itemize}
The scalar $\phi$ stays at integer time (sourced by $\rho$ at integer
time, no inconsistency).  This mirrors standard Yee$+$Boris EM PIC's
half-step interleaving of $\vv{E}$ and $\vv{B}$, transposed to the
potentials.

\paragraph{Result after the half-step fix.}  Rerunning Stage 27 (EM)
and Stage 28 (gravity, with the GEM inductive piece enabled):

\begin{center}
\begin{tabular}{@{}lcc@{}}
\toprule
\textbf{Variant} & \textbf{Pre-fix (integer $\vv{A}$)} & \textbf{Post-fix (half-step $\vv{A}$)} \\
\midrule
EM, no inductive            & $-7.0\%$    & $-6.9\%$  \quad (unchanged) \\
EM, full Lorentz            & $+18.6\%$   & $\mathbf{-9.2\%}$  \quad (inspiral!) \\
Gravity, no inductive       & $-0.59\%$   & $-0.54\%$ \quad (unchanged) \\
Gravity, full Tier-3        & $+23.0\%$   & $\mathbf{-9.1\%}$  \quad (inspiral!) \\
\bottomrule
\end{tabular}
\end{center}

The catastrophic heating is gone in both sectors.  Both sectors show
matching inspiral rates ($\sim -9\%$ per 4 orbits, $\approx -2.3\%$ per
orbit), consistent with a 2D radiation-reaction signal.  Spatial-piece
behaviour is unchanged (the static-background reads are bit-clean
because averaging two identical static buffers gives the same value).

\subsubsection{CFL convergence (Stages 29 and 30)}
\label{sec:stage_cfl_convergence}

The half-step fix yields an inspiral, but is the inspiral physical
radiation reaction or residual numerical truncation?  The test is a
CFL convergence sweep: a physical effect should be CFL-independent at
small $\Delta t$, while a numerical truncation residual should scale
as $\order{\Delta t^2}$ (ratio 4 between consecutive CFL halvings).

\paragraph{Gravity (Stage 29) at $m_{\mathrm{test}} = 10^{-3}$, 4 orbits:}

\begin{center}
\begin{tabular}{@{}cccc@{}}
\toprule
\textbf{CFL} & \textbf{No-inductive drift} & \textbf{Inductive-on drift} & \textbf{Contribution} \\
\midrule
$1/\sqrt{2}      \approx 0.707$ & $-0.535\%$ & $-9.05\%$  & $-8.52\%$ \\
$1/(2\sqrt{2})   \approx 0.354$ & $-0.391\%$ & $-18.29\%$ & $-17.90\%$ \\
$1/(4\sqrt{2})   \approx 0.177$ & $-0.276\%$ & $-19.07\%$ & $-18.80\%$ \\
$1/(8\sqrt{2})   \approx 0.088$ & $-0.175\%$ & $-19.16\%$ & $-18.99\%$ \\
\bottomrule
\end{tabular}
\end{center}

Consecutive-CFL contribution ratios: $0.476,\ 0.952,\ 0.990$.  The
ratio approaches $1.0$ (CFL-independent plateau) as $\Delta t \to 0$
--- the signature of a \emph{physical} effect.  The inductive
contribution converges to $\sim -19\%/4\mathrm{orbits} \approx -4.75\%$
per orbit.

\paragraph{EM (Stage 30) at $q_{\mathrm{test}} = -10^{-3}$, $m =
10^{-3}$, 4 orbits:}

\begin{center}
\begin{tabular}{@{}cccc@{}}
\toprule
\textbf{CFL} & \textbf{No-inductive drift} & \textbf{Inductive-on drift} & \textbf{Contribution} \\
\midrule
$1/\sqrt{2}      \approx 0.707$ & $-6.93\%$ & $-9.19\%$  & $-2.25\%$ \\
$1/(2\sqrt{2})   \approx 0.354$ & $-5.97\%$ & $-19.15\%$ & $-13.18\%$ \\
$1/(4\sqrt{2})   \approx 0.177$ & $-5.40\%$ & $-19.73\%$ & $-14.33\%$ \\
$1/(8\sqrt{2})   \approx 0.088$ & $-4.97\%$ & $-20.47\%$ & $-15.50\%$ \\
\bottomrule
\end{tabular}
\end{center}

EM contribution ratios: $0.171,\ 0.919,\ 0.925$.  Same qualitative
pattern --- ratio approaches $1.0$, plateau emerges.  The EM
no-inductive drift was larger and slower-to-converge than gravity's
in this (pre-v36) table, an asymmetry that was tracked down to two
distinct issues:
\begin{itemize}
  \item The EM scalar gradient lacked TSC/LB variants.  v36 ports
        these from \texttt{grav\_grad\_at} onto
        \texttt{phi\_em\_grad\_at\_total}; Stage 33 confirms the four
        (shape $\times$ interp) combinations give consistent drift
        post-fix.
  \item More fundamentally, the EM source-coupling sign was inverted
        relative to standard Maxwell (\texttt{c\_em} was negative,
        like gravity's $-4\pi G$, instead of $+4\pi k_e$).  v36 fixes
        this; the EM PDE now reads $\Box \phi = -\rho/\epsilon_0$ and
        $\Box \vv{A} = -\vv{J}/(\epsilon_0 c^2)$ (standard Maxwell).
        Stage 34 verifies the full chain end-to-end:
        Poisson sign, $\vv{J}=\rho\vv{v}$ deposition, $\vv{A}$
        wave-equation sign, and Lorentz force direction (including
        like-repel + opposite-attract).
\end{itemize}
\paragraph{Open item -- inductive-piece sign mismatch.}  The post-fix
inductive contribution at the same CFL sweep reads
$+3.5\%,\ +23.4\%,\ +25.4\%,\ +22.3\%$ (drift, not
contribution-only) --- an \emph{outspiral} that is also CFL-converged
(ratio approaching $1.0$).  Pre-v36 the same setup gave inspiral.
The sign flip on the field convention also flips the sign of the
inductive-piece self-coupling, reversing the discrete-PDE
radiation-reaction direction.  Energy conservation demands inspiral
for a bound orbit, so the outspiral indicates a sign error in the
inductive force-law term relative to the new field convention.

\textbf{Diagnostic chain (Stage 34--36).}  The v36 sign convention
is verified end-to-end on the static EM machinery: Poisson sign,
$\vv{J}=\rho\vv{v}$ deposition, $\vv{A}$ wave-equation sign, all
four Lorentz-force direction tests pass standard Maxwell
(Section~\ref{sec:stage_em_chain_signs}, Stage~34).  The pure
$q\vv{v}\times\vv{B}$ piece gives Amp\`ere attraction for parallel
like-currents (Stage~35) with the expected $(v/c)^2$ scaling.  When
the inductive piece is OFF, the EM Kepler orbit is essentially
stable ($\sim +0.3\%$ drift).  Stage~36 confirms that flipping the
sign on $-q\partial_t\vv{A}$ (via the diagnostic
\texttt{gr\_sim\_set\_em\_inductive\_sign(-1)}) turns outspiral
into inspiral:

\begin{center}
\begin{tabular}{@{}lc@{}}
\toprule
inductive\_sign & drift over 4 orbits \\
\midrule
$+1$ (default standard Lorentz)  & $+3.73\%$ \,(\emph{outspiral}) \\
$-1$ (flipped)                    & $-3.45\%$ \,(\emph{inspiral})  \\
\bottomrule
\end{tabular}
\end{center}

So the surgical fix is a one-sign change on the
$-q\partial_t\vv{A}$ term in \texttt{em\_force\_at}.  The deeper
question --- why the discrete PIC self-force interaction needs the
flipped inductive sign in the standard-Maxwell field convention
(but not in the gravity-style anti-Maxwell convention used pre-v36)
--- is documented in \texttt{memory/grlite-em-inductive-sign-open.md}
as the next investigation step.

\paragraph{Interpretation.}  The inductive piece's inspiral is a
genuine physical radiation reaction in the 2D-softened-Newton$+$GEM
toy system.  The magnitude ($\sim 5\%$ per orbit at $v/c \approx 0.22$)
is much larger than a 3D Larmor estimate ($\sim 0.01\%$ per orbit at
the same parameters) because:
\begin{itemize}
  \item The 2D wave equation has a logarithmic Green's function vs.\
    $1/r$ in 3D; radiation scaling with source motion differs
    fundamentally.
  \item Gravity is softened-Newtonian here, not full GR.
  \item The toy is self-consistent for the equations it actually solves;
    those equations are not Maxwell-3D, so quantitative agreement with
    3D Larmor is not expected.
\end{itemize}

The pedagogical conclusion holds: when the full force law is active,
the orbit slowly spirals in.  The classical-atom-collapse picture is
faithfully demonstrated in this sandbox, with a CFL-convergent inspiral
rate that emerges as a real prediction of the discrete PDE.

\subsubsection{Stage 31: Shapiro delay in the EM wave equation}
\label{sec:stage_shapiro_delay}

The doc's headline EM physics addition: light slows in a gravitational
potential, as derived from the null geodesic condition
(Section~\ref{sec:shapiro}).  The implementation replaces the uniform
wave speed $c$ in the EM-field leapfrog with a per-cell
$c_{\mathrm{local}}(x) = c (1 + 2 \Phig(x) / c^2)$ on the Laplacian
term; the source-coupling term keeps the bare $c$, since Shapiro
modifies free propagation rather than source coupling.

Only the three EM fields ($\phi$, $A_x$, $A_y$) are affected; the GEM
fields keep the uniform $c$ (gravity propagates at $c$ regardless of
background gravity, by construction of the toy).  Implementation:
three per-sublattice arrays $c_{\mathrm{local}}^2$ on CORNER, X-EDGE,
Y-EDGE are precomputed at enable-time from the installed $\Phig$
background, sampled at each EM-field's own Yee node positions.

\paragraph{Verification (Stage 31).}  An EM pulse (Gaussian initial
condition on $\phi$ with zero initial time derivative, $\sigma = 4$
cells, amplitude $1$) is launched at $(c_x - L/2,\ c_y)$ for $L = 120$
cells, with a softened point-mass background ($GM = 0.5$, $\epsilon = 4$)
at the grid center.  The arrival of the wavefront at the diametrically
opposite probe is timed via parabolic-interpolated peak amplitude.

\begin{center}
\begin{tabular}{@{}lcc@{}}
\toprule
                    & \textbf{No mass (baseline)} & \textbf{With mass (Shapiro)} \\
\midrule
Peak $|\phi|$       & 0.0652  & 0.0802 \\
Arrival time        & 117.84  & 124.10 \\
\bottomrule
\end{tabular}
\end{center}

Measured Shapiro delay: $\Delta t_{\mathrm{meas}} = +6.26$.  Analytic
prediction (Eq.~\ref{eq:shapiro_delay} integrated along the axial ray
through a softened point mass at impact parameter $b = 0$):

\begin{equation}
  \Delta t_{\mathrm{ana}} = \frac{4 G M}{c^3}\,
    \operatorname{asinh}\!\left(\frac{L}{2 \epsilon}\right) = +6.80.
\end{equation}

Agreement within $\sim 8\%$ --- well within the pedagogical tolerance for
a 2D-radial wavefront whose path samples a band of $\Phig$ values, not
strictly the axial ray.  The amplitude focusing ($0.065 \to 0.080$) is
the real signature of gravitational lensing: $c_{\mathrm{local}}$ is
smaller near the central well, bunching the wavefront on the slow side
of the obstacle.  See \texttt{core/tests/stage31\_shapiro\_delay.c}.

\paragraph{Stability.}  For attractive $\Phig < 0$ backgrounds,
$c_{\mathrm{local}} < c$ everywhere, so the CFL bound stays at
$c\,\Delta t / \Delta x \le 1/\sqrt{2}$ (2D, 5-point Laplacian).
A defensive floor clamps $1 + 2\Phig/c^2 \ge 10^{-3}$ to guard against
pathological strong-field probes outside the weak-field regime in
which the linearized $c_{\mathrm{local}}$ formula is valid; for the
Stage~31 configuration $\min(1 + 2\Phig/c^2) \approx 0.75$, far from
the floor.

\subsubsection{Stage 32: EM binary PIC -- two-body mutual EM dynamics}
\label{sec:stage_em_binary_pic}

Two charged particles mutually orbit via the FDTD perturbation field
alone (no analytic background).  Each deposits ($\rho_q$, $J_{qx}$,
$J_{qy}$) via Esirkepov; each feels the other's $\phi_{\mathrm{em}}$
and $\vv{A}_{\mathrm{em}}$ perturbation through the full EM Lorentz
force.  This is the EM analog of the Stage~12 gravity binary; together
with Stage~27 (single test charge in a softened-charge background) it
exercises the entire EM Lorentz machinery end-to-end with two active
sources.

\paragraph{Sign convention (v36 fix).}  The original code shared a
single source-coupling sign between gravity and EM ($\mathrm{sc}_{em}
= -4\pi k_e$, matching $\mathrm{sc}_{\mathrm{grav}} = -4\pi G$).
For gravity this is the standard Newton Poisson sign (positive masses
attract); but for EM it produced an ``anti-Maxwell'' channel in which
PIC same-sign charges attracted.  The analytic point-charge background
hand-coded with the standard $+k_e Q / r$ shape was inconsistent with
the PIC sign, an asymmetry that didn't matter for Stages 23--27 (which
use analytic backgrounds and tiny PIC perturbations) but did show up
once Stage 32 promoted both bodies to active PIC sources.  v36 flips
$\mathrm{sc}_{em}$ to $+4 \pi k_e$ (and the corresponding $A_{em}$
coupling), so the PIC channel now obeys standard Maxwell:
\emph{like charges repel, opposite charges attract}.  The analytic
$1/r$ background remains standard-EM-signed and is now sign-consistent
with the PIC channel; the two paths still differ in spatial shape
($1/r$ vs.\ box-bounded 2D log), so they should not be mixed in the
same simulation.  Stage 32 uses opposite-sign charges
($q_1 = +Q$, $q_2 = -Q$) for the attractive orbit.

\paragraph{Centripetal balance.}  Identical structure to the gravity
binary with $G_{\mathrm{eff}} m \to k_e Q^2$ and $v_{\mathrm{orb}}$
independent of $r$:

\begin{equation}
  v_{\mathrm{orb}} = Q\sqrt{k_e / m},
  \qquad T = \frac{2\pi r}{v_{\mathrm{orb}}}.
\end{equation}

\paragraph{Verification (Stage 32).}  Three checks mirror Stage 12:

\begin{itemize}
  \item \emph{Free-fall:} two opposite-sign charges released at rest
        attract, falling together from $r = 8$ to $r_{\min} \approx 0$
        within ~100 steps.  Mutual attraction has the right sign.
        Esirkepov violations $= 0$ throughout.
  \item \emph{Short-horizon orbit:} at $r = 8$, $v_{\mathrm{orb}} = 0.1$,
        the empirical $\omega$ matches $v_{\mathrm{orb}}/r$ to $\lesssim 5\%$
        in the first wave-crossing time.
  \item \emph{$v_{\mathrm{orb}}$ independent of $r$:} at $r = 6$ and
        $r = 12$, $v_{\mathrm{meas}} = r \omega$ values match each
        other and the analytic $Q\sqrt{k_e/m}$ to within a few percent
        --- the 2D-log signature on the EM channel.
\end{itemize}

\paragraph{Stage 32b: sign-convention verification.}  Three direct
sign tests on the PIC EM channel:
\begin{itemize}
  \item Two like-sign charges at rest \emph{repel} (Coulomb sign).
  \item Two opposite-sign charges at rest \emph{attract}.
  \item $\phi_{\mathrm{em}}$ from a single $+Q$ source is positive and
        decreases monotonically with distance (the expected
        box-bounded 2D Green's function shape).
\end{itemize}
The clean Amp\`ere PIC-dynamics test (parallel like-currents attract
via $\vv{v}\times\vv{B}$ at $\order{(v/c)^2}$) is deferred: at the
$v/c \gtrsim 0.1$ needed to see a measurable magnetic effect, the
moving-source PIC self-heating (Cherenkov-like artifacts at the
discrete grid scale) dominates over the small Amp\`ere correction.
The $\vv{v}\times\vv{B}$ \emph{force law} is independently verified
in Stage 23 (cyclotron in the analytic uniform-$B$ background).

Long-time orbit stability is left for a future test, with the same
caveats as Stage 12: sub-cell PIC heating and the inductive
radiation-reaction inspiral act on overlapping timescales here; teasing
them apart will need either a weak-coupling regime ($v/c \ll 1\%$) or
a CFL sweep analogous to Stages 29/30.  Stage 32 is the basic
end-to-end EM binary validity check.  See
\texttt{core/scenarios/pic\_binary\_em.c} and
\texttt{core/tests/stage32\_em\_binary\_pic.c}.

\subsection{Stage 10: Perturbation FDTD with moving particle source}
\label{sec:stage_perturb_fdtd}

Enable the full perturbation FDTD: the particle deposits $\rho_{\mathrm{matter}}$ and $\vv{J}_{\mathrm{matter}}$ onto the grid each step (new \texttt{gr\_sim\_set\_particle\_source\_deposition} flag, set on by the Stage~10 scenarios), the wave equations evolve the perturbation potentials, and the particle feels both background and perturbation forces via the same force evaluator used in Stages~7--9.

\subsubsection*{Weak-coupling regime}

A subtle issue with this stage: the Hockney--Eastwood adjoint condition
guarantees zero self-force for a \emph{stationary} particle in the
statically-converged field, but for a moving particle the residual
self-force is a finite mixture of physical radiation reaction and
numerical grid-heating.  At the Stage~7/8 setup ($G_{\mathrm{eff}} M = 1$,
$r = 20$, $m_{\mathrm{test}} = 1$), the grid-heating term is strong enough
to unbind a circular orbit on a timescale much shorter than the orbital
period.

The doc's Stage~10 specification implicitly assumes the weak-coupling
regime $G_{\mathrm{eff}} m_{\mathrm{test}} / (c^2 r) \ll 1$, where both
the radiation-reaction and the grid-heating contributions are
proportional to $m_{\mathrm{test}}$ and become negligible at low
$m_{\mathrm{test}}$.  In the implementation, the
\texttt{pic\_orbiting} scenario defaults to
$m_{\mathrm{test}} = 10^{-3}$ for the orbiter; the Stage~10 test uses
$10^{-6}$ for the orbit-regression phase.  This keeps the dynamics in
the test-particle limit while still exercising the deposition path.

\subsubsection*{Test (single-particle suite)}

A second particle is not needed to verify the deposition$+$FDTD chain;
four phases on a single particle suffice:

\begin{itemize}
  \item \emph{(A) Static-source field shape.}  Stationary particle in
        vacuum.  Time-averaged $\Phig^{\mathrm{pert}}(r)$ should follow
        the 2D Green's function $\Phig = 2 G_{\mathrm{eff}} m \ln r + C$
        (the slope of $\Phig$ vs $\ln r$ equals $2 G_{\mathrm{eff}} m$).
        Parallel to Stage~3, but exercising the auto-deposit path.
  \item \emph{(B) Boosted-Coulomb shape.}  Particle moving at constant
        $v$ in vacuum.  $\Ag^{\mathrm{pert}}_x$ and
        $\Phig^{\mathrm{pert}}$ at the same lab points should have the
        same sign and the same Green's-function support.  The 3D-style
        leading-order ratio $|\Ag|/|\Phig| \approx v/c^2$ does
        \emph{not} hold instantaneously in 2D because of the non-local
        retardation structure of the 2D wave kernel and the fact that
        the auto-deposit moves with the particle.  Stage~5 already
        validates vector-potential generation from a moving source via
        the curl/$B_z$ test; this phase is informational.
  \item \emph{(C) Self-force cancellation.}  Stationary particle in
        vacuum, $10^4$ steps.  The particle's position drift should be
        $\ll 1$ cell.  This is the Hockney--Eastwood theorem operating:
        with CIC deposition paired with CIC-of-FD force interpolation,
        a static particle exerts zero net force on itself in the
        statically-converged field.
  \item \emph{(E) Orbit regression.}  Particle in a circular orbit
        around an analytic softened-point-mass background (Stage~7
        setup) with the perturbation FDTD on and the test-particle
        mass in the weak-coupling regime.  Orbital period should match
        Stage~7's result to within $\sim 1\%$; the residual is the
        sum of the (physical) radiation-reaction shift and the (small,
        $\propto m_{\mathrm{test}}$) numerical heating.  In the
        implementation at $m_{\mathrm{test}} = 10^{-6}$ the rel.\ diff
        is $\approx 0.1\%$.
\end{itemize}

\subsubsection*{Outgoing-wave 1/\textsurd{}r falloff (deferred)}

The original v33 spec for this stage also called for verifying the
$1/\sqrt{r}$ amplitude falloff of outgoing gravitational waves from an
orbiting source.  At the test-particle mass needed for orbit stability
($m_{\mathrm{test}} \sim 10^{-6}$), the radiated amplitude is at the
$10^{-10}$ level relative to the background, well within the
discretization-noise floor of the simple CIC$+$FD setup.  A
quantitative slope-fit is deferred to a later Stage 10b once the web
visualization is online; the radiation pattern is qualitatively visible
in the perturbation field and is the cleanest visual demonstration of
gravitational radiation in the sandbox.

\subsubsection*{The grid-heating problem: why orbits unbind, not spiral in}
\label{sec:stage10_heating}

A physical orbiting mass should LOSE energy to outgoing gravitational
waves and slowly spiral inward (the binary-pulsar story).  In our
simulation we observe the opposite: the orbit \emph{gains} energy and
unbinds.  This is the well-known PIC \textbf{grid-heating} artifact
applied to a gravitational-source setup, with the same sign and
mechanism as in plasma-PIC codes.

The Hockney--Eastwood theorem guarantees zero self-force for a
\emph{stationary} particle in the statically-converged field (verified
quantitatively by Test~C).  For a \emph{moving} particle the
adjoint-condition cancellation fails by a residual proportional to the
deposited density's time and space derivatives.  In a
charge-conserving deposition scheme (e.g.\ Esirkepov 2001 on a Yee
staggered grid) this residual is bit-zero by construction.  In our
simpler CIC$+$FD setup on a cell-centered grid, the residual is small
but nonzero, and its sign systematically pumps energy \emph{into} the
particle rather than out of it.

The numerical-self-force-to-physical-centripetal ratio scales as
$m_{\mathrm{test}}\,r/M_{\mathrm{bg}}$ in our units, while the
condition for the linearized GR equations themselves to hold scales as
$r_s^{\mathrm{test}}/r = 2 G m_{\mathrm{test}}/(c^2 r)$ — the numerical
condition is the tighter one, with an extra factor of $r/r_s^{\mathrm{bg}}$.
For $G = c = M_{\mathrm{bg}} = 1$, $r = 20$: the linearization
threshold is $m \ll 10$ but the numerical-stability threshold is
$m \ll 0.05$.

\paragraph{Cell-centered Esirkepov is not available.}
Esirkepov's exact-discrete-continuity recurrence is built around the Yee
staggered grid ($\rho$ at corners, $J_i$ at face centers) and its
particular divergence stencil.  Our cell-centered storage (Section~\ref{sec:numerical})
uses a different divergence operator (typically 2-point centered, spanning
$2\Delta x$), for which Esirkepov's local recurrence does not directly apply.
Adapting it would be a research-level exercise, and migrating the whole
grid to Yee staggering — which we explicitly considered and rejected in
the Stage 4/5 design — is not worth the refactor for a pedagogical sandbox.

\paragraph{Mitigations available on the cell-centered grid.}
\begin{itemize}
  \item \emph{Stay weak-coupling.}  The path Stage 10 takes.  Doesn't
    model radiation reaction; correctly models the
    static and quasi-static dynamics the sandbox is built for.
  \item \emph{Higher-order kernel.}  Promoting matter deposition from CIC
    ($W_2$, 2-cell support) to TSC / quadratic B-spline ($W_3$,
    3-cell support) reduces the high-spatial-frequency content of the
    deposit and correspondingly reduces heating by ${\sim}10\times$.
    The $W_3$ kernel infrastructure is already planned for Stage 13's
    spin-dipole sources and can be promoted to matter sources when
    convenient.
  \item \emph{Symmetric binomial smoothing.}  Apply a $(1,2,1)/4$
    binomial filter $B$ to the source arrays after deposition AND to
    the field slab before particle interpolation, using the same $B$
    on both sides.  The effective deposition and interpolation kernels
    become $W_2 \star B$ and $B \star W_2$ respectively — equal, so the
    Hockney--Eastwood adjoint condition is \emph{preserved}.  $B$'s
    Fourier transform is $\cos^2(k\Delta x/2)$, which exactly nulls the
    Nyquist mode and strongly suppresses content near it — exactly where
    the heating residual is largest.  Trade-off: the effective particle
    size grows to ${\sim}4\Delta x$ instead of $2\Delta x$, a slight
    pedagogical-visualization cost.  Order-of-magnitude reduction in
    heating, not orders.
  \item \emph{Accept and document.}  Frame Stage 10 as ``the
    perturbation FDTD chain is verified to work in the linearization
    regime; quantitative radiation-reaction physics is outside the
    sandbox's current scope.''  The visual radiation pattern in the
    perturbation field remains a clean qualitative demonstration of
    gravitational waves regardless.
\end{itemize}

Asymmetric smoothing (filtering only the source, OR only the field) is
the wrong path: it shifts the effective deposition kernel without
shifting the interpolation kernel, breaking the adjoint condition and
introducing a static self-force that didn't exist before — trading one
artifact for another, possibly worse, one.  The symmetric form above is
the correct pedagogical recipe.

For the current sandbox, ``stay weak-coupling'' is the working
strategy; the upgrade to $W_3$ matter deposition and/or symmetric
binomial smoothing is the natural extension when Stage 13 brings the
$W_3$ kernel online.

\subsection{Stage 11: Two-particle gravitational interaction}

Two particles mutually attracting via the FDTD perturbation field, no analytic background.

\textbf{Test:}
\begin{itemize}
  \item Two equal masses at rest should fall toward each other with initial acceleration $G_{\mathrm{eff}}M/r^2$.
  \item A circular binary orbit should have the correct Keplerian period.
  \item Results should agree with a direct $N$-body calculation using the same force law, verifying the FDTD intermediary introduces no force errors.
  \item The total canonical momentum $\vv{P}_{\mathrm{total}} = \sum_i\vv{p}_i + c^{-2}\int\vv{S}_{\mathrm{GEM}}\,\dd^2x$ should be conserved.
\end{itemize}

\subsection{Stage 12: Charged particles and electromagnetism}

Add electric charge $q$ to particles.

\textbf{Test:}
\begin{itemize}
  \item A charge at rest should produce the 2D Coulomb potential $\phi \propto \ln r$.
  \item Two opposite charges should attract at the correct Coulomb force.
  \item A moving charge in a uniform background magnetic field $B_0$ (sampled into $\vv{A}^{\mathrm{bg}}$ per Stage~6) should undergo circular cyclotron motion at frequency $\omega_c = qB_0/m\gamma$.
  \item The EM vector potential $\vv{A}$ should correctly encode the magnetic field via the Yee curl.
\end{itemize}

\subsection{Stage 13: EM effective GEM sources}

Enable the $\rho_{\mathrm{EM}}$ and $\vv{J}_{\mathrm{EM}}$ source terms (Eqs.~\ref{eq:rho_em_pot}--\ref{eq:J_em_pot}).

\textbf{Test:} an EM pulse passing a test mass should exert a gravitational impulse via the Poynting flux term $\vv{J}_{\mathrm{EM}} = -(\nabla\phi^{\mathrm{pert}}+\partial_t\vv{A}^{\mathrm{pert}})\times(\nabla\times\vv{A}^{\mathrm{pert}})/(\mu_0 c^2)$.  With boosted coupling constants the deflection should be measurable and match the analytic estimate $\Delta p \sim G_{\mathrm{eff}}\int u_{\mathrm{EM}}\,\dd t / (c^2 b)$ where $b$ is the impact parameter.

\subsection{Stage 14: Intrinsic spin, precession, spin forces}

Add spin $\vv{S}$ to particles.  Implement the spin update (Section~\ref{sec:alg_spin}), the Tulczyjew-Dixon constraint correction, and the spin gradient forces.

\textbf{Test:}
\begin{itemize}
  \item A gyroscope in a Kerr-like gravitomagnetic background (sampled into $\Ag^{\mathrm{bg}}$ per Stage~6) should precess at rate $\Omega_{\mathrm{LT}} = 2G_{\mathrm{eff}}J/(c_{\mathrm{eff}}^2 r^3)$ for a circular orbit at radius $r$.
  \item A spinning charged particle in a background magnetic field $B_0$ should precess at the Larmor rate $\Omega_L = g_s qB_0/2m$.
  \item The spin magnitude $|\vv{S}|$ should be conserved to machine precision throughout.
  \item Thomas precession should contribute the correct rate $\Omega_{\mathrm{Thomas}} = -(\gamma^2/(\gamma+1))(\vv{v}\times\vv{a})/c^2$ for a circular orbit.
\end{itemize}

\subsection{Stage 15: Complete background spacetime library}

Extend the Stage~6 background-generator library with the remaining canonical spacetimes from §\ref{sec:bg}: Schwarzschild-like, Kerr-like, Reissner--Nordstr\"om-like, Kerr--Newman-like, uniform field, binary.  Each is just an additional generator that fills the appropriate combination of $\Phig^{\mathrm{bg}}$, $\Ag^{\mathrm{bg}}$, $\phi^{\mathrm{bg}}$, $\vv{A}^{\mathrm{bg}}$ arrays; no changes to the simulation core are required.

\textbf{Test} each background against its known analytic properties:
\begin{itemize}
  \item Schwarzschild: photon orbit (unstable) at $r = 3G_{\mathrm{eff}}M/c_{\mathrm{eff}}^2$.
  \item Kerr-like: frame dragging rate at given $r$ and spin $J$.
  \item Reissner-Nordström: Coulomb potential from background charge consistent with $\phi^{\mathrm{bg}}$.
  \item Binary: superposition of two Schwarzschild potentials at large separation.
\end{itemize}
Verify the background + perturbation split produces no spurious transients at startup (convergence iteration correctly initialized).

\subsection{Stage 16: Visualization pipeline}

Implement the WebGL rendering passes (Section~\ref{sec:viz}): field textures, quiver geometry, particle sprites.

\textbf{Test:}
\begin{itemize}
  \item Potential arrays display correctly as false-color images with correct dynamic range.
  \item On-demand field computation ($\vv{g}$, $B_{g,z}$, $\vv{E}$, $B_z$ from potentials) matches analytic values for the Stage~7--15 test configurations.
  \item Particle tracks display at the correct positions with fading.
  \item Coordinate time and proper time ticks are correctly spaced along the track.
  \item Spin indicator rotates at the correct precession rate.
  \item Force arrows point in the correct directions and scale correctly with force magnitude.
\end{itemize}

\subsection{Stage 17: Full integration benchmarks}

Run the complete set of pedagogical benchmark scenarios:

\begin{center}
\begin{tabular}{@{}L{4.5cm}L{7cm}@{}}
\toprule
\textbf{Scenario} & \textbf{Expected result} \\
\midrule
Twin paradox & Proper time clocks on different orbits show measurable difference after one orbital period \\
Lense-Thirring ring & Gyroscope ring around spinning mass precesses uniformly at $\Omega_{\mathrm{LT}}$ \\
Gravitational wave dipole & Outgoing $\Phig^{\mathrm{pert}}$ wave from oscillating mass, $1/\sqrt{r}$ amplitude fall-off \\
Shapiro delay & EM pulse delay past central mass matches $\Delta t = -2\int\Phig\,\dd\ell/c^3$ \\
EM wave scattering & EM pulse deflected by charged mass, trajectory bending matches analytic estimate \\
Binary inspiral & Two-body orbit slowly decays if radiation damping term is enabled \\
Spin-orbit coupling & Spinning particle orbit precesses differently from non-spinning particle \\
\bottomrule
\end{tabular}
\end{center}

All scenarios should produce qualitatively correct behavior and agree with the corresponding analytic predictions to within the expected numerical accuracy of the linearized GEM approximation.

%====================================================================
\section{References}
\label{sec:refs}
%====================================================================

This bibliography collects the works used directly in this design
document along with foundational references for the numerical techniques
the sandbox implements.  References cited inline above are marked with
keys here; the remaining entries are recommended further reading for
extending the implementation into stages not yet covered.

\subsubsection*{General relativity, weak-field, post-Newtonian}

\begin{thebibliography}{99}\setlength{\itemsep}{2pt}\small

\bibitem{AliHaimoud2019}
Y.~Ali-Ha\"\i moud, \emph{Introduction to gravitational lensing and the
post-Newtonian expansion}, NYU General Relativity Fall 2019 lecture
notes, Lecture~25.  Provides the linearized 1PN line element (eq.~41)
and the Einstein--Infeld--Hoffmann equation (eq.~37) in the form used
to derive the doc's Tier-2 force and Eq.~(\ref{eq:proper_time_full}).
A copy is in \texttt{docs/lecture25.pdf}.

\bibitem{Will1993}
C.~M.~Will, \emph{Theory and Experiment in Gravitational Physics},
2nd ed., Cambridge University Press, 2018 (1st ed.\ 1993).  The
canonical reference for the post-Newtonian formalism and the
parameterized-PN framework.  EIH 1PN equation in harmonic gauge:
\S6.2 / eq.~6.45.

\bibitem{MTW1973}
C.~W.~Misner, K.~S.~Thorne, J.~A.~Wheeler, \emph{Gravitation},
W.~H.~Freeman, 1973.  Kerr metric and frame dragging: \S33.  Geodesic
equation, linearized theory: chapters 18--19.

\bibitem{Wald1984}
R.~M.~Wald, \emph{General Relativity}, University of Chicago Press,
1984.  Slow-rotation limit and gravitomagnetic effects: \S6.6.
Linearized gravity: chapter 4.

\bibitem{Hartle2003}
J.~B.~Hartle, \emph{Gravity: An Introduction to Einstein's General
Relativity}, Addison-Wesley, 2003.  Lense--Thirring and the
gravitomagnetic limit: \S14.2.

\bibitem{Carroll2004}
S.~M.~Carroll, \emph{Spacetime and Geometry: An Introduction to
General Relativity}, Addison-Wesley, 2004.  Slow rotation and Kerr:
chapter~7.

\bibitem{Weinberg1972}
S.~Weinberg, \emph{Gravitation and Cosmology}, Wiley, 1972.  PN
expansion of the field equations follow Weinberg's chapter 9
treatment, the same one followed in \cite{AliHaimoud2019}.

\subsubsection*{Gravitoelectromagnetism and the clock effect}

\bibitem{Mashhoon2003}
B.~Mashhoon, \emph{Gravitoelectromagnetism: a brief review},
arXiv:gr-qc/0311030, 2003.  GEM-form line element and Lorentz
force, in a factor-of-2-in-force normalization that differs from this
document's factor-of-4-in-force convention by an overall $\Ag$
rescaling.

\bibitem{CohenMashhoon1993}
J.~M.~Cohen, B.~Mashhoon, \emph{Standard clocks, interferometry, and
gravitomagnetism}, Phys.~Lett.~A \textbf{181}, 353--358 (1993).  The
original derivation of the gravitomagnetic clock effect for prograde
and retrograde circular orbits in Kerr.  Direction of the effect
(prograde clocks accumulate proper time more slowly) is set by this
paper.

\bibitem{Iorio2001}
L.~Iorio, \emph{Satellite gravitational orbital perturbations and the
gravitomagnetic clock effect}, Int.~J.~Mod.~Phys.~D \textbf{10}, 465
(2001).  Re-derivation of the clock effect with explicit numerical
estimates for Earth satellites.

\bibitem{Tartaglia2000}
A.~Tartaglia, \emph{Detection of the gravitomagnetic clock effect},
Class.~Quantum Grav.~\textbf{17}, 783 (2000).  Discussion of
experimental detectability.

\bibitem{CiufoliniWheeler1995}
I.~Ciufolini, J.~A.~Wheeler, \emph{Gravitation and Inertia},
Princeton University Press, 1995.  Comprehensive treatment of frame
dragging including \S6.10 on the clock effect.

\bibitem{DamourDeruelle1985}
T.~Damour, N.~Deruelle, \emph{General relativistic celestial
mechanics of binary systems. I.~The post-Newtonian motion},
Ann.~Inst.~Henri Poincar\'e A \textbf{43}, 107 (1985).  The
Damour--Deruelle solution of the EIH equations of motion for binary
systems; used implicitly when the perihelion-precession test of
Stage~8 is compared against $\Delta\phi = 6\pi GM/(c^2 a(1-e^2))$.

\subsubsection*{FDTD, Yee grid, and absorbing boundaries}

\bibitem{Yee1966}
K.~S.~Yee, \emph{Numerical solution of initial boundary value problems
involving Maxwell's equations in isotropic media}, IEEE
Trans.~Antennas Propag.~\textbf{14}, 302 (1966).  The original Yee
staggered grid used for the perturbation FDTD; the layout in
\S\ref{sec:numerical} is its standard 2D form.

\bibitem{TafloveHagness2005}
A.~Taflove, S.~C.~Hagness, \emph{Computational Electrodynamics: The
Finite-Difference Time-Domain Method}, 3rd ed., Artech House, 2005.
Comprehensive FDTD reference; CFL stability, dispersion analysis, and
absorbing-boundary techniques in chapters~4--7.

\bibitem{Berenger1994}
J.-P.~B\'erenger, \emph{A perfectly matched layer for the absorption
of electromagnetic waves}, J.~Comput.~Phys.~\textbf{114}, 185 (1994).
Originator of the PML; the simpler quadratic-profile damping layer
used in \S\ref{sec:abc} is a degenerate matched-impedance case.

\subsubsection*{PIC and particle deposition}

\bibitem{HockneyEastwood1988}
R.~W.~Hockney, J.~W.~Eastwood, \emph{Computer Simulation Using
Particles}, Adam Hilger / IOP, 1988 (corrected reprint of the 1981
edition).  The standard PIC reference.  CIC ($W_2$) and TSC ($W_3$)
deposition kernels and the adjoint condition (\S5.4) cited
throughout \S\ref{sec:deposition}.  The Hockney--Eastwood self-force
cancellation theorem underlies the choice of identical kernels for
deposition and force interpolation.

\bibitem{BirdsallLangdon1991}
C.~K.~Birdsall, A.~B.~Langdon, \emph{Plasma Physics via Computer
Simulation}, Adam Hilger, 1991.  Standard PIC textbook complementary
to \cite{HockneyEastwood1988}; the leapfrog with Boris rotation is
treated in chapter~4.

\bibitem{Boris1970}
J.~P.~Boris, \emph{Relativistic plasma simulation--optimization of a
hybrid code}, Proc.~Fourth Conf.~on Numerical Simulation of Plasmas,
NRL, Washington DC, pp.~3--67 (1970).  Original presentation of the
Boris pusher; the relativistic magnetic-rotation step in
\S\ref{sec:numerical} follows this scheme exactly.

\bibitem{Esirkepov2001}
T.~Zh.~Esirkepov, \emph{Exact charge conservation scheme for
particle-in-cell simulation with an arbitrary form-factor},
Comp.~Phys.~Commun.~\textbf{135}, 144 (2001).  Local-recurrence
deposition that exactly satisfies discrete continuity; reserved for
adoption in later stages when the GEM perturbation FDTD is driven by
moving particle sources (Stage~10+).

\bibitem{VayBoris2008}
J.-L.~Vay, \emph{Simulation of beams or plasmas crossing at
relativistic velocity}, Phys.~Plasmas \textbf{15}, 056701 (2008).
Variant of the Boris pusher that fixes the residual $E\times B$
drift error at ultra-relativistic velocities; not currently used,
but relevant if photon-like ($v/c \to 1$) trajectories are ever
needed.

\end{thebibliography}

%====================================================================
\end{document}
%====================================================================
